% !TeX program = xelatex
% Compile with XeLaTeX or LuaLaTeX because this report contains Unicode
% symbols and box-drawing characters in equations and code examples.
% Suggested command: xelatex -interaction=nonstopmode <filename>.tex
% Run the command twice so that the table of contents is complete.

% Options for packages loaded elsewhere
\PassOptionsToPackage{unicode}{hyperref}
\PassOptionsToPackage{hyphens}{url}
\PassOptionsToPackage{dvipsnames,svgnames,x11names}{xcolor}
\documentclass[
  english,
  11pt,
  a4paper,
]{article}
\usepackage{xcolor}
\usepackage[margin=25mm]{geometry}
\usepackage{pdflscape}
\usepackage{amsmath,amssymb}
\setcounter{secnumdepth}{-\maxdimen} % remove section numbering
\usepackage{iftex}
\ifPDFTeX
  \usepackage[T1]{fontenc}
  \usepackage[utf8]{inputenc}
  \usepackage{textcomp} % provide euro and other symbols
\else % if luatex or xetex
  \usepackage{unicode-math} % this also loads fontspec
  \defaultfontfeatures{Scale=MatchLowercase}
  \defaultfontfeatures[\rmfamily]{Ligatures=TeX,Scale=1}
\fi
\usepackage{lmodern}
\ifPDFTeX\else
  % Menlo includes the box-drawing glyphs used by the architecture and
  % directory-tree diagrams. Fall back to the engine default if unavailable.
  \IfFontExistsTF{Menlo}{\setmonofont{Menlo}[Scale=MatchLowercase]}{}
\fi
\ifPDFTeX\else
  % xetex/luatex font selection
\fi
% Use upquote if available, for straight quotes in verbatim environments
\IfFileExists{upquote.sty}{\usepackage{upquote}}{}
\usepackage{fvextra}
\RecustomVerbatimEnvironment{verbatim}{Verbatim}{%
  breaklines=true,
  breakanywhere=true,
  fontsize=\small
}
\IfFileExists{microtype.sty}{% use microtype if available
  \usepackage[]{microtype}
  \UseMicrotypeSet[protrusion]{basicmath} % disable protrusion for tt fonts
}{}
\makeatletter
\@ifundefined{KOMAClassName}{% if non-KOMA class
  \IfFileExists{parskip.sty}{%
    \usepackage{parskip}
  }{% else
    \setlength{\parindent}{0pt}
    \setlength{\parskip}{6pt plus 2pt minus 1pt}}
}{% if KOMA class
  \KOMAoptions{parskip=half}}
\makeatother
\usepackage{longtable,booktabs,array}
\usepackage{caption}
\captionsetup[table]{skip=6pt}
\newcounter{none} % for unnumbered tables
\usepackage{calc} % for calculating minipage widths
% Correct order of tables after \paragraph or \subparagraph
\usepackage{etoolbox}
\AtBeginEnvironment{longtable}{\small}
\renewcommand{\arraystretch}{1.12}
\setlength{\LTpre}{0.6\baselineskip}
\setlength{\LTpost}{0.8\baselineskip}
\makeatletter
\patchcmd\longtable{\par}{\if@noskipsec\mbox{}\fi\par}{}{}
\makeatother
% Allow footnotes in longtable head/foot
\IfFileExists{footnotehyper.sty}{\usepackage{footnotehyper}}{\usepackage{footnote}}
\makesavenoteenv{longtable}
\ifLuaTeX
\usepackage[bidi=basic,shorthands=off]{babel}
\else
\usepackage[bidi=default,shorthands=off]{babel}
\fi
\ifLuaTeX
  \usepackage{selnolig} % disable illegal ligatures
\fi
\setlength{\emergencystretch}{3em} % prevent overfull lines
\providecommand{\tightlist}{%
  \setlength{\itemsep}{0pt}\setlength{\parskip}{0pt}}
\usepackage{bookmark}
\IfFileExists{xurl.sty}{\usepackage{xurl}}{} % add URL line breaks if available
\urlstyle{same}
% fallback for those not using the hyperref driver hyperxmp:
\makeatletter
\@ifundefined{xmpquote}{\newcommand{\xmpquote}[1]{#1}}{}
\makeatother
\hypersetup{
  pdftitle={Explainable Pedestrian--Vehicle Priority Negotiation at Smart Intersections},
  pdflang={en},
  colorlinks=true,
  linkcolor={blue},
  filecolor={Maroon},
  citecolor={Blue},
  urlcolor={blue},
  pdfcreator={LaTeX via pandoc}}

\title{Explainable Pedestrian--Vehicle Priority Negotiation at Smart
Intersections}
\author{}
\date{}

\begin{document}
\maketitle

{
\setcounter{tocdepth}{2}
\tableofcontents
}
\section{1. Executive Summary}\label{executive-summary}

The strongest defensible version of this research is an
\textbf{integrated algorithm--explanation--human-evaluation study}, not
a claim that pedestrian--vehicle game-theoretic control is itself new.
Wang et al.~already formulated pedestrian and vehicle signal phases as
players in a Nash bargaining game and evaluated the controller in
SUMO.\hyperref[source-1]{[1]} Other work already optimizes
pedestrian/vehicle signal
timing,\hyperref[source-4]{[4]}\,\hyperref[source-5]{[5]}\,\hyperref[source-6]{[6]}\,\hyperref[source-7]{[7]}
studies fairness in traffic
control,\hyperref[source-9]{[9]}\,\hyperref[source-10]{[10]}\,\hyperref[source-11]{[11]}
and explains traffic or automated-driving
decisions.\hyperref[source-2]{[2]}\,\hyperref[source-3]{[3]}\,\hyperref[source-18]{[18]}\,\hyperref[source-19]{[19]}
What remains underexplored is whether a transparent, explicitly
fairness-bounded allocation rule can expose a faithful decision trace
and whether those explanations improve \textbf{objective comprehension,
perceived procedural fairness, and calibrated---not merely
increased---trust}.

The recommended system is a university-scale, reproducible proof of
concept at one four-arm signalized intersection. A top-level controller
treats pedestrian service and vehicle service as two \textbf{virtual
advocacy agents}. This is a normative resource-allocation abstraction;
it does not assert that actual pedestrians or drivers calculate Nash
equilibria. At every admissible decision point, it generates safe
candidate phase plans, keeps only plans within a calibrated
efficiency-loss bound, and chooses among them with weighted Nash
bargaining. Waiting-time debt raises each mode's bargaining weight, hard
maximum-wait constraints prevent starvation, and emergency preemption is
handled lexicographically before ordinary optimization. A subordinate
max-pressure rule allocates vehicle green between north--south and
east--west approaches.

The selected action emits a structured \texttt{DecisionTrace}: input
values and freshness, normalized utilities, binding constraints,
candidate scores, the best rejected alternative, and the nearest
feasible input change that would flip the decision. A deterministic
template engine turns only that trace into factual and counterfactual
explanations. This architecture makes explanation fidelity testable and
avoids an unnecessary language model or post-hoc surrogate.

Evaluation has two linked but separable parts:

\begin{enumerate}
\def\labelenumi{\arabic{enumi}.}
\tightlist
\item
  A paired-seed SUMO experiment compares fixed-time, efficiency-only
  adaptive, and fairness-aware bargaining controllers. Primary outcomes
  are total person delay, pedestrian and vehicle waiting distributions,
  maximum wait, throughput, queue length, starvation violations,
  proportional-fairness measures, safety-constraint violations, and
  compute latency.
\item
  A preregistered within-participant HCI experiment compares
  decision-only, factual explanation, and factual-plus-counterfactual
  explanation interfaces using identical prerecorded scenario traces. It
  measures rationale comprehension, contrastive comprehension, response
  time, perceived procedural/informational fairness, usability, and
  trust calibration across sound and deliberately anomalous decisions.
\end{enumerate}

This design separates three concepts that must not be collapsed:
\textbf{objective allocative fairness} in the controller,
\textbf{procedural transparency} in the explanation, and
\textbf{perceived fairness} in the participant. It also separates
controller effects from interface effects: adding an explanation cannot
change the simulated traffic outcome for the same decision trace.

The recommended contribution is:

\begin{quote}
A starvation-safe, bounded-efficiency weighted-bargaining controller for
pedestrian--vehicle service allocation; a faithful trace-to-explanation
protocol with solver-validated counterfactuals; and a cross-layer
evaluation showing when those explanations help people understand,
judge, and appropriately rely on smart-intersection decisions.
\end{quote}

This is publishable only if the efficiency bound and utility weights are
calibrated and sensitivity-tested, the closest prior bargaining
controller is implemented or faithfully approximated, explanation
fidelity is verified, and the human study measures calibrated reliance
rather than celebrating higher trust as inherently desirable.

\section{2. Top 20 Papers}\label{top-20-papers}

This is a targeted integrative review, not a PRISMA-complete systematic
review. Candidate searches crossed pedestrian-aware/adaptive signal
control, multimodal priority, traffic fairness/equity, bargaining/game
models, explainable traffic control, automated-vehicle explanations,
algorithmic fairness perception, and trust in automation. Peer-reviewed
journal/conference papers were preferred; recent papers were included
when they materially changed the closest-work analysis, while
foundational HCI papers were retained for construct and experiment
design. The ranking favors direct usefulness to the proposed method over
raw citation count because current citation counts are database- and
date-dependent and systematically disadvantage new work.

All titles, author lists, years, venues, and DOI links below were
checked against DOI metadata and, where available, the primary publisher
or author paper. DOI metadata was rechecked on 11 September 2026.
``Year'' follows the issue year; Niels et al.~appeared online in 2023
and in the 2024 issue. The search should be updated immediately before
manuscript submission, and the bibliography export should be checked by
a reference manager rather than copied from this prose.

{\def\LTcaptype{none} % do not increment counter
\begin{longtable}[]{@{}
  >{\raggedleft\arraybackslash}p{(\linewidth - 4\tabcolsep) * \real{0.4000}}
  >{\raggedright\arraybackslash}p{(\linewidth - 4\tabcolsep) * \real{0.3000}}
  >{\raggedright\arraybackslash}p{(\linewidth - 4\tabcolsep) * \real{0.3000}}@{}}
\toprule\noalign{}
\begin{minipage}[b]{\linewidth}\raggedleft
Rank
\end{minipage} & \begin{minipage}[b]{\linewidth}\raggedright
Verified paper
\end{minipage} & \begin{minipage}[b]{\linewidth}\raggedright
Why it is in the top 20
\end{minipage} \\
\midrule\noalign{}
\endhead
\bottomrule\noalign{}
\endlastfoot
1 & Anyou Wang, Ke Zhang, Meng Li, Junqi Shao, and Shen Li (2023),
``Game Theory-Based Signal Control Considering Both Pedestrians and
Vehicles in Connected Environment,'' \emph{Sensors} 23(23), 9438.
\href{https://doi.org/10.3390/s23239438}{DOI} & Closest competitor:
phase players, Nash bargaining, pedestrian and vehicle queues, SUMO. \\
2 & Stefano Giovanni Rizzo, Giovanna Vantini, and Sanjay Chawla (2019),
``Reinforcement Learning with Explainability for Traffic Signal
Control,'' \emph{IEEE ITSC}, 3567--3572.
\href{https://doi.org/10.1109/ITSC.2019.8917519}{DOI} & Direct
traffic-signal XAI using SHAP; exposes the gap between technical
attribution and human-centered explanation evaluation. \\
3 & Wei-Li Liu, Jinghui Zhong, Peng Liang, Jianhua Guo, Huimin Zhao, and
Jun Zhang (2024), ``Towards Explainable Traffic Signal Control for Urban
Networks through Genetic Programming,'' \emph{Swarm and Evolutionary
Computation} 88, 101588.
\href{https://doi.org/10.1016/j.swevo.2024.101588}{DOI} & Inherently
readable symbolic signal rules; strongest recent interpretability
comparator. \\
4 & Qing He, K. Larry Head, and Jun Ding (2014), ``Multi-modal Traffic
Signal Control with Priority, Signal Actuation and Coordination,''
\emph{Transportation Research Part C} 46, 65--82.
\href{https://doi.org/10.1016/j.trc.2014.05.001}{DOI} & Foundational
multimodal, request-based control including pedestrians, cars, and
transit. \\
5 & Wanjing Ma, Dabin Liao, Yue Liu, and Hong Kam Lo (2015),
``Optimization of Pedestrian Phase Patterns and Signal Timings for
Isolated Intersection,'' \emph{Transportation Research Part C} 58,
502--514. \href{https://doi.org/10.1016/j.trc.2014.08.023}{DOI} &
Jointly selects pedestrian phasing and timing while representing safety
exposure, delay, detour, and noncompliance. \\
6 & Tanja Niels, Klaus Bogenberger, Markos Papageorgiou, and Ioannis
Papamichail (2024), ``Optimization-Based Intersection Control for
Connected Automated Vehicles and Pedestrians,'' \emph{Transportation
Research Record} 2678(2), 135--152.
\href{https://doi.org/10.1177/03611981231172956}{DOI} & Rolling-horizon
safe-slot optimization and pedestrian service constraints provide a
strong optimization baseline. \\
7 & Mobin Yazdani, Majid Sarvi, Saeed Asadi Bagloee, Neema Nassir, Jeff
Price, and Hossein Parineh (2023), ``Intelligent Vehicle Pedestrian
Light (IVPL): A Deep Reinforcement Learning Approach for Traffic Signal
Control,'' \emph{Transportation Research Part C} 149, 103991.
\href{https://doi.org/10.1016/j.trc.2022.103991}{DOI} & Calibrated
multimodal simulation and a reward based on total road-user delay; close
on mixed-mode adaptive control. \\
8 & Kangjie Xu, Junqin Huang, Linghe Kong, Jiadi Yu, and Guihai Chen
(2023), ``PV-TSC: Learning to Control Traffic Signals for Pedestrian and
Vehicle Traffic in 6G Era,'' \emph{IEEE Transactions on Intelligent
Transportation Systems} 24(7), 7552--7563.
\href{https://doi.org/10.1109/TITS.2022.3156816}{DOI} & Distributed RL
explicitly integrates pedestrian traffic and multi-intersection
coordination. \\
9 & Majid Raeis and Alberto Leon-Garcia (2021), ``A Deep Reinforcement
Learning Approach for Fair Traffic Signal Control,'' \emph{IEEE ITSC},
2512--2518. \href{https://doi.org/10.1109/ITSC48978.2021.9564847}{DOI} &
Makes fairness an explicit traffic-control objective and exposes the
efficiency--fairness trade-off. \\
10 & Hao Liu and Vikash V. Gayah (2023), ``Total-delay-based Max
Pressure: A Max Pressure Algorithm Considering Delay Equity,''
\emph{Transportation Research Record} 2677(6), 324--339.
\href{https://doi.org/10.1177/03611981221147051}{DOI} & Interpretable
adaptive control that improves individual vehicle delay equity; useful
efficiency baseline. \\
11 & Xinqi Du, Ziyue Li, Cheng Long, Yongheng Xing, Philip S. Yu, and
Hechang Chen (2024), ``FELight: Fairness-Aware Traffic Signal Control
via Sample-Efficient Reinforcement Learning,'' \emph{IEEE Transactions
on Knowledge and Data Engineering} 36(9), 4678--4692.
\href{https://doi.org/10.1109/TKDE.2024.3376745}{DOI} & Recent
fairness-aware RL with counterfactual augmentation, though the
counterfactual is for learning rather than human explanation. \\
12 & Kevin Riehl, Anastasios Kouvelas, and Michail A. Makridis (2026),
``Urban Priority Pass: Fair Signalised Intersection Management
Accounting for Passenger Needs through Prioritisation,''
\emph{Transportation Research Interdisciplinary Perspectives} 37,
101988. \href{https://doi.org/10.1016/j.trip.2026.101988}{DOI} & Recent
priority-negotiation mechanism accounting for passenger needs; useful
caution about entitlement and willingness-to-pay assumptions. \\
13 & Haojie Li, Haodong Hu, Ziqian Zhang, and Yingheng Zhang (2023),
``The Role of Yielding Cameras in Pedestrian--Vehicle Interactions at
Un-signalized Crosswalks: An Application of Game Theoretical Model,''
\emph{Transportation Research Part F} 92, 27--43.
\href{https://doi.org/10.1016/j.trf.2022.11.004}{DOI} & Empirical
pedestrian--driver game model combining observation and questionnaire
data. \\
14 & Fanta Camara, Patrick Dickinson, and Charles Fox (2021),
``Evaluating Pedestrian Interaction Preferences with a Game Theoretic
Autonomous Vehicle in Virtual Reality,'' \emph{Transportation Research
Part F} 78, 410--423.
\href{https://doi.org/10.1016/j.trf.2021.02.017}{DOI} & Tests a
sequential game in VR and shows how preferences can be behaviorally
estimated. \\
15 & Amir Hossein Kalantari, Yue Yang, Yee Mun Lee, Natasha Merat, and
Gustav Markkula (2023), ``Driver--Pedestrian Interactions at
Unsignalized Crossings Are Not in Line With the Nash Equilibrium,''
\emph{IEEE Access} 11, 110707--110723.
\href{https://doi.org/10.1109/ACCESS.2023.3322959}{DOI} & Critical
negative evidence: conventional Nash equilibrium poorly describes
observed paired interactions. \\
16 & Tim Miller (2019), ``Explanation in Artificial Intelligence:
Insights from the Social Sciences,'' \emph{Artificial Intelligence} 267,
1--38. \href{https://doi.org/10.1016/j.artint.2018.07.007}{DOI} &
Establishes that explanations are contrastive, selective, and social;
anchors the explanation design. \\
17 & Danding Wang, Qian Yang, Ashraf Abdul, and Brian Y. Lim (2019),
``Designing Theory-Driven User-Centric Explainable AI,'' \emph{ACM CHI},
Paper 601, 1--15. \href{https://doi.org/10.1145/3290605.3300831}{DOI} &
Maps explanation goals to theories and measurable user outcomes. \\
18 & Jonathan Dodge, Q. Vera Liao, Yunfeng Zhang, Rachel K. E. Bellamy,
and Casey Dugan (2019), ``Explaining Models: An Empirical Study of How
Explanations Impact Fairness Judgment,'' \emph{ACM IUI}, 275--285.
\href{https://doi.org/10.1145/3301275.3302310}{DOI} & Directly shows
that explanation style changes algorithmic fairness judgments. \\
19 & Na Du, Jacob Haspiel, Qiaoning Zhang, Dawn Tilbury, Anuj K.
Pradhan, X. Jessie Yang, and Lionel P. Robert Jr.~(2019), ``Look Who's
Talking Now: Implications of AV's Explanations on Driver's Trust, AV
Preference, Anxiety and Mental Workload,'' \emph{Transportation Research
Part C} 104, 428--442.
\href{https://doi.org/10.1016/j.trc.2019.05.025}{DOI} & Within-subject
simulator evidence on explanation timing, trust, preference, workload,
and anxiety. \\
20 & Zishuo Zhu, Xiaomeng Li, Patricia Delhomme, Ronald Schroeter,
Sebastien Glaser, and Andry Rakotonirainy (2025), ``Human-Centric
Explanations for Users in Automated Vehicles: A Systematic Review,''
\emph{Accident Analysis \& Prevention} 220, 108152.
\href{https://doi.org/10.1016/j.aap.2025.108152}{DOI} & Recent PRISMA
review that consolidates explanation content, timing, modality, and
study-design evidence. \\
\end{longtable}
}

\section{3. Detailed Literature
Analysis}\label{detailed-literature-analysis}

The sequence in this section is thematic rather than the rank order in
Section 2. Each profile extracts the research problem,
method/environment or data, variables, control mechanism,
game/fairness/explanation/human involvement, evaluation, findings,
limitations, and reusable design implication; the complete
presence/absence coding is made explicit in Section 4.

\subsection{3.1 Multimodal and pedestrian-aware signal
control}\label{multimodal-and-pedestrian-aware-signal-control}

\textbf{1. Wang et al.~(2023).} Research problem: allocate green time
among pedestrian and vehicle movements in a connected environment.
Method: four signal phases act as players in a Nash bargaining solution
(NBS), with a fixed 144 s cycle and green bounds of 15--45 s.
Environment/data: synthetic SUMO intersection experiments. Variables:
pedestrian/vehicle queues, saturation flows, green time, demand, and
phase utility. Mechanism: compute an NBS allocation and seek queue
equalization. Game theory: yes, cooperative bargaining. Fairness:
implicit through phase bargaining and queue balancing rather than
separately validated distributive/procedural fairness.
Explainability/HCI: no human-facing explanation and no user study.
Evaluation: delay/queue-related simulation outcomes against an actuated
controller. Main finding: bargaining can improve multimodal timing under
tested demand patterns. Limitation: fixed cycle and synthetic
assumptions; the rational-player semantics attach to phases, not actual
human behavior; no efficiency-loss bound, explanation fidelity test,
perceived fairness, or trust calibration. Reuse: mandatory closest
baseline and the clearest evidence that ``using Nash bargaining for
pedestrian and vehicle traffic'' cannot be the novelty
claim.\hyperref[source-1]{[1]}

\textbf{2. He et al.~(2014).} Problem: serve competing requests from
transit, passenger vehicles, and pedestrians while retaining
coordination. Method: request-based mixed-integer signal control
combining priority, actuation, and coordination. Environment: VISSIM
with ASC/3 controller-in-the-loop. Variables include estimated arrivals,
requested service, phase timing, coordination constraints, and priority
classes. Mechanism: rolling optimization under signal constraints; not
game theoretic. Fairness: priority policy, not a formal equity
construct. Explainability/HCI: absent. Evaluation: simulation and
controller-in-the-loop operational measures. Finding: a common
optimization framework can combine heterogeneous requests. Limitations:
communication/arrival assumptions, computation and calibration demands,
and no perceived legitimacy assessment. Reuse: service-request
abstraction, constraint modeling, and an optimization
comparator.\hyperref[source-4]{[4]}

\textbf{3. Ma et al.~(2015).} Problem: choose a pedestrian phase pattern
as well as signal timings at an isolated intersection. Method: compare
exclusive pedestrian phases and two-way crossing patterns through an
optimization that monetizes pedestrian delay/detour, vehicle delay,
pedestrian--vehicle exposure, and potential noncompliance. Environment:
case-based numerical optimization. Variables: flows, phase durations,
crossing distance, conflict exposure, delays, and behavioral costs.
Mechanism: phase-pattern and timing optimization; no game theory.
Fairness: not explicit, although pedestrian costs receive first-class
representation. Explainability/HCI: absent. Evaluation: operational and
safety-related objective values across demand conditions. Finding: no
one pedestrian phase pattern dominates across all demand. Limitations:
parameterized behavioral cost and case-study transferability. Reuse:
evidence for a constrained action set and for testing multiple
pedestrian-demand regimes.\hyperref[source-5]{[5]}

\textbf{4. Niels et al.~(2024 issue; online 2023).} Problem: coordinate
connected automated vehicles and pedestrians without conventional fixed
phases. Method: rolling-horizon mixed-integer linear programming assigns
conflict-free space--time slots. Environment: Aimsun Next connected to
Python/Gurobi; reported 60 min simulations with 10 min warm-up and five
seeds. Variables: vehicle trajectories/arrival times, pedestrian calls,
safety gaps, waiting constraints, and class weights. Mechanism:
optimization, not a strategic human game. Fairness: bounded pedestrian
waiting and weight choices, not a human-perception construct.
Explainability/HCI: absent. Evaluation: delay and throughput across
penetration/demand settings. Finding: optimization can safely coordinate
both modes when connectivity assumptions hold. Limitations: CAV
dependence, solver cost, and few replications for strong stochastic
inference. Reuse: hard safety constraints, rolling horizon, and maximum
waiting as a non-negotiable design requirement.\hyperref[source-6]{[6]}

\textbf{5. Yazdani et al.~(2023).} Problem: optimize signal control for
interacting vehicles and pedestrians, including realistic pedestrian
behavior. Method: deep reinforcement learning with a total road-user
delay objective. Environment/data: a Melbourne intersection
microsimulation calibrated with video, loop detector, and SCATS-related
parameters; fully actuated comparison. Variables: vehicle and pedestrian
flows/waits, crossing interaction, and jaywalking behavior. Mechanism:
learned policy; no game theory. Fairness: broadly equitable multimodal
service in motivation, but not a formal allocative-fairness construct
with human validation. Explainability/HCI: policy remains opaque; no
explanation study. Evaluation: simulated delay and operational
performance. Finding: a pedestrian-aware learned controller can improve
total user delay in its calibrated setting. Limitations: reward and
calibration transfer, opacity, and no perceived fairness/trust evidence.
Reuse: person-centered efficiency metric and calibration
expectations.\hyperref[source-7]{[7]}

\textbf{6. Xu et al.~(2023).} Problem: scale adaptive control to mixed
pedestrian and vehicle traffic using anticipated high-precision
connectivity. Method: distributed reinforcement learning, with
pedestrian access integrated into multi-intersection control.
Environment: traffic simulation across network settings. Variables:
pedestrian presence/flow, vehicle queues/traffic states, phases, and
coordination information. Mechanism: distributed learned agents; not
bargaining. Fairness: pedestrian inclusion, not a defined fairness
measure. Explainability/HCI: absent. Evaluation: classic control
baselines and ablated policy variants. Finding: explicitly modeling
pedestrian traffic can improve mixed-flow control relative to
vehicle-centric policies. Limitations: reliance on 6G-style
localization/tracking assumptions and opaque policy behavior. Reuse:
evidence that pedestrian inclusion must be evaluated, not treated as an
interface-only feature.\hyperref[source-8]{[8]}

\subsection{3.2 Fairness and priority in traffic
control}\label{fairness-and-priority-in-traffic-control}

\textbf{7. Raeis and Leon-Garcia (2021).} Problem: conventional
traffic-signal RL can distribute delay inequitably while optimizing
throughput. Method: fairness-aware Double DQN. Environment: SUMO single
four-way intersection under Poisson, Markov-modulated Poisson, and
nonhomogeneous Poisson arrivals. Variables: queue/wait/traffic state and
fairness-augmented reward. Mechanism: learned action value; no
pedestrian agent and no bargaining. Fairness: traffic delay/throughput
distribution across movements. Explainability/HCI: absent. Evaluation:
efficiency and fairness under multiple arrival processes. Finding:
fairness can be improved but trades against efficiency depending on
reward design. Limitations: vehicle-only, synthetic, opaque, and
sensitive to scalar reward weights. Reuse: motivates an explicit
Pareto/efficiency bound rather than an arbitrary weighted
sum.\hyperref[source-9]{[9]}

\textbf{8. Liu and Gayah (2023).} Problem: standard max-pressure is
throughput-optimal but can leave large individual delay disparities.
Method: total-delay-based pressure replaces queue-length pressure.
Environment: SUMO at a single intersection and a 4×4 grid, including
partial connected-vehicle information. Variables: cumulative individual
delays, movement demand, downstream state, and phase pressure.
Mechanism: interpretable adaptive heuristic, no game. Fairness: delay
equity among vehicles. Explainability/HCI: no public-facing explanation.
Evaluation: total delay and distributional outcomes. Finding:
delay-based pressure can improve equity with competitive operational
performance. Limitations: no pedestrians, perceived fairness, or
explanation; state observability matters. Reuse: preferred
efficiency-oriented adaptive baseline because its score is transparent
and reproducible.\hyperref[source-10]{[10]}

\textbf{9. Du et al.~(2024; FELight).} Problem: fairness-aware signal RL
is often sample-inefficient and may over-activate fairness
interventions. Method: fairness activation plus counterfactual
augmentation and self-supervised state representation. Environment/data:
benchmark traffic datasets/simulated networks. Variables: traffic state,
vehicle latency, activation threshold, and learned representation.
Mechanism: RL. Fairness: vehicle scheduling/latency fairness.
Explainability/HCI: the counterfactual is a training augmentation, not
an explanation to a person; no user study. Evaluation: reward, fairness,
and sample efficiency against learning baselines. Finding: selectively
learning fairness-relevant cases can improve sample efficiency and
performance. Limitations: black-box policy, vehicle-centric equity,
threshold dependence. Reuse: distinguishes machine-learning
counterfactual augmentation from human-consumable counterfactual
explanation.\hyperref[source-11]{[11]}

\textbf{10. Riehl et al.~(2026).} Problem: ordinary signals ignore
heterogeneous passenger urgency and social value. Method:
priority-pass/auction-like allocation integrated with traffic signal
management. Environment: Manhattan-like microscopic simulations.
Variables: traffic pressure, passenger need/priority, value-of-time or
entitlement parameters, and bids/credits. Mechanism: prioritization and
allocation rather than pedestrian--vehicle bargaining. Fairness:
articulated in terms of access to priority and passenger needs.
Explainability/HCI: no evaluation of public explanations. Evaluation:
delay and allocation outcomes under different priority regimes; reported
improvements are simulation-conditional. Limitations: identification,
truthful preference/need declaration, distributive legitimacy, privacy,
and possible pay-to-win effects. Reuse: warning that individual bidding
is both unnecessary and ethically fraught for a public crossing
controller.\hyperref[source-12]{[12]}

\subsection{3.3 Explainable signal
control}\label{explainable-signal-control}

\textbf{11. Rizzo et al.~(2019).} Problem: RL signal controllers are
hard to inspect. Method: a policy-gradient controller is analyzed with
SHAP feature attributions. Environment/data: a signalized roundabout, 11
phases, and real traffic data used in simulation. Variables:
queue/traffic features, selected phase, learned policy outputs, and SHAP
values. Mechanism: RL plus post-hoc attribution. Game/fairness: neither.
Explainability: global/local feature influence on selected signal
phases. Humans: no controlled user study. Evaluation: traffic efficiency
and technical interpretability examples. Finding: SHAP can expose which
features influenced a learned phase selection. Limitations: attribution
is not automatically a causal reason, no pedestrian equity, fidelity can
be model/perturbation dependent, and no comprehension/trust test. Reuse:
direct benchmark demonstrating that technical XAI alone is not a
human-centered contribution.\hyperref[source-2]{[2]}

\textbf{12. Liu et al.~(2024).} Problem: obtain high-performing signal
controllers whose logic is directly readable. Method: genetic
programming evolves symbolic phase-score expressions. Environment: SUMO
synthetic grids and a Berlin network. Variables: queues/traffic
descriptors used in evolved expressions. Mechanism: white-box symbolic
rule. Fairness/pedestrians: not central. Explainability: intrinsic rule
inspection. Humans: no interface/user study. Evaluation:
delay/throughput plus expression complexity/interpretability
comparisons. Finding: compact rules can remain competitive in tested
networks. Limitations: rule readability is treated as interpretability
without demonstrating stakeholder comprehension; no pedestrian
allocation or perceived fairness. Reuse: establishes an
intrinsic-interpretability alternative and supports keeping the proposed
allocator transparent.\hyperref[source-3]{[3]}

\subsection{3.4 Behavioral game theory for pedestrian--vehicle
interaction}\label{behavioral-game-theory-for-pedestrianvehicle-interaction}

\textbf{13. Li et al.~(2023).} Problem: understand how yielding cameras
and contextual factors change driver--pedestrian yielding at
unsignalized crossings. Method: roadside/aerial observations plus
questionnaires and a quantal-response-equilibrium model. Variables:
demographics, behavior, environment, perceived probabilities/beliefs,
and camera condition. Mechanism: behavioral game.
Fairness/explainability: not an allocation controller or explanation
interface. Humans/data: real observed interactions and self-report.
Evaluation: model fit and behavioral effects. Finding: yielding
technology and beliefs affect interaction, with heterogeneous behavior.
Limitation: unsignalized dyadic encounters do not transfer directly to
group signal scheduling. Reuse: evidence to use bounded/behavioral
models only if describing humans; the proposed group allocator instead
uses game theory normatively.\hyperref[source-13]{[13]}

\textbf{14. Camara et al.~(2021).} Problem: model pedestrian preferences
when interacting with a game-theoretic autonomous vehicle. Method:
sequential chicken-game policy evaluated in VR; Gaussian-process
inference estimates preference parameters. Variables: time-to-arrival,
pedestrian/vehicle actions, collision cost, and delay/preference terms.
Mechanism: sequential game. Fairness: not group allocative fairness.
Explainability: not evaluated. Humans: VR participant experiments.
Evaluation: crossing choices and inferred utilities. Finding: collision
avoidance dominates and participant preferences are heterogeneous.
Limitations: small experiments, VR validity, dyadic AV interaction, and
inferred utility sensitivity. Reuse: shows that utility terms require
empirical estimation rather than convenient
coefficients.\hyperref[source-14]{[14]}

\textbf{15. Kalantari et al.~(2023).} Problem: test whether established
game-theoretic models actually describe driver--pedestrian behavior.
Method: connected Cave Automatic Virtual Environment pedestrian
simulator plus motion-base driving simulator; 32 pairs; four game models
compared. Variables: kinematics, gap acceptance, actions, and behavioral
payoff parameters. Mechanism: empirical model comparison.
Fairness/explainability: neither. Humans: direct paired interaction.
Evaluation: predictive fit to observed behavior. Finding: conventional
Nash-equilibrium predictions were inconsistent with interactions; a
behavioral game formulation fit substantially better. Limitation:
unsignalized, laboratory dyads. Reuse: decisive reason not to claim that
real pedestrians and drivers are rational Nash
agents.\hyperref[source-15]{[15]}

\subsection{3.5 Human-centered explanations, fairness, and
trust}\label{human-centered-explanations-fairness-and-trust}

\textbf{16. Miller (2019).} Problem: AI explanations often ignore
findings about how people ask for and process explanations. Method:
integrative review of philosophy, cognitive psychology, and social
psychology. Data: prior empirical/theoretical literature.
Variables/mechanism: conceptual, not a controller. Explainability:
explanations should be contrastive, selective, and socially situated;
causes people regard as abnormal are often salient. Evaluation:
synthesis rather than a new experiment. Limitation: general guidance
does not select domain-specific content. Reuse: supports reporting the
chosen action, the salient reason, the rejected alternative, and a valid
contrastive threshold.\hyperref[source-16]{[16]}

\textbf{17. Wang et al.~(2019).} Problem: connect XAI interface design
to user goals and behavioral theories. Method: conceptual framework
grounded in explanation and HCI theory, illustrated with design
considerations. Explainability/HCI: central. Evaluation: design
framework rather than a domain deployment. Main finding: explanation
design must begin with user questions and measurable objectives, not
whichever algorithmic output is convenient. Limitation: no
traffic-domain validation. Reuse: maps ``why this priority?'' and ``what
would change it?'' to separate comprehension
outcomes.\hyperref[source-17]{[17]}

\textbf{18. Dodge et al.~(2019).} Problem: determine how explanations
influence fairness judgments of algorithmic decisions. Method:
controlled online experiments with more than 160 participants and four
explanation approaches. Variables: explanation type, model behavior,
fairness judgment, and task context. Finding: explanation form changes
perceived fairness and can surface different fairness concerns; more
explanation is not universally better. Limitation: non-traffic tasks and
crowd-worker sampling. Reuse: direct precedent for manipulating
explanation format and measuring fairness perception separately from
objective fairness.\hyperref[source-18]{[18]}

\textbf{19. Du et al.~(2019).} Problem: determine whether and when AV
explanations influence trust and acceptance. Method: within-subject
driving-simulator experiment with 32 participants and four conditions:
no explanation, pre-action, post-action, and explanation plus approval.
Variables: explanation timing/control, trust, AV preference, anxiety,
and mental workload. Finding: pre-action explanations were associated
with higher trust and preference; anxiety and workload did not differ.
Limitations: small student-heavy sample, in-vehicle actions rather than
public signal allocation, and trust increase rather than calibration.
Reuse: timing matters; for a smart-intersection audit interface, the
explanation should appear with the decision, while detailed
counterfactual information can be inspected after the safety-critical
signal indication.\hyperref[source-19]{[19]}

\textbf{20. Zhu et al.~(2025).} Problem: consolidate evidence on
human-centered explanations for automated-vehicle users. Method:
PRISMA-style systematic review across five databases; 59 studies met
stated inclusion criteria. Variables: explanation content, modality,
timing, context, and user outcomes. Findings: reasons commonly
outperform bare action descriptions; pre-action explanations often help;
multimodal detail is beneficial only when complementary and can overload
users. Limitations: predominance of simulation studies, heterogeneous
measures, and little infrastructure-level allocation. Reuse: supports
concise factual explanations, measured rather than assumed benefits, and
a scoped comparison of no explanation, reason, and counterfactual
reason.\hyperref[source-20]{[20]}

\section{4. Literature Comparison
Matrix}\label{literature-comparison-matrix}

Legend: \textbf{Y} = central and explicit; \textbf{P} =
partial/indirect; \textbf{N} = absent. ``Real data'' means
observational, survey, or calibrated field data---not necessarily
real-world deployment.

\begin{landscape}
{\def\LTcaptype{none} % do not increment counter
\scriptsize
\setlength{\tabcolsep}{2pt}
\begin{longtable}[]{@{}
  >{\raggedleft\arraybackslash}p{(\linewidth - 26\tabcolsep) * \real{0.025}}
  >{\raggedright\arraybackslash}p{(\linewidth - 26\tabcolsep) * \real{0.080}}
  >{\centering\arraybackslash}p{(\linewidth - 26\tabcolsep) * \real{0.035}}
  >{\centering\arraybackslash}p{(\linewidth - 26\tabcolsep) * \real{0.065}}
  >{\centering\arraybackslash}p{(\linewidth - 26\tabcolsep) * \real{0.035}}
  >{\centering\arraybackslash}p{(\linewidth - 26\tabcolsep) * \real{0.065}}
  >{\centering\arraybackslash}p{(\linewidth - 26\tabcolsep) * \real{0.060}}
  >{\centering\arraybackslash}p{(\linewidth - 26\tabcolsep) * \real{0.080}}
  >{\centering\arraybackslash}p{(\linewidth - 26\tabcolsep) * \real{0.050}}
  >{\centering\arraybackslash}p{(\linewidth - 26\tabcolsep) * \real{0.040}}
  >{\centering\arraybackslash}p{(\linewidth - 26\tabcolsep) * \real{0.075}}
  >{\centering\arraybackslash}p{(\linewidth - 26\tabcolsep) * \real{0.055}}
  >{\raggedright\arraybackslash}p{(\linewidth - 26\tabcolsep) * \real{0.225}}
  >{\raggedright\arraybackslash}p{(\linewidth - 26\tabcolsep) * \real{0.110}}@{}}
\toprule\noalign{}
\begin{minipage}[b]{\linewidth}\raggedleft
\#
\end{minipage} & \begin{minipage}[b]{\linewidth}\raggedright
Paper
\end{minipage} & \begin{minipage}[b]{\linewidth}\centering
P--V
\end{minipage} & \begin{minipage}[b]{\linewidth}\centering
Smart intersection
\end{minipage} & \begin{minipage}[b]{\linewidth}\centering
Game
\end{minipage} & \begin{minipage}[b]{\linewidth}\centering
Multi-agent
\end{minipage} & \begin{minipage}[b]{\linewidth}\centering
Fairness
\end{minipage} & \begin{minipage}[b]{\linewidth}\centering
Explanation
\end{minipage} & \begin{minipage}[b]{\linewidth}\centering
HCI/user
\end{minipage} & \begin{minipage}[b]{\linewidth}\centering
Trust
\end{minipage} & \begin{minipage}[b]{\linewidth}\centering
Simulation
\end{minipage} & \begin{minipage}[b]{\linewidth}\centering
Real data
\end{minipage} & \begin{minipage}[b]{\linewidth}\raggedright
Main limitation
\end{minipage} & \begin{minipage}[b]{\linewidth}\raggedright
Relevance
\end{minipage} \\
\midrule\noalign{}
\endhead
\bottomrule\noalign{}
\endlastfoot
1 & Wang 2023 & Y & Y & Y & Y & P & N & N & N & SUMO & N & No
human-facing transparency or user validation & Closest method \\
2 & Rizzo 2019 & N & Y & N & N & N & Y & N & N & Traffic sim. & Y & SHAP
interpretation only; no people/pedestrians & Closest traffic XAI \\
3 & Liu 2024 GP & N & Y & N & N & N & Y & N & N & SUMO & P & Readability
assumed, not tested & White-box comparator \\
4 & He 2014 & Y & Y & N & P & P & N & N & N & VISSIM/SIL & P & Priority
policy, not perceived legitimacy & Multimodal constraints \\
5 & Ma 2015 & Y & Y & N & N & P & N & N & N & Numerical & Case &
Static/calibrated behavioral costs & Phase design \\
6 & Niels 2024 & Y & Y & N & P & P & N & N & N & Aimsun & N &
High-connectivity/solver assumptions & Safety/wait bounds \\
7 & Yazdani 2023 & Y & Y & N & P & P & N & N & N & Calibrated sim. & Y &
Opaque DRL; fairness not formal & Person-delay benchmark \\
8 & Xu 2023 PV-TSC & Y & Y & N & Y & P & N & N & N & Traffic sim. & N &
6G localization and black-box assumptions & Pedestrian-aware RL \\
9 & Raeis 2021 & N & Y & N & P & Y & N & N & N & SUMO & N &
Vehicle-only, reward-weight sensitivity & Fairness baseline \\
10 & Liu \& Gayah 2023 & N & Y & N & P & Y & P & N & N & SUMO & P &
Vehicle-only; no user meaning of equity & Efficiency baseline \\
11 & FELight 2024 & N & Y & N & P & Y & N & N & N & Traffic sim. & P &
``Counterfactual'' is training data, not explanation & Fair-RL
comparator \\
12 & Riehl 2026 & N & Y & P & Y & Y & N & N & N & Microscopic & N &
Entitlement, privacy, truthful-report assumptions & Priority-allocation
caution \\
13 & Li 2023 & Y & N & Y & Y & N & N & Y & P & N & Y & Unsignalized
dyads & Behavioral game evidence \\
14 & Camara 2021 & Y & N & Y & Y & N & N & Y & P & VR & Y & Small VR
dyads & Utility elicitation \\
15 & Kalantari 2023 & Y & N & Y & Y & N & N & Y & N & Coupled simulators
& Y & Unsignalized lab setting & Reject naive Nash realism \\
16 & Miller 2019 & N & N & N & N & P & Y & P & P & N & Review & No
domain implementation & Explanation theory \\
17 & Wang 2019 CHI & N & N & N & N & P & Y & Y & P & N &
Review/ framework & No traffic validation & XAI design framework \\
18 & Dodge 2019 & N & N & N & N & Y & Y & Y & P & Online experiment & Y
& Non-traffic task & Fairness perception \\
19 & Du 2019 & P & N & N & N & N & Y & Y & Y & Driving sim. & Y &
In-vehicle, small sample, raw trust & Explanation timing \\
20 & Zhu 2025 & P & N & N & N & P & Y & Y & Y & Systematic review & Y &
Heterogeneous AV studies & Explanation evidence \\
\end{longtable}
}
\end{landscape}

\subsection{Answers to the ten synthesis
questions}\label{answers-to-the-ten-synthesis-questions}

\textbf{Q1. How is priority determined?} Fixed cycles and actuated calls
remain common. Research controllers use mixed-integer optimization,
pressure/delay scores, learned action values, request priorities, or
bargaining among phases. Pedestrian service is often a constraint or
reward term, not an equal public stakeholder. The closest bargaining
work allocates green by phase utilities; priority-pass work allocates by
passenger need/entitlement.

\textbf{Q2. Which variables recur?} Vehicle queue/occupancy, arrival
flow, cumulative or average delay, downstream capacity, phase elapsed
time, and switching cost dominate. Pedestrian-aware work adds
counts/calls, arrival rate, crossing occupancy, average/maximum waiting,
walking speed, crossing distance, and interaction/noncompliance risk.
Emergency status and sensor confidence appear less consistently and
should be modeled explicitly here.

\textbf{Q3. How is fairness defined?} There is no single accepted
definition. Traffic papers use queue equalization, disparity in
individual delays, Jain-style equality, bounded waiting, priority
access, or fairness penalties in a reward. These are objective
distributions. Algorithmic-fairness perception studies measure
procedural, distributive, or informational justice. The two levels must
be reported separately.\hyperref[source-21]{[21]}

\textbf{Q4. Are pedestrians and vehicles competing agents?} Sometimes,
but usually not as literal human agents. Wang et al.~use signal phases
as bargaining players. Behavioral crossing papers model a pedestrian and
driver/AV as strategic actors. Many control papers instead aggregate
modes into requests or reward components. This project should call its
participants ``virtual mode advocates'' to avoid an empirically false
rational-human claim.

\textbf{Q5. Is game theory useful?} Yes as a transparent
\textbf{normative allocation rule}: weighted Nash bargaining encodes
proportional fairness and exposes utilities and trade-offs. It is not
justified as a descriptive model of real human cognition; Kalantari et
al.~directly challenge conventional Nash predictions. Game theory adds
value only if it yields a clearer fairness contract than a simpler
weighted score and is compared with that score.

\textbf{Q6. How do traffic systems explain decisions?} Usually they do
not. Rizzo et al.~add post-hoc feature attribution to RL; Liu et
al.~evolve readable symbolic rules. Neither evaluates whether citizens
or operators understand a particular priority allocation. Most
controllers expose internal metrics only to researchers.

\textbf{Q7. Which transportation explanation techniques exist?} Feature
attribution, inherently interpretable rules, ``what/why''
natural-language reasons, pre- versus post-action timing, and
occasionally visual/audio modalities. The literature supports selective,
contrastive reasons; it does not justify dumping every sensor value or
claiming that a counterfactual generated by a surrogate is
causal.\hyperref[source-16]{[16]}\,\hyperref[source-20]{[20]}\,\hyperref[source-22]{[22]}

\textbf{Q8. How are trust and fairness evaluated?} Traffic-control
papers mostly do not evaluate either with people. Automated-driving work
uses validated trust questionnaires, preference/acceptance,
workload/anxiety, and behavioral reliance. Algorithmic-decision work
uses perceived justice/fairness ratings and qualitative rationales. The
correct target is calibrated trust: higher confidence in sound decisions
and lower confidence in anomalous ones.\hyperref[source-23]{[23]}

\textbf{Q9. Which interfaces exist?} Simulation dashboards and
engineering consoles show phases, flows, or model features; AV studies
use in-vehicle visual/audio messages and VR. The literature does not
establish a standard citizen-facing interface for infrastructure
allocation. The proposed dashboard should therefore be framed as a
public/operator audit-and-replay artifact, never as a replacement for
legal WALK/DON'T WALK or traffic-light indications.

\textbf{Q10. What gaps remain?} No core paper in the set jointly
supplies: (i) a formal pedestrian--vehicle fairness contract with
starvation and bounded efficiency loss; (ii) faithful per-decision
factual and counterfactual explanations; (iii) an audit interface; and
(iv) a causal user evaluation of comprehension, perceived fairness, and
calibrated trust. The gap is this verified combination and its
evaluation---not the mere co-occurrence of fashionable terms.

\section{5. Research Gap}\label{research-gap}

{\def\LTcaptype{none} % do not increment counter
\begin{longtable}[]{@{}
  >{\raggedright\arraybackslash}p{(\linewidth - 6\tabcolsep) * \real{0.2500}}
  >{\raggedright\arraybackslash}p{(\linewidth - 6\tabcolsep) * \real{0.2500}}
  >{\raggedright\arraybackslash}p{(\linewidth - 6\tabcolsep) * \real{0.2500}}
  >{\raggedright\arraybackslash}p{(\linewidth - 6\tabcolsep) * \real{0.2500}}@{}}
\toprule\noalign{}
\begin{minipage}[b]{\linewidth}\raggedright
Layer
\end{minipage} & \begin{minipage}[b]{\linewidth}\raggedright
Established
\end{minipage} & \begin{minipage}[b]{\linewidth}\raggedright
Underexplored
\end{minipage} & \begin{minipage}[b]{\linewidth}\raggedright
Defensible opportunity
\end{minipage} \\
\midrule\noalign{}
\endhead
\bottomrule\noalign{}
\endlastfoot
Control & Pedestrian-aware optimization and RL can reduce mixed-mode
delay.\hyperref[source-4]{[4]}\,\hyperref[source-5]{[5]}\,\hyperref[source-6]{[6]}\,\hyperref[source-7]{[7]}\,\hyperref[source-8]{[8]}
& Simple, transparent controllers with explicit efficiency-loss
guarantees & Constrain candidate plans to a near-efficient set before
fairness selection \\
Negotiation & Nash bargaining has already allocated pedestrian/vehicle
phase green.\hyperref[source-1]{[1]} & Robust bargaining with hard
starvation constraints, direction-level pressure, uncertainty, and
sensitivity analysis & Weighted mode-advocate bargaining as a normative
public-service rule \\
Fairness & Delay equity, queue balance, Jain indices, and fairness
rewards
exist.\hyperref[source-9]{[9]}\,\hyperref[source-10]{[10]}\,\hyperref[source-11]{[11]}
& Separation of objective, procedural, and perceived fairness & Dual
evaluation: traffic distributions plus human justice judgments \\
Explanation & SHAP and symbolic traffic rules exist; AV explanations
affect
attitudes.\hyperref[source-2]{[2]}\,\hyperref[source-3]{[3]}\,\hyperref[source-19]{[19]}\,\hyperref[source-20]{[20]}
& Faithful instance-level reasons tied to binding traffic constraints
and genuine decision flips & Typed decision trace plus replay-validated
counterfactual contract \\
HCI & Explanation format/timing affects trust and fairness
judgments.\hyperref[source-18]{[18]}\,\hyperref[source-19]{[19]}\,\hyperref[source-20]{[20]}
& Infrastructure-level priority decisions; comprehension and trust
calibration & Controlled audit/replay dashboard experiment with sound
and anomalous decisions \\
Evidence & Simulation and human studies exist separately & Cross-layer
causal chain from controller state to explanation to judgment & Release
scenarios, trace records, stimuli, analysis code, and fidelity tests
together \\
\end{longtable}
}

The integrated combination appears genuinely underinvestigated in the
reviewed literature. That statement is narrower and more supportable
than ``no prior work combines game theory and explanations in smart
intersections.'' A later paper could always exist outside the search
set, so the manuscript should phrase novelty as ``to our knowledge,
among the reviewed peer-reviewed literature'' and update the search
immediately before submission.

The closest competing work is Wang et al.~Its method and this proposal
overlap in multimodal signal control, SUMO, phase/mode utilities, and
Nash bargaining. They differ as follows:

{\def\LTcaptype{none} % do not increment counter
\begin{longtable}[]{@{}
  >{\raggedright\arraybackslash}p{(\linewidth - 4\tabcolsep) * \real{0.3333}}
  >{\raggedright\arraybackslash}p{(\linewidth - 4\tabcolsep) * \real{0.3333}}
  >{\raggedright\arraybackslash}p{(\linewidth - 4\tabcolsep) * \real{0.3333}}@{}}
\toprule\noalign{}
\begin{minipage}[b]{\linewidth}\raggedright
Dimension
\end{minipage} & \begin{minipage}[b]{\linewidth}\raggedright
Wang et al.~2023
\end{minipage} & \begin{minipage}[b]{\linewidth}\raggedright
Proposed work
\end{minipage} \\
\midrule\noalign{}
\endhead
\bottomrule\noalign{}
\endlastfoot
Player interpretation & Four traffic-signal phases & Two virtual mode
advocates at the top level; direction choice delegated to vehicle
pressure \\
Cycle/action & Fixed cycle, bounded phase green allocation &
Receding-horizon safe candidate plans; act only at admissible transition
points \\
Fairness & Queue equalization/bargaining & Hard maximum wait +
waiting-debt weights + proportional fairness + near-efficiency bound \\
Safety/robustness & Signal constraints in simulation & Lexicographic
safety/emergency layer, sensor-quality modes, auditable fallbacks \\
Explanation & None & Faithful trace, rejected alternative, and
solver-validated counterfactual \\
Human evidence & None & Preregistered comprehension, fairness,
usability, and trust-calibration experiment \\
Claim & Operational multimodal control & Cross-layer fairness and
explanation evaluation; only incremental control novelty \\
\end{longtable}
}

\section{6. Novelty Opportunities}\label{novelty-opportunities}

\begin{enumerate}
\def\labelenumi{\arabic{enumi}.}
\tightlist
\item
  \textbf{Bounded-efficiency fair bargaining (incremental algorithmic
  novelty).} Rather than optimize a weighted fairness/efficiency sum
  with arbitrary coefficients, first filter to plans whose efficiency
  loss is within a calibrated tolerance \(\delta\), then maximize a
  waiting-debt-weighted Nash product subject to maximum-wait and
  signal-safety constraints. Evidence must include Pareto and
  sensitivity analyses; otherwise this is merely another hand-tuned
  utility.
\item
  \textbf{Decision-trace explanation contract (technical/explainability
  novelty).} Treat explanation as a verified projection of the actual
  solver record. Every numerical reason maps to a timestamped state
  field; every claimed constraint was binding; every counterfactual is
  replayed through the same controller and actually flips the decision.
  This is stronger than a decorative rationale and easier to audit than
  post-hoc SHAP.
\item
  \textbf{Causal HCI evaluation of explanation depth (evaluation
  novelty).} Hold algorithm, scenario, visual layout, and decision
  constant while varying decision-only versus factual versus
  factual-plus-counterfactual content. Include sound and anomalous
  outcomes so the target is discrimination/calibration rather than
  maximum trust.
\item
  \textbf{Cross-layer reproducibility package (engineering/evaluation
  contribution).} Release controller configs, paired random seeds, raw
  decision traces, generated stimuli, fidelity tests, and analysis
  scripts. The contribution is the benchmark linkage across traffic and
  human outcomes, not the use of React, FastAPI, or SUMO.
\end{enumerate}

The first opportunity is publishable as an algorithmic contribution only
if it outperforms or offers a measurably better fairness contract than
simpler baselines. Opportunities 2 and 3 are the safer central thesis.

\section{7. Final Proposed Research
Contribution}\label{final-proposed-research-contribution}

The proposed research introduces and evaluates \textbf{Explainable Fair
Priority Bargaining (EFPB)}, a transparent intersection controller and
audit interface. EFPB uses a lexicographic decision pipeline: legal
signal safety and active emergency handling; maximum-wait/starvation
protection; an empirically calibrated near-efficient candidate set; then
weighted Nash bargaining between pedestrian and vehicle service.
Vehicle-direction service is resolved by a delay-based max-pressure
rule. Each decision produces a machine-verifiable trace that feeds
factual and contrastive explanation templates. A paired stochastic
simulation study tests traffic performance and a preregistered
repeated-measures experiment tests whether explanation depth improves
accurate understanding and calibrated judgment.

Publication-style contribution statement:

\begin{quote}
This work contributes (1) a transparent, starvation-safe
weighted-bargaining formulation for multimodal intersection service with
an explicit bound on efficiency loss; (2) a faithful
trace-to-explanation architecture that reports selected causes, binding
constraints, rejected alternatives, and validated nearest decision-flip
counterfactuals; and (3) a controlled human-subject evaluation
separating objective allocation fairness from perceived procedural
fairness and distinguishing appropriate trust calibration from
indiscriminate trust increase.
\end{quote}

{\def\LTcaptype{none} % do not increment counter
\begin{longtable}[]{@{}
  >{\raggedright\arraybackslash}p{(\linewidth - 4\tabcolsep) * \real{0.3333}}
  >{\raggedright\arraybackslash}p{(\linewidth - 4\tabcolsep) * \real{0.3333}}
  >{\raggedright\arraybackslash}p{(\linewidth - 4\tabcolsep) * \real{0.3333}}@{}}
\toprule\noalign{}
\begin{minipage}[b]{\linewidth}\raggedright
Contribution class
\end{minipage} & \begin{minipage}[b]{\linewidth}\raggedright
Proposed contribution
\end{minipage} & \begin{minipage}[b]{\linewidth}\raggedright
Novelty status
\end{minipage} \\
\midrule\noalign{}
\endhead
\bottomrule\noalign{}
\endlastfoot
Algorithmic & Safety/starvation-constrained, near-efficient weighted
bargaining with temporal service debt & Incremental relative to existing
pedestrian--vehicle Nash bargaining; must earn novelty through
guarantees/ablations \\
Technical/XAI & Typed decision trace and deterministic
factual/alternative/counterfactual explanation with replay validation &
Potentially novel integration and fidelity contract \\
HCI & Explanation-depth experiment targeting objective mental-model
performance and calibrated judgment & Strong potential novelty in
infrastructure-level priority allocation \\
ITS & A transparent multimodal controller evaluated under demand,
accessibility, emergency, and sensor-failure scenarios &
Application/integration contribution unless expanded beyond one
intersection \\
Experimental & Joined paired-seed traffic benchmark and fixed-replay
human benchmark linked by immutable decision IDs & Potentially novel
cross-layer evaluation \\
Engineering & SUMO/TraCI, FastAPI, React, WebSocket, Parquet, templates
& Enabling implementation, not novelty by itself \\
\end{longtable}
}

Claims to avoid: ``the first game-theoretic pedestrian--vehicle
controller,'' ``humans behave as Nash-rational agents,'' ``explanations
cause trust'' without a controlled study, ``equal Jain score means
fair,'' or ``SUMO results prove field safety.''

\section{8. Research Questions}\label{research-questions}

\textbf{RQ1 --- Allocation.} Compared with fixed-time and
efficiency-only adaptive control, can EFPB reduce severe
pedestrian/vehicle waiting disparities and starvation violations while
keeping total person-delay degradation within its declared efficiency
tolerance?

\textbf{RQ2 --- Robustness.} How sensitive are efficiency and fairness
outcomes to the efficiency tolerance, waiting-debt function, demand mix,
mobility-impaired walking speed, sensor error, and emergency demand?

\textbf{RQ3 --- Explanation comprehension.} Does a faithful factual
explanation improve people's ability to identify why an intersection
made a priority decision, compared with showing the decision and state
alone?

\textbf{RQ4 --- Contrastive understanding.} Does adding a
solver-validated counterfactual improve understanding of what would have
changed the priority decision beyond a factual explanation, and at what
time or workload cost?

\textbf{RQ5 --- Judgment calibration.} Do explanations help people
distinguish defensible priority decisions from injected anomalous
decisions, producing better calibrated trust and fairness judgments
rather than uniformly higher trust?

RQ1--RQ2 support the ITS/control claim; RQ3--RQ5 support the HCI/XAI
claim. No single metric should be used to answer both.

\section{9. Hypotheses}\label{hypotheses}

The HCI hypotheses refer to three interface conditions: \texttt{D}
(decision and state only), \texttt{F} (decision, state, and factual
explanation), and \texttt{FC} (the same plus a nearest feasible
counterfactual). ``Sound'' decisions are genuine EFPB outputs;
``anomalous'' decisions are carefully constructed constraint/priority
violations used only as study stimuli.

\begin{itemize}
\tightlist
\item
  \textbf{H1:} F and FC will increase objective factual-rationale
  accuracy relative to D.
\item
  \textbf{H2:} FC will increase objective contrastive/decision-flip
  accuracy relative to both F and D.
\item
  \textbf{H3:} Explanation condition will interact with decision
  quality: F and FC will improve discrimination between sound and
  anomalous decisions, measured by correctness/confidence calibration
  and appropriate endorsement, rather than simply increasing mean trust.
\item
  \textbf{H4:} For sound decisions, F and FC will improve procedural and
  informational fairness ratings relative to D; for anomalous decisions,
  explanations will not erase the quality penalty.
\item
  \textbf{H5:} FC may impose additional inspection time compared with F;
  its benefit will therefore be evaluated as an accuracy--time
  trade-off, not with the unsupported directional claim that more detail
  must reduce response time.
\end{itemize}

SUS and overall acceptance are secondary/exploratory because a replay
prototype is not a full operational product. Raw trust elevation is not
a success criterion: automation should be distrusted when it is
wrong.\hyperref[source-23]{[23]}

\section{10. Complete Methodology}\label{complete-methodology}

\subsection{Phase 1 --- Requirements and evidence
specification}\label{phase-1-requirements-and-evidence-specification}

Create a preregisterable requirements table before coding. Each row
contains a requirement, source, operationalization, test, and failure
response. The minimum requirements are:

\begin{itemize}
\tightlist
\item
  conflict-free phase transitions, jurisdiction-appropriate
  yellow/all-red and pedestrian clearance;
\item
  maximum pedestrian and vehicle service debt with no starvation in all
  feasible non-emergency states;
\item
  declared efficiency-loss tolerance relative to the best safe predicted
  plan;
\item
  deterministic decisions given the same state/configuration/seed;
\item
  safe fixed-time fallback for stale or invalid sensors;
\item
  emergency preemption and explicit post-preemption recovery;
\item
  explanations whose numbers, comparisons, constraints, and
  counterfactuals are trace-verifiable;
\item
  separation of traffic performance from human-perception outcomes;
\item
  reproducible stochastic scenarios, paired seeds, versioned
  configurations, and immutable study stimuli.
\end{itemize}

Freeze the control and HCI protocols before confirmatory data
collection. Pilot data may choose parameters and power assumptions but
must not enter the confirmatory analyses.

\subsection{Phase 2 --- Intersection
model}\label{phase-2-intersection-model}

Use one isolated four-arm intersection with four inbound through-only
lanes, four outbound lanes, four marked crosswalks, pedestrian waiting
areas, and no turns in the primary experiment. This yields ecological
recognizability without the confounding and safety complexity of
permissive turns. The action plan contains three service phases:

\begin{itemize}
\tightlist
\item
  \texttt{V\_NS}: north--south through vehicles;
\item
  \texttt{V\_EW}: east--west through vehicles;
\item
  \texttt{P\_ALL}: an exclusive all-pedestrian interval serving the four
  crosswalks.
\end{itemize}

Each service phase is followed by its configured transition states;
pedestrians receive WALK, pedestrian change, and buffer/clearance states
as required. SUMO owns the legal signal-state sequence, while the
external controller requests only safe target phases. Crosswalk length,
lane width, speed limit, minimum green, WALK, pedestrian clearance,
yellow, and all-red are configuration inputs derived from a documented
site or governing standard. The U.S. MUTCD 11th edition, for example,
specifies pedestrian timing guidance based on crossing distance and
walking speed and asks for slower speeds where slower pedestrians or
wheelchair users routinely cross; it is an example source, not a
universal substitute for local rules.\hyperref[source-24]{[24]}

Why this model: a two-arm single crossing is easy but weak for ITS
publication; a full turning-movement junction is realistic but makes the
first study mostly a conflict-engineering project. The three-phase
four-arm model gives genuine repeated resource competition and
interpretable explanations. Protected turns and concurrent pedestrian
phases are external-validity extensions after the core study.

\subsection{Phase 3 --- State, sensors, and
calibration}\label{phase-3-state-sensors-and-calibration}

SUMO provides ground-truth positions during development. The controller
sees a sensor abstraction, not privileged simulator internals: stop-line
loops/camera zones for vehicles, waiting-area/crosswalk zones for
pedestrians, SPaT, emergency flag, timestamp, and quality. A
configurable noise layer adds missed detections, false positives, count
error, and latency.

Collect state every \(\Delta t = 1\,\mathrm{s}\) as a software design choice;
recalculate candidate scores at the same cadence but permit a phase
change only after minimum-green and clearance constraints. Recheck this
cadence in a latency sensitivity study. Compute arrival rates over a
preregistered rolling window (candidate windows 30, 60, and 120 s in
pilot), not from a single tick. Normalization scales come from
training/calibration runs or service standards and are frozen before
testing.

\subsection{Phase 4 --- Agent
abstraction}\label{phase-4-agent-abstraction}

Implement two stateless scoring modules:

\begin{itemize}
\tightlist
\item
  \texttt{PedestrianAdvocate} receives pedestrian counts, arrival rate,
  mean/max waiting, waiting-area density, predicted service, sensor
  quality, and any accessibility extension request.
\item
  \texttt{VehicleAdvocate} receives queue length by approach, arrival
  rate, mean/max waiting, occupancy estimate, predicted discharge,
  emergency state, and sensor quality.
\end{itemize}

They do not learn and do not represent individual psychology. They
output a mode utility for each candidate plan, a service-debt weight,
and reason facts. Vehicle direction is selected by total-delay max
pressure inside plans that allocate vehicle service.

\subsection{Phase 5 --- Safe candidate-plan
generation}\label{phase-5-safe-candidate-plan-generation}

At every admissible epoch, enumerate a small finite set of horizon
plans, for example \(P_ALL\rightarrow{}V\), \(V_NS\rightarrow{}V_EW\), \(V_EW\rightarrow{}V_NS\),
continue-current, and plans with one legal phase extension. Generate
durations from a configured grid bounded by minimum/maximum green and
pedestrian clearance. Reject any plan with conflicting greens, truncated
clearance, incompatible emergency movement, impossible transition, or
predicted maximum-wait breach when an alternative feasible plan avoids
it.

The primary horizon and duration grid are selected in pilot profiling. A
practical initial grid is a 60--120 s horizon and 5 s duration
increments, clearly labeled as engineering candidates rather than
literature constants. If enumeration exceeds the real-time budget, prune
dominated plans or formulate the same finite choice as a small
mixed-integer program; do not add deep RL.

\subsection{Phase 6 --- Efficiency filter and
bargaining}\label{phase-6-efficiency-filter-and-bargaining}

Predict each plan's total person delay and mode service fraction with a
deterministic queue-discharge model calibrated from SUMO. Find the
minimum-loss plan, then retain only plans no more than \(\delta\) worse.
Within this set, enforce individual rationality relative to a safe
fallback and maximize the weighted Nash product. Hard starvation
constraints and emergency policy precede both operations. Section 11
supplies the equations.

Choose \(\delta\) using calibration-only Pareto curves: identify the knee or
the smallest tolerance that materially reduces maximum wait/starvation,
then preregister it. Report the full sensitivity set (recommended grid
\texttt{0,\ .02,\ .05,\ .10,\ .20}) so no conclusion rests on one
convenient threshold.

\subsection{Phase 7 --- Signal execution and
fallback}\label{phase-7-signal-execution-and-fallback}

The \texttt{SignalSafetyKernel} receives a target service phase and
returns the legal transition sequence. It refuses unsafe commands.
During an emergency, a verified route/movement request invokes
preemption: finish any non-truncatable pedestrian clearance, transition
through yellow/all-red, serve the emergency-compatible vehicle phase,
and then enter a recovery policy that repays accumulated
pedestrian/vehicle debt. If sensor data exceed age/quality limits,
switch to the documented fixed-time plan and explain the degraded mode.

Emergency detection itself is out of scope: simulation supplies a
signed/ground-truth flag. This avoids pretending that the prototype
validates V2X security or emergency-vehicle recognition.

\subsection{Phase 8 --- Decision trace and explanation
generation}\label{phase-8-decision-trace-and-explanation-generation}

Persist the exact input state, configuration hash, feasible/rejected
candidates, scores, chosen plan, binding constraints, emergency/fallback
status, top reason factors, alternative, and counterfactual search.
Generate text from typed templates only. Validate every explanation
before display; a validation failure shows a neutral message such as
``Detailed rationale unavailable; decision ID \ldots'' rather than
plausible invented prose.

\subsection{Phase 9 --- Simulation
experiment}\label{phase-9-simulation-experiment}

For each scenario/controller/seed tuple:

\begin{enumerate}
\def\labelenumi{\arabic{enumi}.}
\tightlist
\item
  generate identical demand streams once;
\item
  run a 600 s warm-up not used in primary metrics;
\item
  run a 3,600 s evaluation window;
\item
  log 1 Hz states, all decisions, trip/person summaries, constraint
  events, and latency;
\item
  validate the run and aggregate metrics.
\end{enumerate}

The durations are initial protocol choices, not universal constants.
Pilot autocorrelation and rare-event rates determine whether longer runs
are required. Begin with at least 30 paired seeds per controller and add
batches of 10 until the primary-metric 95\% confidence-interval
half-width is below 5\% of its mean or a preregistered cap of 100 is
reached. Use common random numbers: a given seed produces the same
arrivals and walking speeds for every controller.

\subsection{Phase 10 --- Dashboard and stimulus
production}\label{phase-10-dashboard-and-stimulus-production}

Build a React/TypeScript dashboard with an SVG intersection,
phase/status strip, mode-demand panels, decision card, rationale card,
alternative/counterfactual disclosure, data-freshness indicator, and
decision-history timeline. The prominent signal indications remain
authoritative; explanations are labeled ``why the controller chose this
plan.''

Export immutable JSON replay traces (and optional MP4 previews) from
selected validated runs. Human-study conditions render those same traces
locally, preventing network or live-SUMO randomness from confounding
explanation effects.

\subsection{Phase 11 --- HCI pilot and confirmatory
study}\label{phase-11-hci-pilot-and-confirmatory-study}

Conduct cognitive interviews with 6--10 adults to find ambiguous wording
and misunderstood metrics, plus optional formative review by 3--6
traffic practitioners. Then run a 12--18 participant instrument/timing
pilot. Freeze copy, scoring keys, exclusion rules, power simulation, and
analysis before recruiting the confirmatory sample. Section 19 defines
the full design.

\subsection{Phase 12 --- Analysis and
reporting}\label{phase-12-analysis-and-reporting}

Report controller outcomes as paired differences with distributions and
confidence intervals; do not report only grand means. Fit mixed models
for trial-level HCI outcomes, test the preregistered contrasts, apply
multiplicity correction, and publish all null and adverse findings. Join
the technical and HCI results by \texttt{decision\_id} only for
explanation-fidelity and scenario-stratified analyses, not to imply that
perceived fairness validates objective traffic fairness.

\section{11. Mathematical Model}\label{mathematical-model}

\subsection{11.1 Sets, time, and phase state}\label{sets-time-and-phase-state}

Time is sampled at
\[
  t \in \{0,\Delta t,2\Delta t,\ldots\}, \qquad \Delta t = 1\,\mathrm{s}.
\]
Let the modes be \(\mathcal{M}=\{p,v\}\), the vehicle approaches be
\(\mathcal{R}=\{N,S,E,W\}\), the pedestrian crosswalks be
\(\mathcal{C}=\{N,S,E,W\}\), and the service phases be
\(\Phi=\{P_{\mathrm{ALL}},V_{\mathrm{NS}},V_{\mathrm{EW}}\}\). A
candidate plan \(a\in\mathcal{A}_t\) is an ordered sequence of service
phases and durations over horizon \(H\), including all required
transitions. The set \(\mathcal{A}_{\mathrm{safe}}(s_t)\) contains only
plans that are legal in state \(s_t\).

\subsection{11.2 Pedestrian state}\label{pedestrian-state}

For pedestrian \(i\) waiting at crosswalk \(c\), let \(\tau_i^p\) be
the time at which the pedestrian entered the waiting zone. Define
\begin{align}
  w_i^p(t) &= \max\!\left(0,t-\tau_i^p\right) &&[\mathrm{s}],\\
  n_c^p(t) &= \left|Q_c^p(t)\right| &&[\mathrm{persons}],\\
  \bar w_c^p(t) &= \frac{1}{\max(1,n_c^p(t))}
    \sum_{i\in Q_c^p(t)} w_i^p(t) &&[\mathrm{s}],\\
  w_c^{p,\max}(t) &= \max_{i\in Q_c^p(t)} w_i^p(t) &&[\mathrm{s}],\\
  \lambda_c^p(t) &= \frac{A_c^p(t-W,t)}{W} &&[\mathrm{persons/s}],\\
  \rho_c^p(t) &= \frac{n_c^p(t)}{\operatorname{area}_c}
    &&[\mathrm{persons/m^2}].
\end{align}
The maximum is defined as zero when a queue is empty. Here,
\(A_c^p(t-W,t)\) is the number of arrivals in rolling window \(W\).
Aggregate pedestrian demand over the horizon is
\[
  D_p(t)=\sum_{c\in\mathcal{C}}
  \left[n_c^p(t)+H\lambda_c^p(t)\right].
\]
Mean and maximum wait are reported separately. Density is a
safety/crowding descriptor and must not double-count demand unless a
calibrated model explicitly uses it.

\subsection{11.3 Vehicle state}\label{vehicle-state}

A vehicle is queued when its speed falls below declared threshold
\(v_{\mathrm{stop}}\) inside detection zone \(r\). A candidate value of
\(0.1\,\mathrm{m/s}\) may be tested during the pilot; it is not treated
as a universal constant. Define
\begin{align}
  q_r^v(t) &= \sum_{j\in Z_r(t)}
    \mathbf{1}\!\left[v_j(t)<v_{\mathrm{stop}}\right]
    &&[\mathrm{vehicles}],\\
  w_j^v(t) &= \int_{\tau_j^v}^{t}
    \mathbf{1}\!\left[v_j(u)<v_{\mathrm{stop}}\right],du
    &&[\mathrm{s}],\\
  \bar w_r^v(t) &= \frac{1}{\max(1,q_r^v(t))}
    \sum_{j\in Q_r^v(t)} w_j^v(t) &&[\mathrm{s}],\\
  w_r^{v,\max}(t) &= \max_{j\in Q_r^v(t)}w_j^v(t)
    &&[\mathrm{s}],\\
  \lambda_r^v(t) &= \frac{A_r^v(t-W,t)}{W}
    &&[\mathrm{vehicles/s}],\\
  D_v(t) &= \sum_{r\in\mathcal{R}}
    \left[q_r^v(t)+H\lambda_r^v(t)\right].
\end{align}
If occupancy is observed, predicted vehicle delay is multiplied by the
number of occupants. Otherwise, vehicle delay is reported separately
and person-equivalent sensitivity analyses use preregistered occupancy
assumptions; one vehicle is never silently equated with one person.

\subsection{11.4 Data quality and normalization}\label{data-quality-and-normalization}

For sensor feature \(x\), define data age and quality as
\[
  \operatorname{age}_x(t)=t-\operatorname{timestamp}_x,
  \qquad q_x\in[0,1].
\]
The controller rejects or degrades the feature when
\(\operatorname{age}_x(t)>A_x^{\max}\) or \(q_x<Q_x^{\min}\), with
limits established by sensor requirements and failure tests. A
nonnegative feature is normalized using frozen robust scaling:
\[
  \mathcal{N}(x;s_x)=\operatorname{clip}\!\left(\frac{x}{s_x},0,1\right).
\]
The scale \(s_x\) is either an operational limit or a high calibration
quantile. The pilot compares the 90th, 95th, and 99th percentiles. Test
runs are never normalized by their own maxima.

\subsection{11.5 Predicted service and utility}\label{predicted-service-and-utility}

For each plan, a calibrated deterministic queue model predicts the
served demand \(S_m(a,t)\). The bounded service utility is
\[
  u_m(a,t)=\min\!\left(1,
  \frac{S_m(a,t)}{\max(1,D_m(t))}\right), \qquad m\in\mathcal{M}.
\]
This is an opportunity-to-serve fraction, not a model of human
happiness. Let \(a_t^0\) be the safe fallback plan and let
\(d_m(t)=u_m(a_t^0,t)\) be its disagreement utility. The individually
rational set is
\[
  \mathcal{A}_{\mathrm{IR}}(t)=
  \left\{a\in\mathcal{A}_{\mathrm{safe}}(s_t):
  u_m(a,t)\ge d_m(t)\;\forall m\in\mathcal{M}\right\}.
\]
If this set is empty, the trace records the failure and the controller
uses the lexicographic maximum-wait rule instead of concealing the
conflict with a numerical constant.

\subsection{11.6 Waiting debt and bargaining weights}\label{waiting-debt-and-bargaining-weights}

Let the current maximum waits by mode be
\[
  W_p(t)=\max_{c\in\mathcal{C}}w_c^{p,\max}(t), \qquad
  W_v(t)=\max_{r\in\mathcal{R}}w_r^{v,\max}(t).
\]
Given policy service limits \(T_m^{\max}\), define
\begin{align}
  r_m(t) &= \operatorname{clip}\!\left(
    \frac{W_m(t)}{T_m^{\max}},0,1\right),\\
  \widetilde{\omega}_m(t) &= \pi_m
    \left[1+\kappa r_m(t)^{\eta}\right],\\
  \omega_m(t) &= \frac{\widetilde{\omega}_m(t)}
    {\sum_{k\in\mathcal{M}}\widetilde{\omega}_k(t)}.
\end{align}
Here \(\pi_m>0\) is a declared base policy weight. Equal base weights
are the primary starting point unless standards or stakeholder evidence
supports another choice. Parameters \(\kappa\ge0\) and \(\eta\ge1\)
are selected on calibration scenarios and sensitivity-tested.

\subsection{11.7 Safety, starvation, and emergency feasibility}\label{safety-starvation-and-emergency-feasibility}

Let \(C_{\mathrm{safe}}(a,s_t)=1\) only when all conflict matrices,
minimum service intervals, clearances, and transition automata are
satisfied. With predicted maximum wait \(\widehat W_m(a,t+H)\), define
\[
  \mathcal{A}_{\mathrm{starve}}(t)=
  \left\{a\in\mathcal{A}_{\mathrm{IR}}(t):
  \widehat W_m(a,t+H)\le T_m^{\max}\;\forall m\in\mathcal{M}\right\}.
\]
If the set is empty, select the plan that lexicographically minimizes
the number of service-limit breaches, the largest normalized breach,
and then efficiency loss. Safety constraints are never relaxed. When
emergency state \(e(t)=1\), the controller restricts the feasible set
to emergency-compatible plans after completing mandatory clearance;
ordinary bargaining is suspended and displaced service debt is retained
for recovery.

\subsection{11.8 Near-efficient candidate set}\label{near-efficient-candidate-set}

The predicted efficiency loss over the horizon is
\[
  L_{\mathrm{eff}}(a,t)=
  \sum_{m\in\mathcal{M}}c_m\widehat D_m(a,t)
  +c_{\mathrm{sw}}N_{\mathrm{switch}}(a),
\]
where \(\widehat D_m\) is predicted person-delay when occupancy is
known, \(N_{\mathrm{switch}}\) is the number of service-phase changes,
and the \(c\)-terms are declared unit conversions or policy weights.
Let
\begin{align}
  L_{\min}(t) &= \min_{a\in\mathcal{A}_{\mathrm{starve}}(t)}
    L_{\mathrm{eff}}(a,t),\\
  \mathcal{A}_{\delta}(t) &=
    \left\{a\in\mathcal{A}_{\mathrm{starve}}(t):
    L_{\mathrm{eff}}(a,t)\le L_{\min}(t)
    +\delta\max\!\left(L_{\min}(t),L_{\mathrm{floor}}\right)\right\}.
\end{align}
The positive floor avoids a meaningless relative bound when predicted
loss is nearly zero. Both relative and absolute delay differences are
reported, and \(\delta\) is selected only on calibration data.

\subsection{11.9 Weighted Nash bargaining decision}\label{weighted-nash-bargaining-decision}

For bargaining gain \(g_m(a,t)=u_m(a,t)-d_m(t)\), select
\[
  a^*(t)=\operatorname*{arg\,max}_{\substack{
      a\in\mathcal{A}_{\delta}(t)\\ g_m(a,t)\ge0}}
    \sum_{m\in\mathcal{M}}\omega_m(t)
    \log\!\left(g_m(a,t)+\varepsilon\right).
\]
The small numerical constant \(\varepsilon\) is declared in the
configuration and varied to verify rank stability. Ties are resolved by
lower predicted maximum normalized wait, fewer switches, and finally a
stable candidate identifier. The finite search records every feasible
candidate and score. The solution is proportionally fair over the
defined utility gains; that property alone does not establish moral
fairness.

\subsection{11.10 Direction-level vehicle choice}\label{direction-level-vehicle-choice}

For vehicle movement group \(g\in\{\mathrm{NS},\mathrm{EW}\}\), define
total-delay pressure as
\[
  P_g(t)=\sum_{r\in g}
  \left[Q^{\mathrm{up}}_{\mathrm{delay},r}(t)
  -Q^{\mathrm{down}}_{\mathrm{delay},r}(t)\right].
\]
When a candidate allocates vehicle service, the legal direction with
the larger pressure is chosen subject to current-phase extension and
transition cost. If the isolated network omits explicit spillback,
downstream terms use measured receiving-lane occupancy and this
approximation is declared.\hyperref[source-10]{[10]}

\subsection{11.11 Counterfactual optimization}\label{counterfactual-optimization}

Given chosen action \(a^*\) and alternative \(b\), find the smallest
permitted state perturbation that makes \(b\) win:
\begin{align}
  \Delta s^* ={}& \operatorname*{arg\,min}_{\Delta s\in\Omega}
    \left\|D^{-1}\Delta s\right\|_1,\\
  \text{subject to }{}&
    \operatorname{controller}(s_t+\Delta s)=b,\\
  & C_{\mathrm{state}}(s_t+\Delta s)=1,
    \quad\text{with immutable facts unchanged.}
\end{align}
The feasible perturbation set \(\Omega\) changes only meaningful,
controllable variables on their measurement grids. Matrix \(D\)
contains human-readable step sizes such as one person, one vehicle, one
second, or one arrival-rate bin. The actual controller is replayed at
each search point. If no valid flip exists within the declared range,
the interface says so rather than inventing a threshold.

\section{12. Game-Theoretic Model
Selection}\label{game-theoretic-model-selection}

{\def\LTcaptype{none} % do not increment counter
\begin{longtable}[]{@{}
  >{\raggedright\arraybackslash}p{(\linewidth - 6\tabcolsep) * \real{0.2500}}
  >{\raggedright\arraybackslash}p{(\linewidth - 6\tabcolsep) * \real{0.2500}}
  >{\raggedright\arraybackslash}p{(\linewidth - 6\tabcolsep) * \real{0.2500}}
  >{\raggedright\arraybackslash}p{(\linewidth - 6\tabcolsep) * \real{0.2500}}@{}}
\toprule\noalign{}
\begin{minipage}[b]{\linewidth}\raggedright
Candidate
\end{minipage} & \begin{minipage}[b]{\linewidth}\raggedright
Advantages
\end{minipage} & \begin{minipage}[b]{\linewidth}\raggedright
Disadvantages
\end{minipage} & \begin{minipage}[b]{\linewidth}\raggedright
Suitability
\end{minipage} \\
\midrule\noalign{}
\endhead
\bottomrule\noalign{}
\endlastfoot
Non-cooperative Nash & Familiar strategic framing & Empirically weak for
real pedestrian--driver behavior; multiple equilibria; poor
public-service semantics & Reject for controller \\
Behavioral/quantal-response game & Better descriptive fit and allows
bounded rationality & Requires human interaction data and parameter
estimation; unnecessary for signal allocation & Use only in a separate
behavior study \\
Stackelberg game & Represents leader--follower interaction & Casts
infrastructure as strategic leader and people as responders; ethically
and empirically awkward & Reject for core \\
Auction/priority market & Expresses heterogeneous urgency & Needs
bids/credits, identity, incentive compatibility, privacy, and legitimacy
safeguards & Reject for university prototype \\
Multi-agent RL & Can learn complex network coordination & Opaque,
data-intensive, reward-sensitive, and difficult to reproduce/explain &
Comparator from literature, not implementation \\
Max-min allocation & Strong protection for worst-served mode & Can
sacrifice substantial efficiency and switch excessively & Useful
ablation/secondary objective \\
Weighted-sum optimization & Simple and fast & Scale-sensitive
coefficients; can hide severe loss to one mode & Efficiency baseline \\
Weighted Nash bargaining & Transparent proportional-fairness property;
supports disagreement point and inspectable trade-offs &
Utilities/weights still normative; existing pedestrian--vehicle use
means limited novelty & \textbf{Recommended within safety and efficiency
bounds} \\
\end{longtable}
}

The recommended primary mechanism is weighted Nash bargaining over
virtual mode advocates. Its justification is institutional: scarce safe
crossing service must be allocated between stakeholder classes according
to a declared rule. It is not a claim that people negotiate with a
traffic light. The study must compare EFPB against (i) Wang-style
unbounded/equal-weight bargaining and (ii) a simpler weighted-score or
delay-pressure controller. If performance and explanations are
indistinguishable, game theory is unnecessary and the simpler controller
should be preferred.

The disagreement point uses a real safe fallback plan, not zero service,
so ``agreement'' means each mode does at least as well in predicted
service as that fallback. Individual-rationality failure is possible
under acute conflict and must be logged, not mathematically concealed.
Safety, emergency, and starvation rules remain lexicographically
superior to bargaining.

\section{13. Fairness Model}\label{fairness-model}

Fairness is specified as a layered contract:

\begin{enumerate}
\def\labelenumi{\arabic{enumi}.}
\tightlist
\item
  \textbf{Safety fairness:} no mode receives service through unsafe
  conflict; vulnerable users are not asked to trade safety for
  throughput.
\item
  \textbf{Minimum-service fairness:} feasible ordinary operation must
  respect mode maximum-wait limits; violations and causes are countable.
\item
  \textbf{Proportional allocative fairness:} among near-efficient,
  individually rational plans, weighted Nash bargaining balances
  marginal utility gains.
\item
  \textbf{Outcome distribution:} report mode-specific mean, median,
  90th/95th percentile, and maximum waits plus service ratios; do not
  hide tails in a mean.
\item
  \textbf{Procedural fairness:} the same public rule, constraints,
  override policy, and explanation contract apply consistently.
\item
  \textbf{Perceived fairness:} participant judgments are a distinct
  empirical construct and cannot certify the first five.
\end{enumerate}

For run \(r\), let
\[
  z_m = \frac{N_m^{\mathrm{served}}}{N_m^{\mathrm{requested}}},
  \qquad m\in\{p,v\},
\]
which is undefined if there is no demand for mode \(m\). A descriptive
two-mode Jain index is
\[
  J_r = \frac{(z_p+z_v)^2}{2\left(z_p^2+z_v^2\right)}.
\]

Also report normalized wait disparity
\[
  G_r =
  \left|
    \frac{\overline{W}_p}{T_p^{\mathrm{ref}}}
    -
    \frac{\overline{W}_v}{T_v^{\mathrm{ref}}}
  \right|
\]
and tail disparity
\[
  G_r^{95} =
  \left|
    \frac{Q_{0.95}(W_p)}{T_p^{\mathrm{ref}}}
    -
    \frac{Q_{0.95}(W_v)}{T_v^{\mathrm{ref}}}
  \right|.
\]

Reference times are policy or baseline targets fixed before testing.
Jain equality is secondary: \(J=1\) could describe two modes both
receiving disastrously little service. Every fairness plot must
therefore be paired with total delay, service levels, and maximum wait.
Likewise, equal raw seconds are not necessarily equitable because a
stopped vehicle may contain several people and pedestrian
exposure/crossing constraints differ.

The controller implements a constrained rather than unconstrained
trade-off:
\[
  \underset{a\in\mathcal{A}_{\delta}(t)}{\operatorname{minimize}}
  \quad \text{fairness loss}(a),
\]
with safety and starvation enforced first.

Report a Pareto frontier of total person delay versus maximum normalized
wait/disparity. The primary fairness claim is supported only if EFPB
reduces preregistered disparity/tail/starvation outcomes while observed
efficiency remains inside the declared bound across most relevant
scenarios, not only on average.

\section{14. Explainability Model}\label{explainability-model}

\subsection{14.1 Audience and question}\label{audience-and-question}

The primary audience is an adult road user or civic stakeholder
inspecting a live/replayed decision; the secondary audience is a traffic
operator auditing it. The dashboard is not intended for someone to read
while driving and never supersedes the traffic signal. It answers:

\begin{itemize}
\tightlist
\item
  \textbf{What happened?} Which service phase was selected and for how
  long?
\item
  \textbf{Why this?} Which two or three trace facts made the largest
  decision-relevant difference?
\item
  \textbf{Why not the alternative?} Which score or constraint prevented
  the strongest rejected alternative?
\item
  \textbf{What would change it?} What smallest feasible controllable
  state change makes that alternative win?
\item
  \textbf{Was an override active?} Safety, emergency, fallback, and
  data-quality status.
\end{itemize}

\subsection{14.2 Explanation record}\label{explanation-record}

Every \texttt{ExplanationRecord} contains:

\texttt{decision\_id}, \texttt{generated\_at}, \texttt{audience},
\texttt{decision\_summary}, \texttt{reason\_facts{[}{]}},
\texttt{binding\_constraints{[}{]}}, \texttt{alternative},
\texttt{counterfactual}, \texttt{override}, \texttt{data\_freshness},
\texttt{uncertainty\_note}, \texttt{template\_id}, \texttt{trace\_hash},
and \texttt{validation\_status}.

Each reason fact includes \texttt{feature\_id}, display label, raw
value/unit, normalized value, comparison reference, direction of
influence, and source timestamp. Select reasons by exact
leave-one-factor-out score change within the transparent controller,
constrained to facts that are actionable or normatively relevant. Cap
the default display at three reasons; details remain expandable.

\subsection{14.3 Recommended formats}\label{recommended-formats}

\textbf{Factual template}

\begin{quote}
Pedestrian service was selected for the next 25 s. The longest
pedestrian wait is 78 s (78\% of the 100 s service limit), compared with
a 34 s longest vehicle wait, and this plan remains within the 5\%
predicted-delay allowance. No emergency override is active.
\end{quote}

\textbf{Alternative template}

\begin{quote}
North--south vehicle service was the closest alternative. It was not
selected because it would raise the predicted pedestrian wait above the
configured service limit.
\end{quote}

\textbf{Counterfactual template}

\begin{quote}
With all other observed values unchanged, north--south vehicles would be
selected if the longest pedestrian wait were 17 s lower. This result was
replayed through controller version \(\ldots{}\).
\end{quote}

These are format examples, not results or preset thresholds. Actual
numbers come from a stored trace. In emergency mode the first sentence
must say that the emergency policy---not ordinary fairness
bargaining---determined the priority and that mandatory pedestrian
clearance was preserved.

\subsection{14.4 Why templates, not SHAP or an
LLM}\label{why-templates-not-shap-or-an-llm}

The decision engine is finite, symbolic, and logs candidate utilities. A
deterministic template is faithful, fast, multilingual-ready, and
testable. SHAP is valuable for opaque learned policies but can be
mistaken for a causal explanation and introduces background-distribution
choices. A language model adds nondeterminism, hallucination risk,
latency, privacy exposure, and no necessary capability. Natural-language
realization can be added later only as a paraphrase whose semantic slots
are automatically checked against the canonical trace.

\subsection{14.5 Fidelity and quality
tests}\label{fidelity-and-quality-tests}

\begin{itemize}
\tightlist
\item
  \textbf{Value fidelity:} every displayed value equals its trace value
  after declared rounding.
\item
  \textbf{Decision fidelity:} replaying
  \texttt{state\ +\ config\ +\ version} returns the same candidate and
  reason ordering.
\item
  \textbf{Constraint fidelity:} every described constraint has a
  matching evaluated predicate; ``binding'' requires slack within a
  configured numerical tolerance.
\item
  \textbf{Counterfactual validity:} replaying \(s+\Delta s^*\) selects the
  stated alternative; immutable/safety constraints remain true.
\item
  \textbf{Minimality audit:} no tested lower-distance perturbation flips
  the decision under the enumerated search grid.
\item
  \textbf{Coverage:} percentage of ordinary and override decisions for
  which a validated explanation is produced.
\item
  \textbf{Stability:} small measurement perturbations that do not change
  the action should not radically replace all reasons; quantify
  reason-set Jaccard similarity.
\item
  \textbf{Latency:} mean, p95, and p99 generation and counterfactual
  search time.
\item
  \textbf{Human quality:} objective rationale/contrastive accuracy,
  response time, perceived informativeness, and optional Explanation
  Satisfaction Scale; do not substitute self-reported clarity for
  comprehension.\hyperref[source-22]{[22]}
\end{itemize}

\section{15. System Architecture}\label{system-architecture}

\begin{verbatim}
SUMO ground truth
      │
      ▼
Sensor adapter ──► noise/staleness model ──► IntersectionState
                                                  │
                         ┌────────────────────────┼───────────────────────┐
                         ▼                        ▼                       ▼
               PedestrianAdvocate        VehicleAdvocate       Safety/Emergency kernel
                         │                        │                       │
                         └──────────────┬─────────┘                       │
                                        ▼                                │
                              Candidate plan generator ◄─────────────────┘
                                        │
                                        ▼
                      prediction → efficiency filter → bargaining
                                        │
                         ┌──────────────┴──────────────┐
                         ▼                             ▼
                  target signal plan             DecisionTrace
                         │                             │
                         ▼                             ├──► explanation validator
                  TraCI executor                       │          │
                                                       ▼          ▼
                                                 JSONL/Parquet  ExplanationRecord
                                                                  │
                                    ┌─────────────────────────────┼────────────┐
                                    ▼                             ▼            ▼
                              WebSocket API                 live dashboard   study replay
\end{verbatim}

Safety and explanation are not peer opinions in a voting architecture.
The safety kernel constrains what can be chosen; the controller makes
the choice; the explanation reports that choice. The frontend cannot
issue raw signal colors. It can start a configured simulation, select a
scenario, or request replay, while only the backend safety kernel can
translate target service phases into legal SUMO states.

Every run has a \texttt{run\_id}; every control epoch has a
\texttt{decision\_id}; every artifact stores \texttt{scenario\_hash},
\texttt{config\_hash}, \texttt{software\_commit}, and \texttt{seed}. The
trace is append-only during a run. UI text and study responses reference
a decision ID, preventing accidental joining by timestamp alone.

\section{16. Simulation Architecture}\label{simulation-architecture}

\subsection{16.1 Tool selection}\label{tool-selection}

{\def\LTcaptype{none} % do not increment counter
\begin{longtable}[]{@{}
  >{\raggedright\arraybackslash}p{(\linewidth - 12\tabcolsep) * \real{0.1429}}
  >{\raggedright\arraybackslash}p{(\linewidth - 12\tabcolsep) * \real{0.1429}}
  >{\raggedright\arraybackslash}p{(\linewidth - 12\tabcolsep) * \real{0.1429}}
  >{\raggedright\arraybackslash}p{(\linewidth - 12\tabcolsep) * \real{0.1429}}
  >{\raggedright\arraybackslash}p{(\linewidth - 12\tabcolsep) * \real{0.1429}}
  >{\raggedright\arraybackslash}p{(\linewidth - 12\tabcolsep) * \real{0.1429}}
  >{\raggedright\arraybackslash}p{(\linewidth - 12\tabcolsep) * \real{0.1429}}@{}}
\toprule\noalign{}
\begin{minipage}[b]{\linewidth}\raggedright
Tool
\end{minipage} & \begin{minipage}[b]{\linewidth}\raggedright
Pedestrians
\end{minipage} & \begin{minipage}[b]{\linewidth}\raggedright
Signal control/Python
\end{minipage} & \begin{minipage}[b]{\linewidth}\raggedright
Scale
\end{minipage} & \begin{minipage}[b]{\linewidth}\raggedright
Visual realism
\end{minipage} & \begin{minipage}[b]{\linewidth}\raggedright
Reproducibility
\end{minipage} & \begin{minipage}[b]{\linewidth}\raggedright
Verdict
\end{minipage} \\
\midrule\noalign{}
\endhead
\bottomrule\noalign{}
\endlastfoot
SUMO & Native persons, walking areas, crossings & Mature TraCI/libsumo
and traffic-light domains & Good for corridor/city; excellent for one
junction & Strong 2D GUI & Strong, open source, config-based &
\textbf{Select} \\
CityFlow & Primarily vehicles & Fast Python-friendly large-network
control & Excellent & Basic & Strong & Reject: pedestrian model not the
core strength \\
MATSim & Agent-based daily travel demand & Control possible but not
microscopic signal/pedestrian interaction focus & Metropolitan & Limited
micro-visualization & Strong & Reject for intersection controller \\
Custom Python & Only what is built & Full control & Adequate only for
toy model & Must be built & Depends on implementation & Use only for
fast horizon prediction, not evidence simulator \\
\end{longtable}
}

SUMO is the best choice because it directly supports vehicles, persons,
pedestrian crossings, traffic lights, XML scenario configuration, Python
control, trip/person outputs, and visual inspection. The official TraCI
API exposes vehicle, person, edge, induction-loop, lane-area, and
traffic-light domains; SUMO documents custom signal control through
TraCI and a pedestrian-crossing control
tutorial.\hyperref[source-25]{[25]}\,\hyperref[source-26]{[26]}\,\hyperref[source-27]{[27]}

\subsection{16.2 Components}\label{components}

\begin{itemize}
\tightlist
\item
  \texttt{intersection.net.xml}: four arms, walking areas, four
  crossings, signal links, and conflict-tested connections.
\item
  \texttt{base\_tls.add.xml}: long-placeholder service phases plus fixed
  yellow/all-red/pedestrian transitions. The controller never skips
  transitions.
\item
  \texttt{scenario.rou.xml}: vehicles and persons generated before each
  run from a manifest.
\item
  \texttt{detectors.add.xml}: lane-area/induction detectors and
  pedestrian zones or edge-based sensing abstraction.
\item
  \texttt{simulation.sumocfg}: network, routes, additional files, step
  length, output files, and deterministic seed.
\item
  Python runner: starts \texttt{sumo} for batches or \texttt{sumo-gui}
  for inspection, subscribes to state variables, calls the controller,
  applies safe phase targets, and logs.
\item
  Batch mode: begin with TraCI for debuggability; switch to
  \texttt{libsumo} behind the same API only after regression tests,
  because it avoids socket overhead for many
  replications.\hyperref[source-25]{[25]}
\end{itemize}

\subsection{16.3 Network and signal
verification}\label{network-and-signal-verification}

Build the network in NetEdit, then version the generated XML. Use
\texttt{netconvert\ -\/-check-lane-foes.all\ true} and SUMO validation
to inspect conflict relationships. In the GUI, manually step through
every service-to-service transition, including emergency preemption and
recovery. Automated tests assert that incompatible signal links are
never green simultaneously and that every pedestrian already in a
crosswalk receives the required clearance.

Pedestrian timing is a geometry-dependent constraint. In an
MUTCD-governed example, clearance uses crossing distance and
walking-speed guidance, total WALK plus clearance must accommodate a
lower design speed, and slower populations require consideration of a
still lower speed.\hyperref[source-24]{[24]} For another jurisdiction,
replace these values with its standard and retain the
formula/configuration provenance.

\subsection{16.4 TraCI loop}\label{traci-loop}

At simulation second \texttt{t}:

\begin{enumerate}
\def\labelenumi{\arabic{enumi}.}
\tightlist
\item
  read subscriptions and current SPaT;
\item
  update entry/exit registries and waiting histories;
\item
  apply the configured sensor-error transform;
\item
  construct and validate \texttt{IntersectionState};
\item
  if a decision is admissible, enumerate and score plans;
\item
  pass the target to \texttt{SignalSafetyKernel};
\item
  apply legal state/duration commands through the traffic-light domain;
\item
  append state, decision, explanation, timing, and event records;
\item
  publish a compact \texttt{LiveFrame} to the WebSocket;
\item
  call \texttt{simulationStep()}.
\end{enumerate}

Use subscriptions rather than repeated per-object calls for scale. The
official pedestrian tutorial illustrates person waiting and route
inspection via TraCI; adapt it into a testable adapter instead of
scattering raw calls through controller code.\hyperref[source-27]{[27]}

\subsection{16.5 Output validation}\label{output-validation}

A run is invalid if it terminates early, produces collisions/teleports
attributable to controller/network errors, has missing decision records,
violates signal invariants, or has unmatched demand manifests. Teleports
caused by extreme congestion are reported separately and trigger
scenario/network review, not silent deletion. Compare online aggregates
against SUMO tripinfo/personinfo recomputation to detect logger
mistakes.

\section{17. Scenario Design}\label{scenario-design}

Define demand levels relative to calibrated intersection capacity so the
design transfers: \texttt{L}, \texttt{M}, \texttt{H}, and \texttt{X} are
approximately 25\%, 50\%, 75\%, and 90--100\% of empirically estimated
feasible arrival capacity, finalized in pilot. Wait debt is created by
an initial demand history, never by editing a displayed wait after the
fact. ``Expected decision'' below means an invariant or likely priority
under the initialized state; the controller remains allowed to choose
the other mode when full state and constraints justify it.

\begin{landscape}
{\def\LTcaptype{none} % do not increment counter
\footnotesize
\setlength{\tabcolsep}{4pt}
\begin{longtable}[]{@{}
  >{\raggedright\arraybackslash}p{(\linewidth - 10\tabcolsep) * \real{0.1667}}
  >{\raggedright\arraybackslash}p{(\linewidth - 10\tabcolsep) * \real{0.1667}}
  >{\raggedright\arraybackslash}p{(\linewidth - 10\tabcolsep) * \real{0.1667}}
  >{\raggedright\arraybackslash}p{(\linewidth - 10\tabcolsep) * \real{0.1667}}
  >{\raggedright\arraybackslash}p{(\linewidth - 10\tabcolsep) * \real{0.1667}}
  >{\raggedright\arraybackslash}p{(\linewidth - 10\tabcolsep) * \real{0.1667}}@{}}
\toprule\noalign{}
\begin{minipage}[b]{\linewidth}\raggedright
ID
\end{minipage} & \begin{minipage}[b]{\linewidth}\raggedright
Scenario and input conditions
\end{minipage} & \begin{minipage}[b]{\linewidth}\raggedright
Expected decision/invariant
\end{minipage} & \begin{minipage}[b]{\linewidth}\raggedright
Expected fairness behavior
\end{minipage} & \begin{minipage}[b]{\linewidth}\raggedright
Example explanation content
\end{minipage} & \begin{minipage}[b]{\linewidth}\raggedright
Evaluation purpose
\end{minipage} \\
\midrule\noalign{}
\endhead
\bottomrule\noalign{}
\endlastfoot
S1 & Balanced: pedestrian M, vehicle M, symmetric approaches & Alternate
by marginal service gain/current debt & Similar normalized service and
bounded tails & Chosen mode has larger current debt or avoids a switch &
Nominal efficiency/fairness \\
S2 & Heavy pedestrian: pedestrian H, vehicle M & P\_ALL becomes more
frequent; not continuously green & Vehicle limit still enforced & High
pedestrian demand and wait; vehicle bound remains feasible & Pedestrian
responsiveness \\
S3 & Heavy vehicle: pedestrian M, vehicle H & Vehicle phases more
frequent; pedestrian cap remains hard & No pedestrian starvation & Large
vehicle service deficit; pedestrian wait below bound & Vehicle
responsiveness \\
S4 & High pedestrian wait: low current count but initialized max wait
near limit & P\_ALL unless safety/emergency prevents it & Tail
protection dominates count & One person's wait is close to service limit
& Test max-wait vs majority demand \\
S5 & High vehicle wait on NS; moderate pedestrians & V\_NS likely
selected & Vehicle tail debt repaid & NS cumulative delay exceeds EW and
pedestrian debt is feasible & Direction and mode hierarchy \\
S6 & Conflicting high demand: pedestrian H, vehicle H & Alternate near
Pareto-efficient plans & Both service limits; expose unavoidable
exceptions & Both modes are busy; plan is within \(\delta\) and minimizes
normalized tail & Main trade-off stress test \\
S7 & Emergency vehicle request on NS during ordinary phase & Finish
mandatory clearance, then V\_NS & Override logged; recovery repays
displaced mode & Emergency policy active; clearance preserved; recovery
scheduled & Emergency safety/response \\
S8 & Repeated imbalance: pedestrian H for multiple cycles, then vehicle
surge & Debt prevents monopoly and adapts after surge & Weight
trajectory and no starvation & Prior service history raised the
underserviced mode's weight & Temporal/procedural fairness \\
S9 & Near-saturation edge: both X with finite downstream capacity & No
unsafe switch; exceptions may be unavoidable & Count and duration of
service-limit breaches & Demand exceeds configured service envelope;
least-breach plan chosen & Graceful failure, no overclaim \\
S10 & Low demand/reference: pedestrian L, vehicle L & Avoid unnecessary
switching; calls served promptly & High Jain score is not interpreted
alone & Current phase extended because all waits are low & Operational
sanity/control \\
S11 & Mobility-impaired: a configured subset uses slower walking
speed/extension request & Full legal pedestrian clearance; no truncation
& Accessibility request has procedural priority & Extended clearance is
required by active request & Accessibility and timing \\
S12 & Sensor degradation: missed pedestrian detections, stale vehicle
sensor, and recovery levels & Switch to conservative/fixed fallback at
declared threshold & Degraded-mode impact reported by mode & Rationale
unavailable/limited because sensor is stale; fallback active &
Robustness and honest uncertainty \\
\end{longtable}
}
\end{landscape}

Use a factorial core---pedestrian demand \texttt{\{L,M,H\}} × vehicle
demand \texttt{\{L,M,H\}} × controller---with targeted S4--S12 stress
tests. Include asymmetric NS/EW splits and two demand processes (Poisson
and bursty Markov-modulated arrivals). Avoid treating one hand-crafted
scenario per claim as general evidence.

Scenario manifests specify demand process, rates by origin/destination,
pedestrian walking-speed distribution, vehicle type mix, emergency event
schedule, sensor-error model, warm-up/evaluation length, and seed.
Expected explanations are tested by semantic slots, not exact prose.

\section{18. Baselines}\label{baselines}

{\def\LTcaptype{none} % do not increment counter
\begin{longtable}[]{@{}
  >{\raggedright\arraybackslash}p{(\linewidth - 6\tabcolsep) * \real{0.2500}}
  >{\raggedright\arraybackslash}p{(\linewidth - 6\tabcolsep) * \real{0.2500}}
  >{\raggedright\arraybackslash}p{(\linewidth - 6\tabcolsep) * \real{0.2500}}
  >{\raggedright\arraybackslash}p{(\linewidth - 6\tabcolsep) * \real{0.2500}}@{}}
\toprule\noalign{}
\begin{minipage}[b]{\linewidth}\raggedright
ID
\end{minipage} & \begin{minipage}[b]{\linewidth}\raggedright
Controller
\end{minipage} & \begin{minipage}[b]{\linewidth}\raggedright
What it isolates
\end{minipage} & \begin{minipage}[b]{\linewidth}\raggedright
Calibration rule
\end{minipage} \\
\midrule\noalign{}
\endhead
\bottomrule\noalign{}
\endlastfoot
B0 & Fixed-time & Conventional operational reference & Webster/local
timing followed by equal calibration budget; all safety timings
identical \\
B1 & Efficiency-only adaptive & Benefit/cost of adaptivity without
fairness mechanism & Minimize predicted total person delay or use
total-delay max pressure; same horizon/action set \\
B2 & Equal-weight bargaining & Closest game-theory structure & Implement
Wang-style/equal-weight bargaining as faithfully as geometry allows;
disclose deviations \\
B3 & EFPB controller, no user explanation & Incremental control/fairness
effect & Same tuned parameter budget as B1/B2 \\
B4 & EFPB + factual explanation & Interface effect only & Controller and
decisions byte-identical to B3 \\
B5 & EFPB + factual + counterfactual & Additional
contrastive-information effect & Same controller/decisions as B3/B4 \\
\end{longtable}
}

B0--B3 are technical simulation comparators. B3--B5 do \textbf{not} need
separate traffic runs because explanations cannot affect SUMO; reusing
the same traces is a design requirement. The HCI experiment compares UI
presentations corresponding to B3--B5. If time is limited, retain B0,
B1, and B3 technically, but do not drop the decision-only HCI condition.

Apply the same safety kernel, pedestrian timing, candidate duration grid
where applicable, demand manifests, and parameter-selection budget.
Comparing a well-tuned proposal to an untuned default is not valid. The
fixed-time plan may use different logic but must receive documented
calibration on calibration scenarios only.

\section{19. HCI Experiment Design}\label{hci-experiment-design}

\subsection{19.1 Design and participants}\label{design-and-participants}

Use a \(3\text{ (explanation)}\times2\text{ (decision-quality)}\)
repeated-measures
design. Explanation has D, F, and FC levels. Decision quality is sound
EFPB output versus a physically safe but procedurally anomalous output
generated by a controlled fault-injection variant (for example ignoring
a stale count or disabling waiting debt). Explanations in anomaly trials
must faithfully reveal the controller state/rule that produced the
anomaly; never display a fabricated justification.

Recruit 90 usable adults aged 18 or older, aiming for diversity in age,
gender, mobility, driving frequency, pedestrian/transit frequency,
education, and traffic-system familiarity. Recruit approximately
102--108 to absorb preregistered attention/technical exclusions. This is
a planning target, not a substitute for power analysis. After pilot
estimates, simulate power for the smallest primary within-person effect
worth detecting (initial planning value around \(d_z = .30\)) under the
intended mixed models; preregister the achieved target. Oversample
people who regularly walk and drive if feasible; analyze experience as a
measured covariate rather than assigning participants to pretend to be a
driver while actively driving.

\subsection{19.2 Stimuli and
counterbalancing}\label{stimuli-and-counterbalancing}

Each participant completes three six-trial blocks, one per explanation
condition: 18 scored trials total, balanced across quality
(\texttt{3\ sound\ +\ 3\ anomalous} per block) and scenario families.
Create three matched scenario sets with comparable demand, mode
selected, emergency frequency, and difficulty. Rotate condition-to-set
and block order with a Williams/Latin-square schedule. Participants
never see the same exact trace twice.

The layout, visual hierarchy, animation duration, signal state, data
panels, and available viewing time are identical across conditions. In
D, retain the rationale-card footprint with a neutral ``No rationale
shown'' label so layout change does not become the manipulation. Log
disclosure clicks in FC; keep the counterfactual visible by default if
the intended claim concerns exposure rather than optional use.

\subsection{19.3 Procedure}\label{procedure}

\begin{enumerate}
\def\labelenumi{\arabic{enumi}.}
\tightlist
\item
  consent, eligibility, demographics, mobility/traffic experience,
  accessibility needs, and baseline automation familiarity;
\item
  standardized tutorial explaining that the display is an audit
  prototype, not a crossing signal;
\item
  three unscored practice trials covering ordinary, emergency, and
  fallback decisions;
\item
  three counterbalanced condition blocks; each trial shows a 10--15 s
  replay, then objective questions and brief ratings;
\item
  post-block trust and fairness instruments; SUS once after all blocks,
  referring to the overall inspection interface;
\item
  preference/ranking and optional open-ended explanation critique;
\item
  debrief that some physically safe decisions were intentionally
  generated by faulty rules, with the purpose explained.
\end{enumerate}

The trial timer starts when the replay ends and all response controls
become stable; it stops when the participant submits the final objective
answer. Also log first response, changes, counterfactual interaction,
focus loss, and network/rendering timing. Use the browser monotonic
clock (\texttt{performance.now()}), not server wall time.

\subsection{19.4 Trial tasks}\label{trial-tasks}

\begin{itemize}
\tightlist
\item
  identify which road-user group received priority;
\item
  choose the primary reason from one correct and plausible distractors;
\item
  identify whether a safety/emergency/fallback constraint determined the
  choice;
\item
  answer which specified change would most likely flip the decision;
\item
  rate decision appropriateness and confidence (\(0–100\));
\item
  rate brief trial-level transparency/fairness items; full scales appear
  after the block.
\end{itemize}

Primary comprehension outcomes are the factual-reason and contrastive
answers. Priority identification is a manipulation/sanity check because
the signal itself should make it easy.

\subsection{19.5 Instruments}\label{instruments}

\textbf{Trust.} Use the Jian, Bisantz, and Drury Trust in Automated
Systems Scale: 12 items rated on a 7-point response scale, with five
negatively worded distrust items reverse-scored and items
averaged/summed according to the preregistered version. The authors
conceptualized trust/distrust poles along a global dimension. Administer
after each condition block, replacing only the generic referent
``system'' with the intersection decision system; disclose that
adaptation.\hyperref[source-28]{[28]}

\textbf{Perceived fairness.} Adapt only the relevant procedural,
distributive, and informational justice items from Colquitt's validated
four-factor organizational-justice measure. It originally distinguishes
procedural, distributive, interpersonal, and informational justice; this
traffic adaptation is \textbf{not automatically
validated}.\hyperref[source-29]{[29]} Use cognitive interviews, expert
content review, translation/back-translation if needed, a pilot
factor/reliability check, and report McDonald's omega. Do not call a
shortened modified scale ``validated'' without new evidence. Example
constructs, not final items: consistent application of rules
(procedural), balance of delay burdens (distributive), and
adequacy/truthfulness of reasons (informational).

\textbf{Usability.} Administer Brooke's 10-item, five-response SUS after
the full prototype. For odd items use \(response-1\); for even items use
\(5-response\); sum and multiply by \texttt{2.5} to yield \(0–100\).
Interpret only the composite, not individual items, and report
comparisons/uncertainty rather than treating 68 as a universal pass
line.\hyperref[source-30]{[30]}

\textbf{Explanation quality.} Objective comprehension remains primary.
The Explanation Satisfaction Scale or selected preregistered dimensions
from Hoffman et al.~may be secondary if licensing/wording is confirmed;
it should not replace performance measures.\hyperref[source-31]{[31]}

\subsection{19.6 Ethics and
accessibility}\label{ethics-and-accessibility}

Obtain institutional ethics/IRB approval before recruitment. Explain
recording, compensation, withdrawal, data retention, and anomaly
deception/debriefing. Store contact/compensation identifiers separately
from response data. Provide keyboard access, non-color encodings,
captions/text for any audio, reduced-motion mode, readable contrast, and
a comprehension-tested plain-language glossary. The study must not ask
anyone to interact while operating a real vehicle or crossing a real
road.

\section{20. Evaluation Metrics}\label{evaluation-metrics}

\subsection{20.1 Technical traffic
metrics}\label{technical-traffic-metrics}

For completed road users \(i\) of mode \(m\) in evaluation interval
\(T\):

\begin{itemize}
\tightlist
\item
  Mean wait:
  \[
    \overline{W}_m = \frac{1}{N_m}\sum_{i=1}^{N_m} W_i
    \quad [\mathrm{s}].
  \]
\item
  Tail wait:
  \[
    W_m^{95} = Q_{0.95}\!\left(\{W_i\}_{i=1}^{N_m}\right);
  \]
  bootstrap its uncertainty at run level.
\item
  Maximum wait:
  \[
    W_m^{\max} = \max_{1\le i\le N_m} W_i;
  \]
  also report exceedance counts because maxima are unstable.
\item
  Total delay:
  \[
    D_m = \sum_{i=1}^{N_m}
    \left(T_i^{\mathrm{travel}}-T_i^{\mathrm{free\ flow}}\right),
  \]
  measured in person-seconds, or vehicle-seconds when occupancy is
  unavailable.
\item
  Time-average queue:
  \[
    \overline{Q}_m = \frac{1}{|T|}\int_T Q_m(t)\,\mathrm{d}t,
  \]
  approximated by the 1~Hz sum.
\item
  Throughput:
  \[
    X_m = \frac{N_m^{\mathrm{exit}}}{|T|/3600}
    \quad [\mathrm{users/h}].
  \]
\item
  Service ratio:
  \[
    z_m = \frac{N_m^{\mathrm{served}}}{N_m^{\mathrm{requested}}}.
  \]
\item
  Starvation violations:
  \[
    V_m = \sum_{e\in\mathcal{E}_m}
      \mathbf{1}\!\left[W_e>T_m^{\max}\right],
  \]
  plus seconds above the limit.
\item
  Phase switches:
  \[
    S = \text{number of service-phase changes};
  \]
  transition subphases are not double-counted.
\item
  Emergency response:
  \[
    R_e = t_{\mathrm{compatible\ green}}
          -t_{\mathrm{authenticated\ request}};
  \]
  report its distribution and mandatory-clearance component.
\item
  Safety violations: count of conflict-matrix violations, truncated
  clearances, red-entry events attributable to controller, collisions,
  and invalid commands. Expected software invariant count is zero;
  simulation surrogate safety does not imply field safety.
\item
  Controller latency: wall-clock \(t_{\mathrm{return}}-t_{\mathrm{call}}\)
  per admissible
  epoch, with mean/p95/p99 and deadline misses.
\item
  Fairness: \(J\), \(G\), \(G^{95}\), mode
  distributions, and Pareto position from Section 13.
\item
  Efficiency-bound compliance:
  \[
    B = \frac{L_{\mathrm{eff}}(a^*)-L_{\min}}
             {\max\!\left(L_{\min},L_{\mathrm{floor}}\right)};
  \]
  count \(B>\delta\) outside override/fallback cases.
\end{itemize}

Primary technical endpoints: total person delay (or separate mode delay
if occupancy unknown), pedestrian 95th-percentile wait, vehicle
95th-percentile wait, maximum normalized wait, and starvation
violations. Predeclare one fairness endpoint and one efficiency endpoint
for confirmatory claims; treat the rest as explanatory.

\subsection{20.2 Explanation metrics}\label{explanation-metrics}

\begin{itemize}
\tightlist
\item
  Record fidelity rate:
  \texttt{valid\ explanations\ /\ attempted\ explanations}.
\item
  Counterfactual validity rate:
  \texttt{replayed\ flips\ /\ displayed\ counterfactuals}.
\item
  Constraint precision:
  \texttt{correctly\ claimed\ binding\ constraints\ /\ all\ claimed\ binding\ constraints}.
\item
  Coverage by decision class: ordinary, emergency, degraded, and
  recovery.
\item
  Stability: mean Jaccard similarity of top-reason sets under
  non-decision-changing sensor perturbations.
\item
  Explanation latency: mean/p95/p99 milliseconds.
\end{itemize}

The required target for factual and counterfactual validity is 100\% in
released stimuli. Coverage may be below 100\% if the UI honestly
abstains; report abstention rather than filling gaps with invented
prose.

\subsection{20.3 HCI metrics}\label{hci-metrics}

\begin{itemize}
\tightlist
\item
  Factual comprehension \(Y_{\mathrm{fact}}\in\{0,1\}\) from the primary-reason
  question.
\item
  Contrastive comprehension \(Y_{\mathrm{cf}}\in\{0,1\}\) from the valid
  decision-flip question.
\item
  Constraint comprehension \(Y_{\mathrm{con}}\in\{0,1\}\).
\item
  Response time \texttt{RT} in milliseconds from stable question display
  to final submission; primary analysis uses \texttt{log(RT)} on
  correct, focused trials and a sensitivity analysis includes all valid
  trials.
\item
  Trial confidence \(c\in[0,1]\) and correctness \(y\in\{0,1\}\); Brier
  score \(\mathrm{BS}=(c-y)^2\), lower is better.
\item
  Appropriate endorsement: endorse a sound decision or reject an
  anomalous one.
\item
  Quality discrimination: within-person difference in
  appropriateness/confidence between sound and anomalous decisions;
  optional AUROC treats confidence/endorsement as the score.
\item
  Trust: reverse-code designated Jian items, then compute the
  preregistered composite and reliability.
\item
  Fairness: procedural, distributive, and informational adapted
  composites; report each separately and their omega.
\item
  SUS: standard composite \(0\text{--}100\).
\item
  Acceptance/preference: secondary, clearly identified as custom or
  adapted.
\end{itemize}

\section{21. Statistical Analysis}\label{statistical-analysis}

\subsection{21.1 Simulation analysis}\label{simulation-analysis}

The experimental unit is a complete run, not a one-second record.
Analyze scenario-level run metrics with controller as fixed effect and
scenario plus paired seed as blocking/random effects. A suitable model
is

\texttt{metric\ \textasciitilde{}\ controller\ *\ scenario\_family\ +\ (1\textbar{}seed)},

adding demand levels and process where factorial. For skewed positive
metrics use log/Gamma models; for violation counts use negative-binomial
or hurdle models; for zero software safety violations report exact
counts and invariant-test coverage. Present paired controller
differences per seed with 95\% cluster/bootstrap confidence intervals.
Check residuals and influence; use rank-based paired sensitivity if
assumptions fail. Apply Holm correction across the preregistered primary
controller contrasts.

Do not treat 3,600 time points from one run as 3,600 independent
observations. If time-series behavior is analyzed, use block bootstrap
or explicit autocorrelation and label it secondary.

\subsection{21.2 HCI analysis}\label{hci-analysis}

Independent variables are explanation condition (\texttt{D/F/FC}),
decision quality (\texttt{sound/anomalous}), scenario family, and their
preregistered interactions. Covariates include block/order, scenario
set, traffic familiarity, walking/driving frequency, and accessibility
status only when theoretically justified. Do not data-mine demographic
interactions.

Fit:

\begin{itemize}
\tightlist
\item
  factual and contrastive correctness: logistic generalized linear mixed
  models;
\item
  log response time: linear mixed model;
\item
  appropriate endorsement: logistic mixed model;
\item
  appropriateness/confidence: cumulative-link mixed model for ordinal
  ratings or a preregistered continuous sensitivity model;
\item
  Brier score and composite trust/fairness: linear mixed models if
  diagnostics support them, otherwise robust/ordinal alternatives.
\end{itemize}

The maximal intended random structure is participant and scenario
intercepts with participant slopes for explanation/quality:
\texttt{(1 + explanation + quality \textbar{} participant) +
(1 \textbar{} scenario)}. Simplify correlations/slopes by a preregistered
sequence only if convergence fails; report the final structure. Planned
contrasts are \(F-D\), \(FC-D\), and \(FC-F\). H3/H4 depend on the
\(\text{explanation}\times\text{quality}\) interaction; a main effect
alone does not establish calibration.

Use two-sided \(\alpha=0.05\), 95\% confidence intervals, and Holm adjustment
within the five hypothesis families. Report odds ratios for logistic
models, standardized mean differences or model-standardized coefficients
for continuous outcomes, estimated marginal means, and raw descriptive
distributions. A statistically significant tiny effect is not
automatically useful; report smallest-effect-of-interest or equivalence
intervals chosen before data collection.

\subsection{21.3 Data quality and
missingness}\label{data-quality-and-missingness}

Preregister exclusions: duplicate identity, ineligibility, failure of
both tutorial checks after retraining, technical failure affecting
stimulus, implausibly fast completion defined from pilot rather than
convenience, and instructed-response failure. Do not exclude someone for
disagreeing with the controller. Preserve all valid trials from included
participants; model isolated missing responses where possible and report
reasons. Analyze a full valid sample plus a preregistered
attention-compliant sensitivity set.

Report missingness by condition, reliability with McDonald's omega and
confidence intervals, item distributions, manipulation checks, and all
deviations from preregistration. If the adapted justice scale fails its
intended structure, analyze items or defensible subscales exploratorily
and narrow claims.

\section{22. Ablation Studies}\label{ablation-studies}

{\def\LTcaptype{none} % do not increment counter
\begin{longtable}[]{@{}
  >{\raggedright\arraybackslash}p{(\linewidth - 6\tabcolsep) * \real{0.2500}}
  >{\raggedright\arraybackslash}p{(\linewidth - 6\tabcolsep) * \real{0.2500}}
  >{\raggedright\arraybackslash}p{(\linewidth - 6\tabcolsep) * \real{0.2500}}
  >{\raggedright\arraybackslash}p{(\linewidth - 6\tabcolsep) * \real{0.2500}}@{}}
\toprule\noalign{}
\begin{minipage}[b]{\linewidth}\raggedright
Ablation
\end{minipage} & \begin{minipage}[b]{\linewidth}\raggedright
Change from EFPB
\end{minipage} & \begin{minipage}[b]{\linewidth}\raggedright
Primary question
\end{minipage} & \begin{minipage}[b]{\linewidth}\raggedright
Essential outcomes
\end{minipage} \\
\midrule\noalign{}
\endhead
\bottomrule\noalign{}
\endlastfoot
A1 Efficiency-only & Remove bargaining/debt; choose \texttt{L\_min} & Is
fairness machinery doing anything? & Delay, tails, starvation,
disparity \\
A2 No starvation constraint & Retain bargaining but remove hard wait cap
& Does bargaining alone protect rare/high-wait users? & Limit violations
and maximum wait \\
A3 No efficiency filter & Bargain across all safe plans & Does
\(\mathcal{A}_\delta\)
prevent unacceptable efficiency sacrifice? & Total delay, switches,
Pareto position \\
A4 Equal/static weights & Set \(\kappa=0\), equal \(\pi{}\) & Does temporal
service debt matter? & Repeated-imbalance recovery, tails \\
A5 Weighted sum & Replace Nash log objective with tuned linear utility &
Is bargaining needed beyond a simple score? & Same outcomes plus
decision agreement \\
A6 Sensor-quality fallback off & Continue ordinary optimization with
corrupted state & Does conservative fallback reduce unsafe/misleading
behavior? & invalid decisions, performance, abstention \\
A7 Decision-only UI & Hide explanation & Does rationale affect people
while traffic is unchanged? & H1, H3, H4 \\
A8 Factual-only UI & Remove counterfactual from FC & Does contrastive
content add benefit/cost? & H2, H5 \\
\end{longtable}
}

A1--A5 are the minimal technical set for a strong algorithm paper,
although A6 may be restricted to targeted robustness scenarios. A7--A8
are HCI manipulations using the same EFPB trace, not separate
controllers. Avoid a combinatorial experiment across every component;
one-factor removals and a small parameter grid are sufficient.

Sensitivity analyses vary \(\delta\), \(\kappa\), \(\eta\), horizon \texttt{H},
normalization scale, arrival-window length, pedestrian walking-speed
distribution, vehicle occupancy, sensor error, and random seed. Show
whether the recommended configuration lies on a stable region rather
than a narrow optimum.

\section{23. Implementation
Architecture}\label{implementation-architecture}

\subsection{23.1 Minimal stack}\label{minimal-stack}

\begin{itemize}
\tightlist
\item
  Python 3.12+ for domain logic, TraCI/libsumo, FastAPI, Pydantic,
  NumPy/Pandas/Polars, PyArrow, and testing.
\item
  Eclipse SUMO pinned to one release; the official download page and
  Python tooling support \texttt{eclipse-sumo}, \texttt{traci},
  \texttt{sumolib}, and \texttt{libsumo}
  packages.\hyperref[source-32]{[32]}
\item
  React + TypeScript + Vite for the dashboard; SVG/CSS for the
  intersection. A canvas/game engine is unnecessary for one junction.
\item
  FastAPI REST for run control/replay metadata and WebSocket for compact
  live frames.
\item
  JSON/YAML for configuration, JSONL for lossless nested traces, Parquet
  for analytical tables, and CSV only for portable summaries.
\item
  R with \texttt{lme4}, \texttt{ordinal}, \texttt{emmeans},
  \texttt{performance}, and \texttt{effectsize} for confirmatory mixed
  models; Python handles preprocessing and technical metrics.
\item
  \texttt{pytest}, mypy/pyright, Ruff, Vitest, Playwright, and schema
  tests. Docker is optional after the local pipeline works.
\end{itemize}

The backend is one process plus one SUMO child process per run. Do not
introduce microservices, message queues, cloud orchestration, a language
model, or RL. Batch experiments launch isolated workers only after
verifying that seeds and output paths cannot collide.

\subsection{23.2 Algorithms}\label{algorithms}

\textbf{A1 --- Traffic-state collection}

\begin{verbatim}
INPUT subscriptions, previous registries, timestamp t
read current person/vehicle/detector/SPaT values
update first-seen, waiting-start, service, and exit registries
compute counts, rolling arrivals, queue lengths, mean/max waits, density
attach source timestamp and quality to every aggregate
state ← validate(IntersectionState(...))
RETURN state
\end{verbatim}

\textbf{A2 --- Pedestrian advocate}

\begin{verbatim}
INPUT state, candidate plans, configuration
debt_p ← clip(max_ped_wait / ped_wait_limit, 0, 1)
weight_p_raw ← base_p * (1 + kappa * debt_p^eta)
FOR plan IN candidates:
    served_p ← predict_ped_discharge(plan, state)
    utility_p[plan] ← min(1, served_p / max(1, ped_horizon_demand))
    facts_p[plan] ← typed factors used in prediction and weight
RETURN utility_p, weight_p_raw, facts_p
\end{verbatim}

\textbf{A3 --- Vehicle advocate}

\begin{verbatim}
INPUT state, candidate plans, configuration
compute cumulative-delay pressure for NS and EW
debt_v ← clip(max_vehicle_wait / vehicle_wait_limit, 0, 1)
weight_v_raw ← base_v * (1 + kappa * debt_v^eta)
FOR plan IN candidates:
    direction ← legal argmax(NS_pressure, EW_pressure) where vehicle slot exists
    served_v ← predict_vehicle_discharge(plan, direction, state)
    utility_v[plan] ← min(1, served_v / max(1, vehicle_horizon_demand))
RETURN utility_v, weight_v_raw, direction, facts_v
\end{verbatim}

\textbf{A4 --- Utility and loss calculation}

\begin{verbatim}
FOR plan IN safe_candidates:
    simulate deterministic queue discharge over horizon H
    calculate mode service fractions u_p, u_v
    calculate predicted person/mode delay and switch count
    loss[plan] ← weighted_delay + switch_cost
    assert every value has unit and provenance
RETURN CandidateEvaluation[]
\end{verbatim}

\textbf{A5 --- Fairness calculation}

\begin{verbatim}
INPUT completed/requested/service events for one run
calculate waits by mode and normalized mean/tail disparity
calculate service ratios z_p, z_v
IF both ratios defined and not both zero: Jain ← (z_p+z_v)^2/(2*(z_p^2+z_v^2))
count wait-limit violations and duration above limits
RETURN FairnessResult plus distributions (never Jain alone)
\end{verbatim}

\textbf{A6 --- Priority negotiation}

\begin{verbatim}
INPUT state, evaluations, fallback utilities, config
safe ← candidates passing SignalSafetyKernel predicates
IF emergency: RETURN emergency_policy(safe)
IR ← plans with each utility at least its fallback utility
bounded ← plans satisfying predicted wait limits
IF bounded is empty: RETURN least_breach_lexicographic(IR or safe)
Lmin ← minimum efficiency loss in bounded
near ← plans with loss <= Lmin + delta*max(Lmin, Lfloor)
normalize advocate weights
score each plan by sum_m weight_m*log(utility_m-disagreement_m+epsilon)
winner ← deterministic argmax(score, tail_wait, switches, candidate_id)
RETURN decision and complete candidate table
\end{verbatim}

\textbf{A7 --- Emergency override}

\begin{verbatim}
IF no authenticated emergency: RETURN none
identify requested compatible vehicle group
finish any mandatory current pedestrian/vehicle clearance
transition through configured yellow/all-red sequence
serve emergency phase until release condition or maximum policy duration
record response components and displaced-service debt
enter recovery; never clear or reset ordinary wait histories
\end{verbatim}

\textbf{A8 --- Signal controller}

\begin{verbatim}
INPUT chosen target, current SignalMachineState
IF transition not admissible: hold and return next admissible time
path ← prevalidated transition graph[current_service][target_service]
FOR each transition state:
    apply state/duration through TraCI
    assert conflict matrix and minimum interval
RETURN execution receipt
\end{verbatim}

\textbf{A9 --- Explanation generator}

\begin{verbatim}
INPUT immutable DecisionTrace, audience, explanation condition
decision sentence ← template(trace.chosen_plan)
reasons ← top valid trace factors, maximum 3
alternative sentence ← template(trace.best_rejected, trace.rejection_reason)
IF condition includes counterfactual:
    cf ← A10(trace)
render canonical slots; do not add unbound claims
validation ← check values, predicates, hashes, and replay
RETURN explanation if valid ELSE abstention record
\end{verbatim}

\textbf{A10 --- Counterfactual search}

\begin{verbatim}
INPUT trace, target alternative, permitted feature steps/ranges
FOR distance FROM 1 TO max_distance:
    FOR perturbation IN deterministic_order(all combinations at distance):
        candidate_state ← apply perturbation to controllable fields only
        IF state constraints fail: CONTINUE
        replay ← controller(candidate_state, same config/version)
        IF replay.choice == alternative: RETURN perturbation + replay proof
RETURN no_counterfactual_within_range
\end{verbatim}

Use a priority queue rather than materializing all combinations;
terminate at the first valid distance under deterministic ordering.

\textbf{A11 --- Scenario generator}

\begin{verbatim}
INPUT ScenarioConfig, seed
rng ← seeded independent streams for vehicles, pedestrians, speeds, emergencies, sensors
sample complete arrival/event manifests before any controller runs
materialize SUMO routes/person trips and event JSON
hash manifest and validate counts, routes, walking connectivity, time bounds
RETURN immutable ScenarioManifest
\end{verbatim}

\textbf{A12 --- Experiment runner and data logger}

\begin{verbatim}
FOR scenario IN preregistered scenarios:
  FOR seed IN seed_list:
    manifest ← load or generate once
    FOR controller IN randomized controller order:
      create unique run directory; save resolved config and hashes
      launch SUMO; execute A1→A8 each step
      append state batches to Parquet and traces/events to JSONL atomically
      close SUMO; validate outputs; write RunSummary(status)
aggregate only valid runs; retain invalid-run audit log
\end{verbatim}

\textbf{A13 --- Evaluation metric calculator}

\begin{verbatim}
INPUT validated run directory
read demand manifest, SUMO trip/person output, states, decisions, events
cross-check counts and time window
compute per-user waits/delays, run-level traffic/fairness/latency metrics
replay a sample or all explanation records and counterfactuals
write tidy run_metrics.parquet and validation_report.json
RETURN success only when schema and invariant checks pass
\end{verbatim}

\subsection{23.3 What must be researched versus
built}\label{what-must-be-researched-versus-built}

{\def\LTcaptype{none} % do not increment counter
\begin{longtable}[]{@{}
  >{\raggedright\arraybackslash}p{(\linewidth - 10\tabcolsep) * \real{0.350}}
  >{\centering\arraybackslash}p{(\linewidth - 10\tabcolsep) * \real{0.130}}
  >{\centering\arraybackslash}p{(\linewidth - 10\tabcolsep) * \real{0.100}}
  >{\centering\arraybackslash}p{(\linewidth - 10\tabcolsep) * \real{0.160}}
  >{\centering\arraybackslash}p{(\linewidth - 10\tabcolsep) * \real{0.130}}
  >{\centering\arraybackslash}p{(\linewidth - 10\tabcolsep) * \real{0.130}}@{}}
\toprule\noalign{}
\begin{minipage}[b]{\linewidth}\raggedright
Component
\end{minipage} & \begin{minipage}[b]{\linewidth}\centering
Research first
\end{minipage} & \begin{minipage}[b]{\linewidth}\centering
Code
\end{minipage} & \begin{minipage}[b]{\linewidth}\centering
External dataset
\end{minipage} & \begin{minipage}[b]{\linewidth}\centering
Simulation
\end{minipage} & \begin{minipage}[b]{\linewidth}\centering
User study
\end{minipage} \\
\midrule\noalign{}
\endhead
\bottomrule\noalign{}
\endlastfoot
Local signal/\allowbreak pedestrian timing and emergency rules & Y &
Config/\allowbreak tests &
N & Y & N \\
Simplified intersection/conflict matrix & Y & Y & Optional & Y & N \\
Arrival/walking/discharge calibration & Y & Y & Recommended & Y & N \\
Fixed-time and efficiency baselines & Y & Y & N & Y & N \\
Bargaining utilities, disagreement, and efficiency tolerance & Y & Y &
Calibration only & Y & N \\
Objective fairness metrics/service limits & Y & Y & Optional & Y & N \\
Explanation content and wording & Y & Y & N & Replay source & Pilot +
confirmatory \\
Counterfactual algorithm/fidelity tests & Y & Y & N & Controller replay
& Y \\
Dashboard & Formative HCI & Y & N & Replay/live & Y \\
Trust/\allowbreak fairness/\allowbreak usability instruments & Y &
Scoring/\allowbreak rendering & N & N &
Y \\
Statistical models/power & Y & Y & Pilot responses & Run outputs & Y \\
Deployment-grade sensing/V2X/security & Future research & No in v1 &
Would be required & Fault model only & N \\
\end{longtable}
}

``Research'' here includes standards review, stakeholder interviews,
parameter calibration, and measurement validation. It is not a license
to choose values informally. Conversely, implementation work such as
React styling is engineering evidence only when it enables a controlled
or reproducible evaluation.

\section{24. Coding Roadmap}\label{coding-roadmap}

\subsection{Task 01 --- Project bootstrap and
intersection}\label{task-01-project-bootstrap-and-intersection}

\begin{itemize}
\tightlist
\item
  \textbf{Purpose:} establish a reproducible package and a
  conflict-checked three-phase SUMO network.
\item
  \textbf{Files to create:} \path|pyproject.toml|, lock file,
  \path|.gitignore|, \path|configs/intersection.yaml|,
  \path|sumo/net/intersection.net.xml|,
  \path|sumo/tls/base_tls.add.xml|,
  \path|sumo/detectors/detectors.add.xml|, \path|sumo/base.sumocfg|,
  initial test fixtures.
\item
  \textbf{Files to modify:} none initially; later modify only through
  reviewed config/network changes.
\item
  \textbf{Technology:} Python, SUMO/NetEdit/netconvert, YAML, pytest.
\item
  \textbf{Inputs/outputs:} geometry/standard requirements → runnable
  network, conflict matrix, timing configuration.
\item
  \textbf{Algorithm/pseudocode:} construct links and phases;
  \texttt{for\ each\ phase\ pair:\ assert\ no\ foe\ links\ are\ simultaneously\ permissive}.
\item
  \textbf{Implementation notes:} keep turns disabled; store timing
  source/provenance beside each value; give signal links stable names.
\item
  \textbf{Test criteria:} SUMO validation passes; vehicles/persons route
  end to end; all transitions visually checked; automated
  conflict/clearance tests pass.
\item
  \textbf{Expected result:} a deterministic fixed-time smoke run with no
  controller-induced collision, teleport, or invalid phase.
\item
  \textbf{How to run:} \texttt{sumo-gui -c}
  \path|sumo/base.sumocfg| and \texttt{pytest}
  \path|tests/integration/test_network.py|.
\end{itemize}

\subsection{Task 02 --- State collector and sensor
abstraction}\label{task-02-state-collector-and-sensor-abstraction}

\begin{itemize}
\tightlist
\item
  \textbf{Purpose:} create a typed, unit-safe state independent of raw
  TraCI calls.
\item
  \textbf{Files to create:}
  \path|src/intersection_priority/domain/models.py|,
  \path|sensing/traci_adapter.py|,
  \path|sensing/state_collector.py|,
  \path|sensing/normalization.py|, \path|sensing/noise.py|, unit
  tests.
\item
  \textbf{Files to modify:} \texttt{configs/intersection.yaml},
  \texttt{pyproject.toml}.
\item
  \textbf{Technology:} TraCI subscriptions, Pydantic/dataclasses,
  deque-based rolling windows.
\item
  \textbf{Inputs/outputs:} SUMO subscription frames + registry →
  \texttt{IntersectionState} each second.
\item
  \textbf{Algorithm/pseudocode:} A1; update object histories before
  aggregate features and attach freshness/quality.
\item
  \textbf{Implementation notes:} use SI units internally; never infer a
  zero count from missing data; distinguish empty, missing, and stale.
\item
  \textbf{Test criteria:} hand-calculated fixtures match;
  duplicated/disappearing IDs handled; stale/missing values trigger
  correct status.
\item
  \textbf{Expected result:} replaying the same subscription fixture
  yields byte-equivalent canonical JSON.
\item
  \textbf{How to run:}
  \texttt{pytest\ tests/unit/test\_state\_collector.py}.
\end{itemize}

\subsection{Task 03 --- Advocates and horizon
predictor}\label{task-03-advocates-and-horizon-predictor}

\begin{itemize}
\tightlist
\item
  \textbf{Purpose:} compute mode demand, service debt, candidate
  utility, and direction pressure.
\item
  \textbf{Files to create:} \texttt{agents/pedestrian.py},
  \texttt{agents/vehicle.py}, \texttt{prediction/discharge.py},
  \texttt{prediction/calibration.py}, tests.
\item
  \textbf{Files to modify:} controller YAML schemas.
\item
  \textbf{Technology:} pure Python/NumPy; no learning framework.
\item
  \textbf{Inputs/outputs:} state + candidate plans + calibrated
  saturation/walking parameters → utilities, weights, predictions,
  facts.
\item
  \textbf{Algorithm/pseudocode:} A2--A4.
\item
  \textbf{Implementation notes:} keep functions pure; retain unit
  annotations; compare predictions to held-out SUMO rollouts.
\item
  \textbf{Test criteria:} utility is bounded; more feasible service
  never lowers its mode utility; prediction error is reported by demand
  stratum.
\item
  \textbf{Expected result:} candidate evaluation table with traceable
  intermediate terms.
\item
  \textbf{How to run:} \texttt{pytest}
  \path|tests/unit/test_advocates.py|
  \path|tests/integration/test_predictor.py|.
\end{itemize}

\subsection{Task 04 --- Safety, baselines, and EFPB
controller}\label{task-04-safety-baselines-and-efpb-controller}

\begin{itemize}
\tightlist
\item
  \textbf{Purpose:} implement the legal transition kernel, B0--B3,
  emergency/recovery, and deterministic bargaining.
\item
  \textbf{Files to create:} \texttt{controllers/base.py},
  \texttt{fixed\_time.py}, \texttt{efficiency.py},
  \texttt{equal\_bargaining.py}, \texttt{efpb.py},
  \texttt{signal\_safety.py}, \texttt{fairness/weights.py},
  \texttt{fairness/metrics.py}, tests.
\item
  \textbf{Files to modify:} \texttt{configs/controllers/*.yaml}.
\item
  \textbf{Technology:} finite-state machine, finite candidate
  enumeration, NumPy.
\item
  \textbf{Inputs/outputs:} typed state/evaluations →
  \texttt{PriorityDecision} and \texttt{ExecutionReceipt}.
\item
  \textbf{Algorithm/pseudocode:} A5--A8.
\item
  \textbf{Implementation notes:} safety kernel has no
  controller-specific branches; log empty feasible sets; deterministic
  tie-breaks; never reset wait debt after emergency.
\item
  \textbf{Test criteria:} property tests generate arbitrary action
  sequences with zero forbidden conflicts; known fixtures select
  expected winners; all baselines share timing constraints.
\item
  \textbf{Expected result:} paired smoke scenarios complete with
  reproducible decisions and no invariant violations.
\item
  \textbf{How to run:} \texttt{pytest}
  \path|tests/unit/test_bargaining.py|
  \path|tests/integration/test_signal_safety.py|.
\end{itemize}

\subsection{Task 05 --- Trace, explanation, and counterfactual
validation}\label{task-05-trace-explanation-and-counterfactual-validation}

\begin{itemize}
\tightlist
\item
  \textbf{Purpose:} make every displayed reason auditable and every
  counterfactual decision-valid.
\item
  \textbf{Files to create:} \texttt{explanation/trace.py},
  \texttt{generator.py}, \texttt{templates.py},
  \texttt{counterfactual.py}, \texttt{validators.py}, golden
  fixtures/tests.
\item
  \textbf{Files to modify:} \texttt{domain/models.py}, controller code
  to emit complete traces.
\item
  \textbf{Technology:} typed Python, Jinja2 or format templates
  restricted to declared slots, canonical JSON hashing.
\item
  \textbf{Inputs/outputs:}
  \texttt{DecisionTrace} + condition \(\rightarrow\)
  \texttt{ExplanationRecord}, or
  explicit abstention.
\item
  \textbf{Algorithm/pseudocode:} A9--A10.
\item
  \textbf{Implementation notes:} canonical text is generated after
  validation; counterfactual fields have domain ranges and immutable
  flags.
\item
  \textbf{Test criteria:} mutation tests catch wrong numbers/directions;
  all released counterfactuals replay-flip; golden text changes only
  with approved template version.
\item
  \textbf{Expected result:} 100\% fidelity for ordinary study stimuli
  and honest abstention for unsupported cases.
\item
  \textbf{How to run:} \texttt{pytest}
  \path|tests/unit/test_explanations.py|
  \path|tests/regression/test_explanation_golden.py|.
\end{itemize}

\subsection{Task 06 --- Runner, scenarios, logging, and
metrics}\label{task-06-runner-scenarios-logging-and-metrics}

\begin{itemize}
\tightlist
\item
  \textbf{Purpose:} generate common-random-number manifests and execute
  reproducible batch experiments.
\item
  \textbf{Files to create:} \path|simulation/runner.py|,
  \path|simulation/scenarios.py|, \path|simulation/logging.py|,
  \path|experiments/generate_manifests.py|,
  \path|experiments/run_batch.py|,
  \path|analysis/technical_metrics.py|, schemas/tests.
\item
  \textbf{Files to modify:} experiment and scenario YAML files.
\item
  \textbf{Technology:} SUMO, TraCI/libsumo, PyArrow/Parquet, JSONL,
  process pool after serial verification.
\item
  \textbf{Inputs/outputs:} experiment config → validated run directories
  and tidy run metrics.
\item
  \textbf{Algorithm/pseudocode:} A11--A13.
\item
  \textbf{Implementation notes:} write to a temporary run-specific file
  then atomically rename; never overwrite a completed run; record
  failures.
\item
  \textbf{Test criteria:} serial versus parallel outputs match by hash
  for a small fixture; rerunning a completed ID is refused or verified;
  online/offline metrics agree.
\item
  \textbf{Expected result:} one command executes every controller on
  identical demand streams.
\item
  \textbf{How to run:} \texttt{python}
  \path|experiments/run_batch.py| \texttt{--config}
  \path|configs/experiments/pilot.yaml| \texttt{--workers 1}.
\end{itemize}

\subsection{Task 07 --- API and
dashboard}\label{task-07-api-and-dashboard}

\begin{itemize}
\tightlist
\item
  \textbf{Purpose:} show current state, decision, reason, alternative,
  uncertainty, and history without granting unsafe control authority.
\item
  \textbf{Files to create:} \texttt{api/main.py},
  \texttt{api/schemas.py}, \texttt{api/routes.py},
  \texttt{api/websocket.py}, \texttt{dashboard/} Vite app, React
  components, frontend tests.
\item
  \textbf{Files to modify:} trace/replay serializers and CORS/dev
  config.
\item
  \textbf{Technology:} FastAPI, Uvicorn, React, TypeScript, Vite, SVG,
  WebSocket.
\item
  \textbf{Inputs/outputs:} live/replay frames → accessible interactive
  visualization and logged UI events.
\item
  \textbf{Algorithm/pseudocode:}
  \(on frame: validate schema \rightarrow{} update immutable store \rightarrow{} render; never infer missing rationale client-side\).
\item
  \textbf{Implementation notes:} data panels use units/reference bars;
  signal colors also have labels/shapes; explanation conditions share
  layout.
\item
  \textbf{Test criteria:} contract tests, keyboard/navigation checks,
  axe accessibility scan, reconnect/replay, and screenshot regression at
  common widths.
\item
  \textbf{Expected result:} a browser can inspect live or prerecorded
  decisions with stable decision IDs.
\item
  \textbf{How to run:}
  \texttt{uvicorn\ intersection\_priority.api.main:app\ -\/-reload\ -\/-port\ 8000}
  and \texttt{npm\ -\/-prefix\ dashboard\ run\ dev}.
\end{itemize}

\subsection{Task 08 --- Study stimulus and response
system}\label{task-08-study-stimulus-and-response-system}

\begin{itemize}
\tightlist
\item
  \textbf{Purpose:} freeze matched stimuli and collect trial-level
  responses/timing without live simulation confounds.
\item
  \textbf{Files to create:} \texttt{study/build\_stimuli.py},
  \texttt{study/stimuli/manifest.json}, \texttt{study/app/} or approved
  survey integration, \texttt{study/protocol.md},
  \texttt{study/scoring.py}, tests.
\item
  \textbf{Files to modify:} dashboard components to accept
  \texttt{studyMode}; API response schema if self-hosted.
\item
  \textbf{Technology:} same React renderer, immutable JSON replays,
  \texttt{performance.now()}, approved survey/secure endpoint.
\item
  \textbf{Inputs/outputs:} validated trace pool + counterbalancing
  schedule → stimulus packs and pseudonymous responses.
\item
  \textbf{Algorithm/pseudocode:} stratify traces → match
  difficulty/features → rotate set/condition with Williams schedule →
  hash and freeze.
\item
  \textbf{Implementation notes:} no exact trace repeats per participant;
  record focus/render problems; decision-quality faults must remain
  physically safe.
\item
  \textbf{Test criteria:} balance report passes; scoring key
  independently verified; condition changes only explanation fields;
  resume/duplicate protections work.
\item
  \textbf{Expected result:} deterministic 18-trial sessions with
  complete provenance.
\item
  \textbf{How to run:}
  \texttt{python\ study/build\_stimuli.py\ -\/-config\ configs/study/pilot.yaml}
  then the documented local/study host command.
\end{itemize}

\subsection{Task 09 --- Statistical analysis and
figures}\label{task-09-statistical-analysis-and-figures}

\begin{itemize}
\tightlist
\item
  \textbf{Purpose:} produce auditable technical and HCI results from
  immutable processed data.
\item
  \textbf{Files to create:} \texttt{analysis/prepare\_data.py},
  \texttt{analysis/technical\_models.R},
  \texttt{analysis/hci\_models.R}, \texttt{analysis/figures.R},
  \texttt{analysis/report.qmd}, tests/synthetic fixtures.
\item
  \textbf{Files to modify:} preregistered contrast/config files only
  before data lock.
\item
  \textbf{Technology:} Python validation; R \texttt{lme4},
  \texttt{ordinal}, \texttt{emmeans}, \texttt{boot}, Quarto.
\item
  \textbf{Inputs/outputs:} validated Parquet/CSV → model objects,
  tables, plots, session info, rendered report.
\item
  \textbf{Algorithm/pseudocode:} validate schemas → filter by declared
  rules → fit declared models → contrasts/Holm → diagnostics → effect
  sizes/CIs.
\item
  \textbf{Implementation notes:} scripts read, never edit, raw data;
  generate every table/plot programmatically.
\item
  \textbf{Test criteria:} synthetic data recover known effects; rerun
  hashes match; model convergence/diagnostic status is printed.
\item
  \textbf{Expected result:} one command regenerates manuscript-ready
  results after data lock.
\item
  \textbf{How to run:} \texttt{python\ analysis/prepare\_data.py} then
  \texttt{Rscript\ analysis/technical\_models.R} and
  \texttt{Rscript\ analysis/hci\_models.R}.
\end{itemize}

\subsection{Task 10 --- Reproducibility and
release}\label{task-10-reproducibility-and-release}

\begin{itemize}
\tightlist
\item
  \textbf{Purpose:} package code, metadata, de-identified data, and
  documented limitations.
\item
  \textbf{Files to create:} \texttt{README.md}, \texttt{CITATION.cff},
  \texttt{LICENSE}, \texttt{CHANGELOG.md}, \texttt{environment/},
  \texttt{docs/data\_dictionary.md}, \texttt{docs/reproduction.md}, CI
  workflow, archival manifest.
\item
  \textbf{Files to modify:} all dependency locks and configuration
  hashes at release.
\item
  \textbf{Technology:} Git, GitHub/GitLab, OSF/Zenodo or institutional
  repository, optional Docker.
\item
  \textbf{Inputs/outputs:} frozen repository + artifacts → versioned
  research release and DOI where available.
\item
  \textbf{Algorithm/pseudocode:} clean checkout → install locked
  dependencies → smoke simulation → tests → miniature batch → analysis →
  compare expected hashes.
\item
  \textbf{Implementation notes:} license only material owned by the
  project; do not redistribute restricted trajectory data or participant
  identifiers.
\item
  \textbf{Test criteria:} an uninvolved collaborator reproduces the
  miniature result from the guide; archive manifest verifies.
\item
  \textbf{Expected result:} reviewers can audit the full
  decision-to-explanation chain.
\item
  \textbf{How to run:} \texttt{make\ reproduce-mini} or its documented
  platform-neutral command sequence.
\end{itemize}

\section{25. File/Folder Structure}\label{filefolder-structure}

\begin{verbatim}
intersection-priority/
├── README.md
├── CITATION.cff
├── LICENSE
├── pyproject.toml
├── uv.lock                         # or another single chosen Python lock
├── Makefile
├── configs/
│   ├── intersection.yaml
│   ├── controllers/
│   │   ├── fixed_time.yaml
│   │   ├── efficiency.yaml
│   │   ├── equal_bargaining.yaml
│   │   └── efpb.yaml
│   ├── scenarios/
│   │   ├── factorial.yaml
│   │   └── stress.yaml
│   ├── experiments/
│   │   ├── smoke.yaml
│   │   ├── pilot.yaml
│   │   └── confirmatory.yaml
│   └── study/
│       ├── pilot.yaml
│       └── confirmatory.yaml
├── sumo/
│   ├── net/intersection.net.xml
│   ├── tls/base_tls.add.xml
│   ├── detectors/detectors.add.xml
│   ├── routes/.gitkeep
│   └── base.sumocfg
├── src/intersection_priority/
│   ├── domain/models.py
│   ├── sensing/{traci_adapter,state_collector,normalization,noise}.py
│   ├── agents/{pedestrian,vehicle}.py
│   ├── prediction/{discharge,calibration}.py
│   ├── fairness/{weights,metrics}.py
│   ├── controllers/{base,fixed_time,efficiency,equal_bargaining,efpb,signal_safety}.py
│   ├── explanation/{trace,generator,templates,counterfactual,validators}.py
│   ├── simulation/{runner,scenarios,logging}.py
│   └── api/{main,schemas,routes,websocket}.py
├── dashboard/
│   ├── package.json
│   ├── package-lock.json
│   └── src/
│       ├── components/{IntersectionView,StatePanel,DecisionCard,ReasonCard,History}.tsx
│       ├── lib/{api,schemas,format}.ts
│       └── {App,main}.tsx
├── experiments/{generate_manifests,run_batch}.py
├── study/
│   ├── protocol.md
│   ├── build_stimuli.py
│   ├── scoring.py
│   ├── stimuli/manifest.json
│   └── app/
├── analysis/
│   ├── prepare_data.py
│   ├── technical_metrics.py
│   ├── technical_models.R
│   ├── hci_models.R
│   ├── figures.R
│   └── report.qmd
├── data/
│   ├── raw/.gitkeep               # gitignored; restricted datasets/participant exports
│   ├── manifests/
│   ├── runs/.gitkeep              # gitignored large run outputs
│   ├── processed/
│   └── release/
├── docs/{requirements,data_dictionary,reproduction}.md
└── tests/
    ├── unit/
    ├── integration/
    ├── regression/
    ├── property/
    └── fixtures/
\end{verbatim}

\section{26. APIs/Data Structures}\label{apisdata-structures}

\subsection{26.1 Core data structures}\label{core-data-structures}

{\def\LTcaptype{none} % do not increment counter
\begin{longtable}[]{@{}
  >{\raggedright\arraybackslash}p{(\linewidth - 2\tabcolsep) * \real{0.5000}}
  >{\raggedright\arraybackslash}p{(\linewidth - 2\tabcolsep) * \real{0.5000}}@{}}
\toprule\noalign{}
\begin{minipage}[b]{\linewidth}\raggedright
Type
\end{minipage} & \begin{minipage}[b]{\linewidth}\raggedright
Required fields
\end{minipage} \\
\midrule\noalign{}
\endhead
\bottomrule\noalign{}
\endlastfoot
\texttt{Measurement{[}T{]}} & \texttt{value:T}, \texttt{unit:str},
\texttt{observed\_at:float}, \texttt{age\_s:float},
\texttt{quality:float}, \texttt{source\_id:str},
\texttt{status:OK\textbackslash{}\textbar{}MISSING\textbackslash{}\textbar{}STALE\textbackslash{}\textbar{}DEGRADED} \\
\texttt{PedestrianState} & counts/arrival rates/mean and max
waits/density by crosswalk; crossing occupancy; accessibility extension;
service history \\
\texttt{VehicleState} & queues/arrival rates/mean and max
waits/cumulative delay/downstream occupancy by approach; emergency
request \\
\texttt{SignalState} & current SUMO phase, semantic service phase,
elapsed/remaining time, next admissible transition, SPaT timestamp \\
\texttt{IntersectionState} & \texttt{state\_id}, \texttt{sim\_time\_s},
pedestrian, vehicle, signal, sensor health, active overrides, schema
version \\
\texttt{PlanStep} & semantic phase, duration, transition path, served
movements \\
\texttt{CandidatePlan} & \texttt{candidate\_id}, ordered steps, horizon,
source/config hash \\
\texttt{CandidateEvaluation} & utilities/disagreement gains, mode
weights, predicted service/delay/tail wait, efficiency loss, feasibility
predicates, Nash score, reason facts \\
\texttt{PriorityDecision} & \texttt{decision\_id}, chosen
candidate/target phase, decision kind, alternatives, binding
constraints, created/effective times, controller/version/config hash \\
\texttt{DecisionTrace} & input state hash/snapshot, every candidate
evaluation, tie-break path, execution receipt, provenance \\
\texttt{FairnessResult} & service ratios, Jain (nullable), normalized
mean/tail gaps, max waits, violations, units/reference values \\
\texttt{Counterfactual} & target action, feature deltas, distance,
search bounds, replay decision/hash, valid/minimality status \\
\texttt{ExplanationRecord} & canonical text/slots, facts, constraints,
alternative, counterfactual, template/trace hash, validation status \\
\texttt{ScenarioConfig} & demand processes/rates, routes, walking
speeds, events, noise, durations, seed policy \\
\texttt{ScenarioManifest} & sampled immutable arrivals/events, generator
version, seed-stream IDs, hash \\
\texttt{RunMetadata} & run/scenario/controller/seed IDs, resolved
config, software/environment versions, start/end/status/failure \\
\texttt{LiveFrame} & compact state, current signal, latest
decision/explanation, history cursor; no raw candidate table unless
requested \\
\texttt{StudyTrial} & stimulus/decision IDs, condition, quality, correct
answers, order/set, asset hashes \\
\texttt{StudyResponse} & pseudonymous participant/session/trial IDs,
answers, confidence/ratings, monotonic timings, focus/render flags,
submitted timestamp \\
\end{longtable}
}

Represent physical units in field names or metadata and validate bounds.
Store nullable/unknown explicitly; never encode missing sensor data as
zero. Version every externally stored schema.

\subsection{26.2 REST and WebSocket API}\label{rest-and-websocket-api}

{\def\LTcaptype{none} % do not increment counter
\begin{longtable}[]{@{}
  >{\raggedright\arraybackslash}p{(\linewidth - 6\tabcolsep) * \real{0.300}}
  >{\raggedright\arraybackslash}p{(\linewidth - 6\tabcolsep) * \real{0.190}}
  >{\raggedright\arraybackslash}p{(\linewidth - 6\tabcolsep) * \real{0.200}}
  >{\raggedright\arraybackslash}p{(\linewidth - 6\tabcolsep) * \real{0.310}}@{}}
\toprule\noalign{}
\begin{minipage}[b]{\linewidth}\raggedright
Method/path
\end{minipage} & \begin{minipage}[b]{\linewidth}\raggedright
Function
\end{minipage} & \begin{minipage}[b]{\linewidth}\raggedright
Input
\end{minipage} & \begin{minipage}[b]{\linewidth}\raggedright
Output/notes
\end{minipage} \\
\midrule\noalign{}
\endhead
\bottomrule\noalign{}
\endlastfoot
\texttt{GET} \path|/api/v1/health| & process/SUMO status & none & versions,
active run, degraded components \\
\texttt{GET} \path|/api/v1/state| & latest intersection state & optional
\texttt{run\_id} & \texttt{IntersectionState}; \texttt{404} if none \\
\texttt{GET} \path|/api/v1/decisions/{decision_id}| & decision and compact
candidate summary & path ID & \texttt{PriorityDecision}; details require
audit route \\
\texttt{GET} \path|/api/v1/decisions/{decision_id}/trace| & audit trace &
path ID & operator/research authorization; immutable trace \\
\texttt{GET} \path|/api/v1/decisions/{decision_id}/explanation| & validated
explanation & \(condition=D\|F\|FC\), audience &
\texttt{ExplanationRecord}/abstention \\
\texttt{POST} \path|/api/v1/simulations| & start configured run &
scenario/controller IDs, seed & \texttt{202}, run ID; only whitelisted
configs \\
\texttt{POST} \path|/api/v1/simulations/{run_id}/stop| & graceful stop &
reason & closes/logs; no arbitrary process kill \\
\texttt{GET} \path|/api/v1/simulations/{run_id}| & run status & ID &
metadata, progress, validation status \\
\texttt{POST} \path|/api/v1/replays| & start immutable replay & stimulus/run
ID & replay session/cursor \\
\texttt{GET} \path|/api/v1/scenarios| & list approved scenarios & filters &
IDs and metadata, not writable paths \\
\texttt{POST} \path|/api/v1/study/responses| & store a study response if
self-hosted & \texttt{StudyResponse} & idempotent receipt; reject
duplicate trial key \\
\texttt{WS} \path|/api/v1/live?run_id=...| & stream live/replay frames & auth/session
& schema-versioned frames; bounded rate/backpressure \\
\end{longtable}
}

The API accepts semantic phase/config identifiers, never shell commands,
file paths, or raw signal bit strings. Validate all input, bind
development services to localhost, and add study authentication/HTTPS
before remote recruitment.

\subsection{26.3 Storage decision}\label{storage-decision}

A relational database is unnecessary for the simulation core. Runs are
immutable, append-heavy, and naturally partitioned:

\begin{verbatim}
data/runs/{scenario_id}/{controller_id}/{seed}/{run_id}/
  metadata.json
  resolved_config.yaml
  manifest.json
  states.parquet
  decisions.jsonl
  explanations.jsonl
  events.jsonl
  sumo_tripinfo.xml
  sumo_personinfo.xml
  validation_report.json
  run_metrics.parquet
\end{verbatim}

Parquet supports typed columnar analysis; JSONL preserves nested trace
events and partial-run recovery. Write append batches and atomically
finalize. A self-hosted remote HCI study benefits from a small encrypted
PostgreSQL/SQLite-backed response service for idempotent transactions,
but an institutionally approved survey platform is preferable. Export
de-identified responses to a frozen analytical table. Keep
consent/contact/payment mappings outside the repository and outside
decision logs.

\section{27. Dataset Strategy}\label{dataset-strategy}

The controller does not need a training dataset. Synthetic SUMO demand
is sufficient for causal controller comparisons, but a field dataset
materially improves parameter plausibility and scenario validity. The
recommended design is \textbf{Option C-lite: synthetic experimental
control calibrated with a public signalized-intersection dataset}.

{\def\LTcaptype{none} % do not increment counter
\begin{longtable}[]{@{}
  >{\raggedright\arraybackslash}p{(\linewidth - 6\tabcolsep) * \real{0.2500}}
  >{\raggedright\arraybackslash}p{(\linewidth - 6\tabcolsep) * \real{0.2500}}
  >{\raggedright\arraybackslash}p{(\linewidth - 6\tabcolsep) * \real{0.2500}}
  >{\raggedright\arraybackslash}p{(\linewidth - 6\tabcolsep) * \real{0.2500}}@{}}
\toprule\noalign{}
\begin{minipage}[b]{\linewidth}\raggedright
Option
\end{minipage} & \begin{minipage}[b]{\linewidth}\raggedright
Strength
\end{minipage} & \begin{minipage}[b]{\linewidth}\raggedright
Weakness
\end{minipage} & \begin{minipage}[b]{\linewidth}\raggedright
Decision
\end{minipage} \\
\midrule\noalign{}
\endhead
\bottomrule\noalign{}
\endlastfoot
A. Fully synthetic SUMO & Complete control, easy factorial scenarios,
reproducible & Weak realism and external validity; arbitrary
arrivals/walking speeds & Acceptable minimum for proof of concept, with
narrow claims \\
B. Replay public trajectories directly & Natural interactions and
distributions & Signal policy cannot be changed causally; missing
SPaT/queues in many datasets & Use for descriptive validation, not
controller comparison \\
C. Calibrate SUMO from field data & Retains causal controller
experiments while grounding demand/behavior & Mapping/calibration work
and license/access constraints & \textbf{Recommended} \\
\end{longtable}
}

\begin{landscape}
{\def\LTcaptype{none} % do not increment counter
\footnotesize
\setlength{\tabcolsep}{4pt}
\begin{longtable}[]{@{}
  >{\raggedright\arraybackslash}p{(\linewidth - 10\tabcolsep) * \real{0.100}}
  >{\centering\arraybackslash}p{(\linewidth - 10\tabcolsep) * \real{0.140}}
  >{\centering\arraybackslash}p{(\linewidth - 10\tabcolsep) * \real{0.120}}
  >{\raggedright\arraybackslash}p{(\linewidth - 10\tabcolsep) * \real{0.250}}
  >{\raggedright\arraybackslash}p{(\linewidth - 10\tabcolsep) * \real{0.250}}
  >{\raggedright\arraybackslash}p{(\linewidth - 10\tabcolsep) * \real{0.140}}@{}}
\toprule\noalign{}
\begin{minipage}[b]{\linewidth}\raggedright
Dataset
\end{minipage} & \begin{minipage}[b]{\linewidth}\centering
Pedestrians + vehicles
\end{minipage} & \begin{minipage}[b]{\linewidth}\centering
Signal status
\end{minipage} & \begin{minipage}[b]{\linewidth}\raggedright
Principal format
\end{minipage} & \begin{minipage}[b]{\linewidth}\raggedright
Access/license status to verify
\end{minipage} & \begin{minipage}[b]{\linewidth}\raggedright
Realistic role here
\end{minipage} \\
\midrule\noalign{}
\endhead
\bottomrule\noalign{}
\endlastfoot
IMPTC & Y, including multiple VRU classes & Y, 1 Hz signal groups &
Repository-defined synchronized trajectory/context files; JSON signal
context; samples and schemas provided & Data described as non-commercial
research; repository code is Apache-2.0---do not assume the code license
licenses all data & Primary behavioral/demand calibration \\
SinD v1 & Y, plus non-motorized modes & Y & CSV trajectories + Lanelet2
HD map and metadata & Public samples; full data requested by educational
email; confirm terms on receipt & Geographic/behavioral sensitivity \\
U.S. DOT ISC Stage 1B & Y, plus bicycles & Y, SPaT and calls &
Multi-sensor packages, calibration/metadata, 3D ground truth, run labels
& U.S. public domain; sample/full validation downloads & Sensor and
controlled safety robustness \\
inD & Y, plus bicycles & Not consistently suitable for required timing &
CSV trajectory/recording metadata & Non-commercial access terms must be
confirmed & Secondary speed/interaction sanity check \\
\end{longtable}
}
\end{landscape}

\subsection{IMPTC --- primary
recommendation}\label{imptc-primary-recommendation}

The official IMPTC repository describes an instrumented public German
intersection with synchronized camera/LiDAR data, 25 Hz trajectories,
weather/context, and 1 Hz signal status for vehicle and pedestrian
signal groups. It reports eight hours, more than 2,500
vulnerable-road-user and 20,000 vehicle trajectories overall; the public
sequence-focused release contains 270 sequences and explicitly includes
pedestrian, cyclist, wheelchair, stroller, and vehicle tracks. The data
are offered for non-commercial research and the repository provides
samples, formats, and download instructions.\hyperref[source-33]{[33]}

Use: derive walking-speed distributions, pedestrian/vehicle arrival
processes by time block, queue-discharge/saturation estimates,
signal-group mappings, and sensor-error plausibility. Map the chosen
subset to the simplified SUMO geometry; do not claim exact digital-twin
calibration. Record dataset version, selected sequence IDs,
transformation code, and license. Because access/license details can
change, archive the terms accepted at download time without
redistributing restricted files.

\subsection{SinD --- strong alternative/sensitivity
source}\label{sind-strong-alternativesensitivity-source}

SinD v1 is a drone trajectory dataset at a signalized intersection in
China with diverse vehicles and pedestrians, traffic-light states, HD
map, and CSV trajectories; the ITSC paper reports seven hours and more
than 13,000 traffic participants. Full access is requested with an
educational email; public samples and format documentation are in the
official repository.\hyperref[source-34]{[34]} A later SinD 2.0 release
adds multi-city data, but its 2026 core paper is currently described as
an arXiv preprint, so use the peer-reviewed v1 citation for the main
study and treat v2 as optional external
validation.\hyperref[source-35]{[35]}

Use: test whether parameters calibrated in one region change conclusions
under heterogeneous walking, non-motorized traffic, and
signal-compliance patterns. Limitation: geographic behavior and geometry
differ, and full access is not guaranteed.

\subsection{U.S. DOT Intersection Safety Challenge --- optional
robustness
data}\label{u.s.-dot-intersection-safety-challenge-optional-robustness-data}

The official Stage 1B dataset covers a controlled signalized four-way
smart intersection with 20 roadside sensors, SPaT, vehicle/pedestrian
calls, ground-truth 3D boxes, and conflict labels. The current
validation release reports 41 runs and about 30 GB and is U.S.
public-domain data.\hyperref[source-36]{[36]} Use the small sample first
for sensor schema/failure and safety-event validation. It is controlled
conflict data rather than natural demand, so it should not calibrate
ordinary arrival rates.

\subsection{inD --- secondary only}\label{ind-secondary-only}

The inD dataset contains naturalistic vehicle, bicycle, and pedestrian
trajectories at German intersections.\hyperref[source-37]{[37]} It is
useful for walking/driving distributions and interaction sanity checks,
but not the primary timing dataset if SPaT is unavailable for the
required sites. Confirm the current non-commercial license and access
terms before use.

Minimum-data path: use only official IMPTC samples to validate the
transform, then publish synthetic calibrated distributions---not raw
restricted data. If no full dataset is obtained, state
``parameter-informed by public samples/literature'' and bound claims to
simulation.

\section{28. Reproducibility Plan}\label{reproducibility-plan}

\begin{itemize}
\tightlist
\item
  Preregister primary controller endpoints, HCI hypotheses/contrasts,
  exclusion rules, stopping rules, parameter selection, and deviations.
\item
  Freeze a calibration split and never tune on confirmatory scenario
  seeds or participant data.
\item
  Store the exact SUMO binary/package version, Python/R/Node versions,
  OS/CPU, dependency locks, locale/time zone, and Git commit.
\item
  Derive independent RNG streams from a master seed with named keys
  (\texttt{vehicle\_arrivals}, \texttt{ped\_arrivals}, \texttt{speeds},
  \texttt{emergency}, \texttt{sensor\_noise}); store generated
  manifests.
\item
  Use paired seeds and randomized controller execution order; set one
  process/thread policy when nondeterministic libraries are involved.
\item
  Resolve all YAML inheritance before a run and store the resolved
  configuration plus SHA-256 hash.
\item
  Make raw run and study exports read-only; transformations write new
  versioned files and retain a data lineage table.
\item
  Publish schemas, data dictionary, scenario manifests, controller
  configs, de-identified run-level data, validated decision
  traces/stimuli, scoring keys after study completion, and analysis
  code.
\item
  Do not publish compensation/identity mappings, unrestricted licensed
  datasets, secrets, or potentially re-identifying free text without
  review.
\item
  Use continuous integration for unit/integration/golden/schema tests
  and a miniature deterministic experiment.
\item
  Archive a tagged release and checksum manifest in a durable
  repository; document how restricted data can be requested.
\end{itemize}

For stochastic precision, begin with 30 paired seeds and use the
preregistered CI half-width rule described in Section 10. Compute
run-level paired differences, then a percentile or BCa bootstrap that
resamples seed blocks within scenario. Publish both the stopping
trajectory and final estimate so sequential checking is transparent.

For the HCI study, export the randomization schedule and stimulus hashes
before recruitment. Version questionnaire wording and translations. Save
analysis package/session information, model formulas, convergence
diagnostics, and a machine-readable result table. An independent analyst
should reproduce primary tables from the frozen de-identified data
without opening the dashboard.

\section{29. Step-by-Step Execution
Guide}\label{step-by-step-execution-guide}

These commands are the intended interface for the proposed repository;
they become executable after the roadmap is implemented. Pin the exact
releases in lock files at project start. As of the research date, the
official SUMO download page lists release 1.27.1 and documents
pip-distributed tools; confirm again when implementation
begins.\hyperref[source-32]{[32]}

\subsection{29.1 Install}\label{install}

\begin{verbatim}
git clone <repository-url> intersection-priority
cd intersection-priority
python3.12 -m venv .venv
source .venv/bin/activate
python -m pip install --upgrade pip
python -m pip install -e '.[dev,sumo]'
npm --prefix dashboard ci
Rscript -e 'if (!requireNamespace("renv", quietly=TRUE)) install.packages("renv")'
Rscript -e 'renv::restore()'
\end{verbatim}

If SUMO is installed via its official pip distribution rather than the
project extra:

\begin{verbatim}
python -m pip install 'eclipse-sumo==1.27.1'
sumo --version
python -c 'import traci, sumolib; print("TraCI and sumolib import OK")'
\end{verbatim}

On a system package installation, set \texttt{SUMO\_HOME} to that
installation's tools directory as documented by SUMO. Do not hard-code a
developer's home path.

\subsection{29.2 Verify the environment and
network}\label{verify-the-environment-and-network}

\begin{verbatim}
ruff check .
mypy src
pytest -q
sumo -c sumo/base.sumocfg --no-step-log true --duration-log.statistics true
sumo-gui -c sumo/base.sumocfg
\end{verbatim}

Manually inspect every phase and crossing once, then save the signed
verification checklist. A green test suite does not replace geometry
inspection.

\subsection{29.3 Generate immutable demand
manifests}\label{generate-immutable-demand-manifests}

\begin{verbatim}
python experiments/generate_manifests.py \
  --config configs/experiments/pilot.yaml \
  --output data/manifests/pilot
python -m intersection_priority.simulation.scenarios validate \
  data/manifests/pilot
\end{verbatim}

The output is a manifest per scenario/seed plus SUMO route/person files
and hashes. Generate once; all controllers read the same files.

\subsection{29.4 Smoke and calibrate}\label{smoke-and-calibrate}

\begin{verbatim}
python experiments/run_batch.py \
  --config configs/experiments/smoke.yaml \
  --controllers fixed_time efficiency equal_bargaining efpb \
  --workers 1
python analysis/prepare_data.py --experiment smoke
\end{verbatim}

Then run calibration-only sweeps for discharge parameters, \(\delta\), debt,
horizon, and normalization. Freeze the selected configuration and
rationale before confirmatory runs.

\subsection{29.5 Run confirmatory
simulation}\label{run-confirmatory-simulation}

\begin{verbatim}
python experiments/run_batch.py \
  --config configs/experiments/confirmatory.yaml \
  --manifest-dir data/manifests/confirmatory \
  --workers 4
python analysis/prepare_data.py --experiment confirmatory --strict
Rscript analysis/technical_models.R confirmatory
\end{verbatim}

Inspect \texttt{validation\_report.json} and the invalid-run table
before any inferential output. Do not simply rerun failed seeds until a
favorable result appears; diagnose, document, and apply a uniform rule.

\subsection{29.6 Run dashboard and create study
stimuli}\label{run-dashboard-and-create-study-stimuli}

\begin{verbatim}
uvicorn intersection_priority.api.main:app --host 127.0.0.1 --port 8000
npm --prefix dashboard run dev
python study/build_stimuli.py \
  --config configs/study/pilot.yaml \
  --source data/runs/confirmatory \
  --output study/stimuli/pilot
pytest -q tests/regression/test_stimuli.py
\end{verbatim}

Conduct cognitive interviews and the instrument pilot. Revise only using
pilot evidence, then generate and hash the confirmatory stimulus set.
Obtain ethics approval before recruiting participants.

\subsection{29.7 Collect and analyze the HCI
study}\label{collect-and-analyze-the-hci-study}

\begin{verbatim}
python study/scoring.py \
  --responses data/raw/study_export.csv \
  --stimuli study/stimuli/confirmatory/manifest.json \
  --output data/processed/study_trials.parquet
python analysis/prepare_data.py --study confirmatory --strict
Rscript analysis/hci_models.R confirmatory
quarto render analysis/report.qmd
\end{verbatim}

The end-to-end flow is:

\begin{verbatim}
requirements/standards
→ network and safe phase machine
→ field-informed calibration
→ immutable stochastic manifests
→ paired baseline/EFPB simulations
→ trace and explanation validation
→ frozen matched study replays
→ ethics-approved HCI experiment
→ locked preprocessing
→ mixed-model technical/HCI analyses
→ sensitivity/ablation results
→ de-identified reproducibility release
\end{verbatim}

\section{30. Expected Results}\label{expected-results}

No empirical results exist yet. The following are
\textbf{evidence-informed expectations and decision criteria}, not
findings or numeric forecasts.

{\def\LTcaptype{none} % do not increment counter
\begin{longtable}[]{@{}
  >{\raggedright\arraybackslash}p{(\linewidth - 4\tabcolsep) * \real{0.3333}}
  >{\raggedright\arraybackslash}p{(\linewidth - 4\tabcolsep) * \real{0.3333}}
  >{\raggedright\arraybackslash}p{(\linewidth - 4\tabcolsep) * \real{0.3333}}@{}}
\toprule\noalign{}
\begin{minipage}[b]{\linewidth}\raggedright
Claim under test
\end{minipage} & \begin{minipage}[b]{\linewidth}\raggedright
Pattern that would support it
\end{minipage} & \begin{minipage}[b]{\linewidth}\raggedright
Pattern that would weaken/refute it
\end{minipage} \\
\midrule\noalign{}
\endhead
\bottomrule\noalign{}
\endlastfoot
Bounded-efficiency fairness & EFPB reduces tail disparity/starvation
versus B1 while realized delay usually remains within the declared bound
& Fairness gains appear only at large delay cost, or realized loss
repeatedly exceeds \(\delta\) because prediction is inaccurate \\
Debt mechanism & EFPB recovers the underserviced mode faster than static
weighting in S8 & Equal weights perform the same, indicating needless
complexity \\
Bargaining value & EFPB yields better Pareto/tail outcomes or clearer
stable policy than a tuned weighted sum & Weighted sum matches all
outcomes and is easier to explain \\
Safety/fallback & Zero conflict/clearance violations and honest
degraded-mode abstention under S7/S11/S12 & Emergency truncates
clearance, stale data are treated as zero, or explanations conceal
fallback \\
Factual explanation & F and FC improve reason-question accuracy over D &
Only subjective clarity rises; objective accuracy is unchanged or
worse \\
Counterfactual & FC improves decision-flip accuracy enough to justify
added time & FC slows users without improving contrastive accuracy or
increases confusion \\
Calibrated trust & Explanations increase the sound--anomalous
discrimination gap and lower Brier score & Trust rises equally for bad
and good decisions, or anomalous decisions become more persuasive \\
Perceived fairness & Sound decisions with faithful reasons receive
higher procedural/informational justice, while faulty decisions remain
penalized & Explanation merely beautifies bad outcomes or distributive
ratings ignore traffic outcomes \\
\end{longtable}
}

Likely operational behavior is a trade-off: B1 may have the lowest mean
delay in some saturated scenarios, while EFPB should improve
maximum/tail waits and reduce long service gaps. It is plausible that
factual reasons help rationale comprehension and counterfactuals help
contrastive questions, but explanation detail may increase response
time. Literature on transparency warns against assuming monotonic
benefits; too much information can reduce trust or
usability.\hyperref[source-20]{[20]}\,\hyperref[source-38]{[38]}

Report null results as informative. If EFPB offers no benefit over a
simpler controller but faithful explanations help calibration, the paper
becomes primarily an HCI/XAI contribution. If explanations do not help
but the bounded controller is strong, publish the transportation method
and treat the interface study as negative evidence. Do not force both
halves into positive claims.

\section{31. Threats to Validity}\label{threats-to-validity}

{\def\LTcaptype{none} % do not increment counter
\begin{longtable}[]{@{}
  >{\raggedright\arraybackslash}p{(\linewidth - 4\tabcolsep) * \real{0.3333}}
  >{\raggedright\arraybackslash}p{(\linewidth - 4\tabcolsep) * \real{0.3333}}
  >{\raggedright\arraybackslash}p{(\linewidth - 4\tabcolsep) * \real{0.3333}}@{}}
\toprule\noalign{}
\begin{minipage}[b]{\linewidth}\raggedright
Validity threat
\end{minipage} & \begin{minipage}[b]{\linewidth}\raggedright
Consequence
\end{minipage} & \begin{minipage}[b]{\linewidth}\raggedright
Mitigation/evidence
\end{minipage} \\
\midrule\noalign{}
\endhead
\bottomrule\noalign{}
\endlastfoot
Utility/weight arbitrariness & Apparent fairness is encoded by
convenient coefficients & Requirements provenance, calibration-only
selection, equal-weight primary case, Pareto/sensitivity analysis,
stakeholder review \\
Horizon predictor error & Declared efficiency bound fails in realized
SUMO outcomes & Held-out rollout validation, realized-bound reporting,
uncertainty margin, shorter replanning \\
Simplified no-turn geometry & Findings miss common vehicle--pedestrian
conflicts & State scope clearly; add protected-turn sensitivity or
second network before strong ITS claims \\
All-pedestrian phase & Results may not transfer to concurrent pedestrian
phasing & Treat as controlled archetype; replicate one
concurrent-crossing configuration later \\
Simulator behavior & SUMO car-following/walking assumptions determine
outcomes & Calibrate from IMPTC/SinD, vary behavior parameters, report
simulator/version, avoid field-safety claims \\
Real-data mapping & Field distribution is distorted by simplified
geometry & Publish transformation and goodness-of-fit; use data for
ranges, not false digital-twin language \\
Sensor perfection & Controller depends on unavailable state & Sensor
abstraction, noise/staleness/fallback scenarios, separate ground truth
from observed state \\
Emergency idealization & Unrealistic response/security claims & Simulate
authenticated event only; exclude detection/V2X security from claimed
contribution \\
Baseline unfairness & Proposal wins against weak implementations &
Shared constraints/action budget; calibration parity; closest bargaining
and strong adaptive baseline \\
Objective-fairness construct & Low disparity may still be unjust or
universally bad & Multiple mode/tail/service metrics, efficiency bound,
stakeholder review, never use Jain alone \\
Explanation confound & Layout/detail/time rather than reasons causes
effects & Identical renderer/replays, placeholder card, matched sets,
condition-to-set rotation \\
Demand/decision confound & Explainable condition happens to contain
easier cases & Within-person matched scenario sets; multilevel model
with scenario effects \\
Practice/order effects & Later conditions score better &
Williams/Latin-square counterbalancing, practice before scored trials,
order covariate \\
Self-report/common-method bias & Trust/fairness rise without better
decisions & Objective questions, response time, confidence/Brier score,
sound/anomalous trials \\
Scale adaptation & Organizational justice items may not measure
intersection fairness & Cognitive interviews, expert review, pilot,
omega/factor analysis, narrow labeling \\
Participant sample & Online/student sample misses older, disabled, or
low-literacy users & Quotas/outreach, accessibility, demographics,
subgroup estimates only if powered \\
Deception/demand effects & Anomalies or study purpose bias responses &
Ethics review, credible cover wording, physically safe faults, debrief,
suspicion check \\
Multiple testing & Selective false positives & Preregistered
primaries/contrasts, Holm correction, full result table \\
Simulation pseudoreplication & Artificially tiny p-values & Run is unit;
paired seeds/random effects; no 1 Hz observation inflation \\
Researcher degrees of freedom & Tuning and exclusions chase results &
Frozen configs, preregistration, immutable manifests, deviation log \\
\end{longtable}
}

\section{32. Research Limitations}\label{research-limitations}

\begin{itemize}
\tightlist
\item
  The core experiment is one isolated intersection, not a coordinated
  urban network; network spillback and corridor progression remain
  future work.
\item
  Virtual mode advocates aggregate heterogeneous people. They cannot
  express individual disability, trip purpose, passenger occupancy,
  economic disadvantage, or emergency need unless those data are
  explicitly and ethically modeled.
\item
  Weighted Nash bargaining is a normative allocation convention. It does
  not prove that the selected outcome is socially just or that road
  users endorse its utilities.
\item
  Maximum-wait limits and base weights depend on policy and
  jurisdiction. A simulation cannot legitimate them; stakeholder and
  standards processes are needed.
\item
  The horizon model approximates SUMO discharge. The formal
  near-efficiency bound applies to predicted loss; realized performance
  must be separately reported.
\item
  The all-pedestrian phase simplifies conflict logic and may overstate
  pedestrian interruption cost compared with well-designed concurrent
  phases.
\item
  Public trajectory datasets were collected in specific countries,
  intersections, times, and sensor systems. Calibration improves
  plausibility but not universal representativeness.
\item
  The dashboard is an audit/replay interface, not a roadside eHMI and
  not a driver-facing display. Its findings do not establish safe glance
  behavior while driving or crossing.
\item
  Short laboratory/online exposure measures initial comprehension and
  trust, not long-term adaptation, gaming, habituation, or community
  legitimacy.
\item
  Adapted justice measures require local validation; perceived fairness
  can be influenced by wording even when outcomes are inequitable.
\item
  Intentional anomaly trials test calibration under controlled faults,
  not the full range of real controller failures.
\item
  No field hardware, certified controller, cybersecurity,
  privacy-preserving sensing, accessibility certification, or regulatory
  approval is delivered.
\end{itemize}

These limitations are compatible with a strong paper if claims remain at
the level of a controlled, reproducible proof of concept.

\section{33. Publication Potential}\label{publication-potential}

The project has credible publication potential, but the appropriate
paper depends on which evidence becomes strongest.

{\def\LTcaptype{none} % do not increment counter
\begin{longtable}[]{@{}
  >{\raggedright\arraybackslash}p{(\linewidth - 6\tabcolsep) * \real{0.2500}}
  >{\raggedright\arraybackslash}p{(\linewidth - 6\tabcolsep) * \real{0.2500}}
  >{\raggedright\arraybackslash}p{(\linewidth - 6\tabcolsep) * \real{0.2500}}
  >{\raggedright\arraybackslash}p{(\linewidth - 6\tabcolsep) * \real{0.2500}}@{}}
\toprule\noalign{}
\begin{minipage}[b]{\linewidth}\raggedright
Venue type
\end{minipage} & \begin{minipage}[b]{\linewidth}\raggedright
Fit
\end{minipage} & \begin{minipage}[b]{\linewidth}\raggedright
Evidence reviewers will expect
\end{minipage} & \begin{minipage}[b]{\linewidth}\raggedright
Main risk
\end{minipage} \\
\midrule\noalign{}
\endhead
\bottomrule\noalign{}
\endlastfoot
\emph{Transportation Research Part C} & Integrated emerging-technology
control with transportation implications; its scope explicitly includes
ITS, real-time operations, multimodal transport, pedestrians, and
open-science data.\hyperref[source-39]{[39]} & Strong baselines,
calibration, sensitivity, realized efficiency/fairness trade-offs,
preferably a second geometry & Single-junction/HCI prototype may look
underpowered as a control contribution \\
\emph{IEEE Transactions on Intelligent Transportation Systems} &
Sensing/control/simulation involving multimodal transport and road users
is explicitly in scope.\hyperref[source-40]{[40]} & Technical rigor,
robust/failure tests, scalability or broader network, reproducibility &
Algorithm may appear incremental after Wang et al.~and existing fair
control \\
\emph{Transportation Research Part F} & Behavioral and psychological
aspects of traffic and transport.\hyperref[source-41]{[41]} & Strong
theory, powered HCI study, validated/adapted measures, road-user
implications & Detailed control algorithm may distract; use as stimulus
infrastructure \\
ACM AutomotiveUI & Human-centered mobility interfaces explicitly include
pedestrians and drivers; archival papers emphasize scientific excellence
and open materials.\hyperref[source-42]{[42]} & Rigorous user study,
accessibility, ecological interface rationale, supplemental
video/materials & An operator/audit dashboard needs clear
automotive/mobility interaction relevance \\
ACM IUI & Human-centered AI, XAI, trust, and interactive intelligent
systems are explicit topics.\hyperref[source-43]{[43]} & Generalizable
explanation contract, objective comprehension/calibration, strong
mixed-model analysis & Traffic contribution alone is insufficient;
explanation insight must generalize \\
ACM CHI or ACM TiiS & Human-centered XAI and sociotechnical fairness &
Multiple studies or unusually strong theory/generalization, open
instrument/stimuli & A single domain prototype and modest sample are
likely too narrow \\
\end{longtable}
}

Most realistic publication paths:

\begin{enumerate}
\def\labelenumi{\arabic{enumi}.}
\tightlist
\item
  \textbf{HCI-first:} submit a rigorous three-condition
  explanation/calibration paper to IUI, AutomotiveUI, or TR-F. EFPB is a
  transparent, high-fidelity stimulus generator; the novelty centers on
  the trace contract and calibrated judgment.
\item
  \textbf{ITS-first:} expand to a second topology or small corridor, add
  calibrated field distributions and stronger closest baselines, then
  target T-ITS or TR-C. The HCI study supports deployability but is not
  the sole novelty.
\item
  \textbf{Two-paper path:} publish the bounded controller/benchmark and
  the human explanation study separately only if each has enough
  independent depth and no duplicated contribution/data claim.
\end{enumerate}

What makes it publishable: a precise gap against Wang et al., hard
separation of controller and UI effects, fair baseline tuning, robust
sensitivity, validated explanations, powered objective human outcomes,
and open artifacts. What makes it weak: inventing utility weights,
reporting only mean delay and SUS, measuring only higher trust,
live-random scenarios across UI conditions, presenting template text as
XAI novelty, or claiming deployment safety from SUMO.

\section{34. Recommended Next Steps}\label{recommended-next-steps}

\begin{enumerate}
\def\labelenumi{\arabic{enumi}.}
\tightlist
\item
  Decide the primary paper identity: HCI-first is recommended for the
  current title and achievable university scope.
\item
  Write a one-page fairness policy: protected service limits, base
  weights, emergency/recovery rule, efficiency tolerance selection, and
  who has authority to change them.
\item
  Choose the governing signal standard/jurisdiction and derive
  geometry-dependent timing before building the network.
\item
  Reproduce the closest Wang et al.~controller and one delay-pressure
  baseline on the simplified intersection before implementing EFPB.
\item
  Build and property-test the safety kernel and fixed-time smoke
  scenario.
\item
  Implement the typed state, finite candidate evaluator, and immutable
  trace; validate horizon predictions on held-out SUMO rollouts.
\item
  Run a calibration-only Pareto study to select \(\delta\), debt parameters,
  horizon, normalization, and run length.
\item
  Implement counterfactual replay and mutation-test explanation fidelity
  before designing polished copy.
\item
  Acquire IMPTC samples/full access, document terms, and build a
  transparent parameter-extraction notebook/script; use SinD for
  sensitivity if feasible.
\item
  Build the same-layout D/F/FC replay interface, then conduct cognitive
  interviews and accessibility review.
\item
  Run the HCI pilot, simulate mixed-model power, preregister, obtain
  ethics approval, and freeze stimuli before confirmatory recruitment.
\item
  Run confirmatory paired-seed simulations and the user study; release
  results even if one component is null.
\end{enumerate}

A realistic sequence for one researcher is 24--32 weeks: 3--4 weeks
evidence/standards and protocol, 3--4 network/state, 4--5
controllers/traces, 3 explanation/fidelity, 3 batch calibration, 3
dashboard/stimuli, 3 pilot/ethics iteration, 4--6 confirmatory
collection/analysis, and 3--4 writing/release. Ethics review and dataset
access can run in parallel but should be treated as schedule risks.

\section{Sources}\label{sources}

\protect\phantomsection\label{source-1}{} \textbf{{[}1{]}} Wang, A.,
Zhang, K., Li, M., Shao, J., \& Li, S. (2023). ``Game Theory-Based
Signal Control Considering Both Pedestrians and Vehicles in Connected
Environment.'' \emph{Sensors}, 23(23), 9438.
\url{https://doi.org/10.3390/s23239438}

\protect\phantomsection\label{source-2}{} \textbf{{[}2{]}} Rizzo, S. G.,
Vantini, G., \& Chawla, S. (2019). ``Reinforcement Learning with
Explainability for Traffic Signal Control.'' \emph{IEEE ITSC},
3567--3572. \url{https://doi.org/10.1109/ITSC.2019.8917519}

\protect\phantomsection\label{source-3}{} \textbf{{[}3{]}} Liu, W.-L.,
Zhong, J., Liang, P., Guo, J., Zhao, H., \& Zhang, J. (2024). ``Towards
Explainable Traffic Signal Control for Urban Networks through Genetic
Programming.'' \emph{Swarm and Evolutionary Computation}, 88, 101588.
\url{https://doi.org/10.1016/j.swevo.2024.101588}

\protect\phantomsection\label{source-4}{} \textbf{{[}4{]}} He, Q., Head,
K. L., \& Ding, J. (2014). ``Multi-modal Traffic Signal Control with
Priority, Signal Actuation and Coordination.'' \emph{Transportation
Research Part C}, 46, 65--82.
\url{https://doi.org/10.1016/j.trc.2014.05.001}

\protect\phantomsection\label{source-5}{} \textbf{{[}5{]}} Ma, W., Liao,
D., Liu, Y., \& Lo, H. K. (2015). ``Optimization of Pedestrian Phase
Patterns and Signal Timings for Isolated Intersection.''
\emph{Transportation Research Part C}, 58, 502--514.
\url{https://doi.org/10.1016/j.trc.2014.08.023}

\protect\phantomsection\label{source-6}{} \textbf{{[}6{]}} Niels, T.,
Bogenberger, K., Papageorgiou, M., \& Papamichail, I. (2024).
``Optimization-Based Intersection Control for Connected Automated
Vehicles and Pedestrians.'' \emph{Transportation Research Record},
2678(2), 135--152. First published online 2023.
\url{https://doi.org/10.1177/03611981231172956}

\protect\phantomsection\label{source-7}{} \textbf{{[}7{]}} Yazdani, M.,
Sarvi, M., Bagloee, S. A., Nassir, N., Price, J., \& Parineh, H. (2023).
``Intelligent Vehicle Pedestrian Light (IVPL): A Deep Reinforcement
Learning Approach for Traffic Signal Control.'' \emph{Transportation
Research Part C}, 149, 103991.
\url{https://doi.org/10.1016/j.trc.2022.103991}

\protect\phantomsection\label{source-8}{} \textbf{{[}8{]}} Xu, K.,
Huang, J., Kong, L., Yu, J., \& Chen, G. (2023). ``PV-TSC: Learning to
Control Traffic Signals for Pedestrian and Vehicle Traffic in 6G Era.''
\emph{IEEE Transactions on Intelligent Transportation Systems}, 24(7),
7552--7563. \url{https://doi.org/10.1109/TITS.2022.3156816}

\protect\phantomsection\label{source-9}{} \textbf{{[}9{]}} Raeis, M., \&
Leon-Garcia, A. (2021). ``A Deep Reinforcement Learning Approach for
Fair Traffic Signal Control.'' \emph{IEEE ITSC}, 2512--2518.
\url{https://doi.org/10.1109/ITSC48978.2021.9564847}

\protect\phantomsection\label{source-10}{} \textbf{{[}10{]}} Liu, H., \&
Gayah, V. V. (2023). ``Total-delay-based Max Pressure: A Max Pressure
Algorithm Considering Delay Equity.'' \emph{Transportation Research
Record}, 2677(6), 324--339.
\url{https://doi.org/10.1177/03611981221147051}

\protect\phantomsection\label{source-11}{} \textbf{{[}11{]}} Du, X., Li,
Z., Long, C., Xing, Y., Yu, P. S., \& Chen, H. (2024). ``FELight:
Fairness-Aware Traffic Signal Control via Sample-Efficient Reinforcement
Learning.'' \emph{IEEE Transactions on Knowledge and Data Engineering},
36(9), 4678--4692. \url{https://doi.org/10.1109/TKDE.2024.3376745}

\protect\phantomsection\label{source-12}{} \textbf{{[}12{]}} Riehl, K.,
Kouvelas, A., \& Makridis, M. A. (2026). ``Urban Priority Pass: Fair
Signalised Intersection Management Accounting for Passenger Needs
through Prioritisation.'' \emph{Transportation Research
Interdisciplinary Perspectives}, 37, 101988.
\url{https://doi.org/10.1016/j.trip.2026.101988}

\protect\phantomsection\label{source-13}{} \textbf{{[}13{]}} Li, H., Hu,
H., Zhang, Z., \& Zhang, Y. (2023). ``The Role of Yielding Cameras in
Pedestrian--Vehicle Interactions at Un-signalized Crosswalks: An
Application of Game Theoretical Model.'' \emph{Transportation Research
Part F}, 92, 27--43. \url{https://doi.org/10.1016/j.trf.2022.11.004}

\protect\phantomsection\label{source-14}{} \textbf{{[}14{]}} Camara, F.,
Dickinson, P., \& Fox, C. (2021). ``Evaluating Pedestrian Interaction
Preferences with a Game Theoretic Autonomous Vehicle in Virtual
Reality.'' \emph{Transportation Research Part F}, 78, 410--423.
\url{https://doi.org/10.1016/j.trf.2021.02.017}

\protect\phantomsection\label{source-15}{} \textbf{{[}15{]}} Kalantari,
A. H., Yang, Y., Lee, Y. M., Merat, N., \& Markkula, G. (2023).
``Driver--Pedestrian Interactions at Unsignalized Crossings Are Not in
Line With the Nash Equilibrium.'' \emph{IEEE Access}, 11,
110707--110723. \url{https://doi.org/10.1109/ACCESS.2023.3322959}

\protect\phantomsection\label{source-16}{} \textbf{{[}16{]}} Miller, T.
(2019). ``Explanation in Artificial Intelligence: Insights from the
Social Sciences.'' \emph{Artificial Intelligence}, 267, 1--38.
\url{https://doi.org/10.1016/j.artint.2018.07.007}

\protect\phantomsection\label{source-17}{} \textbf{{[}17{]}} Wang, D.,
Yang, Q., Abdul, A., \& Lim, B. Y. (2019). ``Designing Theory-Driven
User-Centric Explainable AI.'' \emph{ACM CHI}, 1--15.
\url{https://doi.org/10.1145/3290605.3300831}

\protect\phantomsection\label{source-18}{} \textbf{{[}18{]}} Dodge, J.,
Liao, Q. V., Zhang, Y., Bellamy, R. K. E., \& Dugan, C. (2019).
``Explaining Models: An Empirical Study of How Explanations Impact
Fairness Judgment.'' \emph{ACM IUI}, 275--285.
\url{https://doi.org/10.1145/3301275.3302310}

\protect\phantomsection\label{source-19}{} \textbf{{[}19{]}} Du, N.,
Haspiel, J., Zhang, Q., Tilbury, D., Pradhan, A. K., Yang, X. J., \&
Robert, L. P. Jr.~(2019). ``Look Who's Talking Now: Implications of AV's
Explanations on Driver's Trust, AV Preference, Anxiety and Mental
Workload.'' \emph{Transportation Research Part C}, 104, 428--442.
\url{https://doi.org/10.1016/j.trc.2019.05.025}

\protect\phantomsection\label{source-20}{} \textbf{{[}20{]}} Zhu, Z.,
Li, X., Delhomme, P., Schroeter, R., Glaser, S., \& Rakotonirainy, A.
(2025). ``Human-Centric Explanations for Users in Automated Vehicles: A
Systematic Review.'' \emph{Accident Analysis \& Prevention}, 220,
108152. \url{https://doi.org/10.1016/j.aap.2025.108152}

\protect\phantomsection\label{source-21}{} \textbf{{[}21{]}} Starke, C.,
Baleis, J., Keller, B., \& Marcinkowski, F. (2022). ``Fairness
Perceptions of Algorithmic Decision-Making: A Systematic Review of the
Empirical Literature.'' \emph{Big Data \& Society}, 9.
\url{https://doi.org/10.1177/20539517221115189}

\protect\phantomsection\label{source-22}{} \textbf{{[}22{]}} Wang, X.,
\& Yin, M. (2022). ``Effects of Explanations in AI-Assisted Decision
Making: Principles and Comparisons.'' \emph{ACM Transactions on
Interactive Intelligent Systems}, 12, 1--36.
\url{https://doi.org/10.1145/3519266}

\protect\phantomsection\label{source-23}{} \textbf{{[}23{]}} Lee, J. D.,
\& See, K. A. (2004). ``Trust in Automation: Designing for Appropriate
Reliance.'' \emph{Human Factors}, 46, 50--80.
\url{https://doi.org/10.1518/hfes.46.1.50_30392}

\protect\phantomsection\label{source-24}{} \textbf{{[}24{]}} Federal
Highway Administration. (2023). \emph{Manual on Uniform Traffic Control
Devices, 11th ed., Part 4}, pedestrian signal timing guidance.
\href{https://mutcd.fhwa.dot.gov/pdfs/11th_Edition/part4.pdf}{Official
PDF}

\protect\phantomsection\label{source-25}{} \textbf{{[}25{]}} Eclipse
SUMO. ``Interfacing TraCI from Python.''
\href{https://sumo.dlr.de/docs/TraCI/Interfacing_TraCI_from_Python.html}{Official
documentation}

\protect\phantomsection\label{source-26}{} \textbf{{[}26{]}} Eclipse
SUMO. ``Traffic Lights.''
\href{https://sumo.dlr.de/docs/Simulation/Traffic_Lights.html}{Official
documentation}

\protect\phantomsection\label{source-27}{} \textbf{{[}27{]}} Eclipse
SUMO. ``TraCI Pedestrian Crossing Tutorial.''
\href{https://sumo.dlr.de/docs/Tutorials/TraCIPedCrossing.html}{Official
documentation}

\protect\phantomsection\label{source-28}{} \textbf{{[}28{]}} Jian,
J.-Y., Bisantz, A. M., \& Drury, C. G. (2000). ``Foundations for an
Empirically Determined Scale of Trust in Automated Systems.''
\emph{International Journal of Cognitive Ergonomics}, 4(1), 53--71.
\url{https://doi.org/10.1207/S15327566IJCE0401_04}

\protect\phantomsection\label{source-29}{} \textbf{{[}29{]}} Colquitt,
J. A. (2001). ``On the Dimensionality of Organizational Justice: A
Construct Validation of a Measure.'' \emph{Journal of Applied
Psychology}, 86(3), 386--400.
\url{https://doi.org/10.1037/0021-9010.86.3.386}

\protect\phantomsection\label{source-30}{} \textbf{{[}30{]}} Brooke, J.
(1996). ``SUS: A `Quick and Dirty' Usability Scale.'' In \emph{Usability
Evaluation in Industry}, 189--194.
\href{https://hci-studies.org/methods-and-measures/downloads/SUS_Brooke1996.pdf}{Author-version
PDF}

\protect\phantomsection\label{source-31}{} \textbf{{[}31{]}} Hoffman, R.
R., Mueller, S. T., Klein, G., \& Litman, J. (2023). ``Measures for
Explainable AI: Explanation Goodness, User Satisfaction, Mental Models,
Curiosity, Trust, and Human--AI Performance.'' \emph{Frontiers in
Computer Science}, 5, 1096257. The paper provides the Explanation
Satisfaction Scale and cautions that self-reported satisfaction is only
one part of evaluation. \url{https://doi.org/10.3389/fcomp.2023.1096257}

\protect\phantomsection\label{source-32}{} \textbf{{[}32{]}} Eclipse
SUMO. ``Downloads.'' Current release and Python package options.
\href{https://sumo.dlr.de/docs/Downloads.php}{Official documentation}

\protect\phantomsection\label{source-33}{} \textbf{{[}33{]}} Hetzel, M.,
Reichert, H., Reitberger, G., Fuchs, E., Doll, K., \& Sick, B. (2023).
``The IMPTC Dataset: An Infrastructural Multi-Person Trajectory and
Context Dataset.'' \emph{IEEE Intelligent Vehicles Symposium}, 1--7.
\href{https://doi.org/10.1109/IV55152.2023.10186776}{DOI};
\href{https://github.com/kav-institute/imptc-dataset}{official
repository}

\protect\phantomsection\label{source-34}{} \textbf{{[}34{]}} Xu, Y.,
Shao, W., Li, J., Yang, K., Wang, W., Huang, H., Lv, C., \& Wang, H.
(2022). ``SIND: A Drone Dataset at Signalized Intersection in China.''
\emph{IEEE ITSC}, 2471--2478.
\href{https://doi.org/10.1109/ITSC55140.2022.9921959}{DOI};
\href{https://github.com/SOTIF-AVLab/SinD}{official repository}

\protect\phantomsection\label{source-35}{} \textbf{{[}35{]}} Li, Y., et
al.~(2026). ``SinD 2.0: A Multi-City UAV Dataset with Semantic Risk
Annotations for SOTIF-Oriented Safety Validation at Signalized
Intersections.'' arXiv:2607.16943.
\href{https://github.com/SOTIF-AVLab/SinD}{Official repository and
status}

\protect\phantomsection\label{source-36}{} \textbf{{[}36{]}} U.S.
Department of Transportation. (2026). ``Intersection Safety Challenge
Stage 1B Validation Data.''
\href{https://catalog.data.gov/dataset/intersection-safety-challenge-stage-1b-validation-data}{Official
Data.gov record}

\protect\phantomsection\label{source-37}{} \textbf{{[}37{]}} Bock, J.,
Krajewski, R., Moers, T., Runde, S., Vater, L., \& Eckstein, L. (2020).
``The inD Dataset: A Drone Dataset of Naturalistic Road User
Trajectories at German Intersections.'' \emph{IEEE Intelligent Vehicles
Symposium}, 1929--1934.
\url{https://doi.org/10.1109/IV47402.2020.9304839}

\protect\phantomsection\label{source-38}{} \textbf{{[}38{]}} Kizilcec,
R. F. (2016). ``How Much Information? Effects of Transparency on Trust
in an Algorithmic Interface.'' \emph{ACM CHI}.
\url{https://doi.org/10.1145/2858036.2858402}

\protect\phantomsection\label{source-39}{} \textbf{{[}39{]}} Elsevier.
``Transportation Research Part C: Emerging Technologies---Aims and
Scope.''
\href{https://shop.elsevier.com/journals/transportation-research-part-c-emerging-technologies/0968-090X}{Official
journal page}

\protect\phantomsection\label{source-40}{} \textbf{{[}40{]}} IEEE
Intelligent Transportation Systems Society. ``IEEE Transactions on
Intelligent Transportation Systems---Scope.''
\href{https://ieee-itss.org/pub/t-its/}{Official page}

\protect\phantomsection\label{source-41}{} \textbf{{[}41{]}} Elsevier.
``Transportation Research Part F: Traffic Psychology and
Behaviour---Aims and Scope.''
\href{https://shop.elsevier.com/journals/transportation-research-part-f-traffic-psychology-and-behaviour/1369-8478}{Official
journal page}

\protect\phantomsection\label{source-42}{} \textbf{{[}42{]}} ACM
AutomotiveUI 2026. ``Papers'' and conference scope.
\href{https://www.auto-ui.org/26/papers/}{Official call}

\protect\phantomsection\label{source-43}{} \textbf{{[}43{]}} ACM IUI
2026. ``Call for Papers.''
\href{https://iui.acm.org/2026/call-for-papers/}{Official call}

\section{``If I Were Implementing This Research From
Zero''}\label{if-i-were-implementing-this-research-from-zero}

I would frame the project as an HCI-first smart-infrastructure study
supported by a credible control contribution. I would not say that
drivers and pedestrians themselves are rational Nash agents, and I would
not sell ``Nash bargaining at an intersection'' as new. The closest
paper already does that. My narrow claim would be that a public
allocation rule can be made starvation-safe, efficiency-bounded,
faithfully explainable per decision, and evaluated for both traffic
outcomes and appropriate human reliance.

I would build one four-arm SUMO intersection with \texttt{V\_NS},
\texttt{V\_EW}, and an exclusive \texttt{P\_ALL} phase, initially
omitting turns. SUMO would own the transition safety logic. A small
Python controller would enumerate safe plans over a short horizon.
Safety, pedestrian clearance, emergency service, sensor fallback, and
maximum-wait constraints would be hard rules. It would identify the
lowest predicted-delay plan, retain only candidates within calibrated
tolerance \(\delta\), and select from that set with waiting-debt-weighted
Nash bargaining. Vehicle direction would use total-delay pressure. If a
tuned weighted sum matched it, I would simplify the final system and
report that result honestly.

I would use IMPTC as the first calibration source because it contains
vehicles, vulnerable road users, and signal status at a real
instrumented intersection. SinD would be a useful regional sensitivity
check. I would use those data to bound arrivals, walking speeds, and
discharge behavior---not pretend the simplified network is a digital
twin. All confirmatory traffic comparisons would use pre-generated
manifests and common random seeds.

The first software milestone would be the safety kernel and fixed-time
baseline, not the dashboard. The second would be a typed state collector
with explicit missing/stale values. The third would be the efficiency
baseline and closest equal-weight bargaining comparator. Only then would
I implement EFPB and validate horizon predictions. Every decision would
save a complete candidate table and configuration hash.

The explanation system would be deterministic. It would say what was
chosen, show at most three decision-relevant facts, identify the closest
rejected plan and its constraint/score reason, and optionally show the
nearest feasible state change that truly flips the controller after
replay. If replay does not validate the counterfactual, the interface
would abstain. I would mutation-test the explanation engine by
deliberately changing numbers, comparison directions, and constraints
and ensuring validation fails.

For the dashboard, I would use React/TypeScript/Vite and a FastAPI
backend with a WebSocket. I would reuse the exact same renderer for live
demonstrations and fixed study replays. The interface would be an audit
tool for road users/operators, not something a driver reads while moving
and not a replacement for WALK/DON'T WALK indications.

Technically, I would compare fixed-time, delay-pressure/efficiency-only,
closest equal-weight bargaining, and EFPB under paired seeds. I would
report mode-specific means and tails, total delay, queues, throughput,
service ratios, maximum waits, starvation, emergency response, phase
switches, safety invariants, fairness distributions, bound compliance,
and latency. The decisive plots would be paired differences and the
efficiency--tail-fairness Pareto frontier. I would run the
no-starvation, no-efficiency-bound, static-weight, and weighted-sum
ablations.

For the human study, I would use 90 usable adults as an initial target,
finalized by pilot-based mixed-model power. Each person would see
decision-only, factual, and factual-plus-counterfactual conditions in
counterbalanced blocks, with matched but non-repeated replay scenarios.
Half of trials would contain sound decisions and half physically safe,
procedurally anomalous decisions from a fault-injection controller. The
main outcomes would be reason accuracy, decision-flip accuracy, log
response time, appropriate endorsement, confidence/Brier score, adapted
procedural/distributive/informational justice, Jian trust, and overall
SUS.

I would analyze trial correctness with logistic mixed models, response
time and suitable composites with mixed models, ordinal ratings with
cumulative-link mixed models, and test the
explanation-by-decision-quality interaction. I would correct planned
contrasts with Holm, report effect sizes and confidence intervals, and
preregister exclusions. The key success is not higher trust; it is
better understanding and better discrimination between good and bad
automation.

The largest risks would be arbitrary policy weights, a weak closest
baseline, an inaccurate horizon model, simulation-only generalization,
an unvalidated fairness questionnaire, and a polished explanation that
persuades without informing. I would mitigate them with policy
provenance, calibration-only tuning, sensitivity analysis, a reproduced
prior bargaining baseline, held-out prediction checks, cognitive
interviews, objective comprehension, anomalous trials, and full trace
release.

If time forced cuts, I would preserve the safety kernel, three technical
controllers, trace fidelity, D/F/FC replay study, objective
comprehension, trust calibration, and reproducibility package. I would
cut deep learning, individual bidding, a database for simulation, cloud
infrastructure, turns, network coordination, and generative text. That
smaller project would still make a coherent publication claim:
\textbf{faithful explanations of a transparent, fairness-constrained
smart-intersection allocator can be evaluated as part of the control
system rather than added as decorative interface text.}

\end{document}
